% Reduction of the Wagstaff Chebyshev sufficiency conjecture to the
% Single-Pass Conjecture.
% v4.0 addendum 2 (2026-07-16, cp415): pre-release inspection fixes.
%   cor:ladder: the three formerly PRP-gated rungs 83341/117241/127033
%   are in fact certified -- W_83339/W_117239/W_127031 carry ECPP proofs
%   (2014, 2022, 2023; t5k top-20 Wagstaff) -- so the ladder has eleven
%   rungs and is complete (every other certified Wagstaff exponent q has
%   q+2 composite). rem:shape-inputs ledger completed: danger-census
%   bullet added (used in (a5)/(b1)/(b5)); the fully-algebraic item list
%   corrected ((0) modulo the Wieferich clause, (b1) modulo the census
%   floor, (a3) modulo the 3200p instantiation). sec:empirical
%   selection-bias note recounted: 184,938 = exponents carrying >= 2
%   candidates (one recorded pair each; 369,876 tested members), now
%   consistent with the distribution table. thm:quartic rank-order
%   paragraph: the two-component reading of the s-coefficient argument
%   made explicit. thm:caseA order parenthetical corrected.
%   rem:pell-wss-global: surplus lem:v2-class citation dropped (the CRT
%   argument is self-contained).
% v4.0 addendum (2026-07-16, cp414): healing pass restoring the v3.8
%   content whose loss was a real compromise, in strengthened v4.0 form:
%   new subsection sec:variants in the reduction section --
%   rem:counting-form (quantitative N_in <= Pi + 1, strictly sharper than
%   v3.8's E_3 <= Pi + W since no unbounded W term; N_in named to avoid
%   colliding with the Pell exclusion numbers N_d; input-free pair-level
%   form replacing the old triple-level rem:input-free) +
%   cor:triple-erosion ((T) = |T_p^>| <= 2 and (C) = Wieferich-paired
%   corner, both implied by PPPC; proof via branch decomposition, no
%   bucket taxonomy) + rem:erosion-terminal (corner doubly conditioned;
%   terminality + function-field realization). Also: super-Wieferich
%   double-wall paragraph after cor:multiplicity-wieferich; phantom-scope
%   parenthetical at sec:x1 head + "of composite W_p" in thm:platinum's
%   verification; M2 self-consistency sentence restored (0.07 / 3
%   predicted vs 0 observed); Moreover divisor note (one modexp).
%   Second adversarial pass on the new blocks: all SOUND; also
%   prime-power wording in the double-wall paragraph; bridge sentence
%   before thm:spc-intro; spelling unified (-ize); Convention paragraph
%   (certificate-backed proofs vs stated computational inputs) added to
%   sec:toolchain.
% Version 4.0 (2026-07-16, cp413): restructure of the spine around the
%   Single-Pass Conjecture. New Section 4 "Local passes and the Single-Pass
%   Conjecture": the passer set P_p (Definition def:local-pass), containment
%   and restriction lemmas, Conjecture SPC (|P_p| <= 1 for every prime
%   p >= 5), and Theorem thm:spc-reduction (Theorem A): SPC implies
%   Conjecture conj:suff, unconditionally and with no computational input
%   (the only external ingredient is the Nagell-Ljunggren lemma
%   lem:w1-empty, moved into Section 2). Branch decomposition of a double
%   pass (split members = compatible triples, inert members = danger
%   triples); Theorem B (thm:sharpened-chain) re-proved directly from the
%   branch decomposition, 3x shorter; the second-moment counting bound, the
%   W-residual trichotomy and the triple form (buckets E_3, W, W_1, W_2,
%   W_2', X_{E_3}, PPPC3) are cut -- the two-distinct-factor input now
%   comes unconditionally from lem:w1-empty, making the multiplicity
%   buckets unnecessary; the X1 subsection and the multiplicity/Wieferich
%   constraints relocated into the cross-case section; the PPPC heuristic
%   (M1-M3 model, all calibrated numbers unchanged) and the NCT
%   verification appendix compressed, with an SPC-level closing sentence
%   added to the heuristic and the danger-census artifact
%   (cp412_danger_census.json) added to the supplementary-bundle list.
%   Hypotheses of the working chain unchanged: PPPC + MPC^> + X1 with
%   Platinum + discharges as inputs.
% v3.8 addendum 2 (2026-07-16, cp412): Proposition prop:danger-census-400
%   (complete danger-triple census for admissible 43 < d <= 400: 34 rungs,
%   complete V_{2d} factorizations, 226 primes APR-CL, 340 tests; exactly
%   2 real triples -- (57, 171369731867, 85684865933) and (67,
%   148937135261632265803, 24822855876938710967), both closed at secondary
%   factors -- plus phantoms (57,683,11), (331,43691,17)) + Corollary
%   cor:inert-floor-400 (every inert factor of a global passer has
%   d_r > 400) + rem:two-slices. Diagonal vanishing extended 124 -> 400;
%   cor:diagonal-k500 upgraded to k <= 1604; shape corollary (b1)/(b5),
%   Table 1, abstract, unconditional-contributions list updated. T1-verified
%   (independent re-implementation, 0 defects in the survivor set).
% v3.8 addendum (2026-07-16): Proposition prop:ladder (self-supporting
%   exponents: W_{p-2} prime => Condition (II) is an unconditional
%   primality test at p; certified ladder p in {5,7,13,19,103,193,349,
%   3541} + three PRP-gated rungs 83341/117241/127033) + cor:ladder +
%   rem:ladder, placed before the R3 subsection (cp410).
% Version 3.8 (2026-07-15). Built on v3.7. v3.8: new Corollary
%   cor:counterexample-shape ('Shape of a counterexample', unconditional
%   modulo the named computational inputs) + rem:shape-inputs + discussion
%   at the head of sec:open-problems; the one-sentence PPPC pair heuristic
%   in sec:heuristic upgraded to the explicit M1-M3 model with the
%   bracketed expected count (cp395).
% Version 3.7 (2026-07-15). Built on v3.6.
%
% v3.7: W_1 bucket closed unconditionally via Nagell-Ljunggren at x = -2
% (Bennett-Levin); sharpened chain now PPPC + MPC^> + X1, triple chain
% MPC^> + PPPC3 + (W_2'=0) + X1; SW3 removed; Vrba-Reix framing corrected
% (literal Reix iteration lives in Q(sqrt(-7)) and is incomparable with
% Condition (II); companion note). NCT: d = 227, 283, 331, 379 closed by
% complete fixed-d certificates enabled by the Pell half-split
% V_2d = 2(2P_d-1)(2P_d+1) (cp367/cp372) + GNFS on the last two
% half-cofactors C104/C106 (cp374-cp376); every parity-unblocked value
% of (200,400] is now closed: NCT unconditional for every odd d <= 400.
%
% Version 3.6 (2026-06-12). Built on v3.5.
%
% v3.6 integrates a verification-gated research round on the W-residual /
% PPPC interface. The main reduction statement and all residual hypotheses
% of v3.5 are unchanged; the new material erodes the role of PPPC in the
% sharpened chain from a pair count to a triple count. Changes:
%   - New subsection "The triple form: eroding the pair count" at the end
%     of Section 6 (after the sharpened-chain remark): the above-cutoff
%     triple count Pi_3 = sum_p binom(|T_p^>|, 3) and Conjecture PPPC3
%     (Pi_3 = 0; implied by PPPC, converse not asserted; distinct from the
%     cutoff-free triple residual Pi_3' of the computational-input-free
%     remark, which above-cutoff PPPC does not imply); the Wieferich-paired
%     corner W_2' (exactly two distinct prime factors, both > 10^12, W_p
%     not squarefree) and Conjecture W_2' = 0 (implied by either PPPC or
%     W = 0; the non-squarefree clause is redundant for all-inert
%     two-factor W_p by mod-8 parity and kept for robustness).
%     Theorem [Triple trichotomy]: deterministically, given the Platinum
%     enumeration, E_3 <= Pi_3 + W_1 + W_2' + X_{E_3}, where X_{E_3}
%     counts witness exponents lying in E_3; given also the discharge
%     verifications, X_{E_3} <= 1 (only 10916765939 can contribute) and
%     X1 forces X_{E_3} = 0; every W_2'-exponent carries a repeated
%     base-2 Wieferich prime = 3 mod 8 above 2^64 plus a second large
%     locally passing factor and has |T_p^>| = 2 exactly (the
%     global-passer Pell-Wall-Sun-Sun statement kept as a remark, as in
%     the W-residual trichotomy). Corollary [Triple-sharpened chain]:
%     Conjecture [suff] from MPC^> + PPPC3 + (W_1 = 0) + (W_2' = 0) + X1,
%     with the same two stated computational inputs as the sharpened
%     chain and no new input; the conjunction is implied by the
%     sharpened-chain conjunction {PPPC, MPC^>, W_1 = 0, X1}, hence also
%     by {PPPC, NCT-general, W = 0}. A closing remark records the
%     heuristic reading (a triple coincidence in place of a pair), the
%     empirical records (|T_p^>| <= 1 wherever computed, so Pi, Pi_3 and
%     W_2' all vanish as far as computed), and that PPPC3 meets the
%     structural obstruction unchanged (the erosion shrinks the
%     footprint, not the wall).
%   - Abstract, S1.2 (Stage C) and the Open problems intro announce the
%     triple form in one sentence each.
%   - NCT range extension into (200, 400]: fifteen new complete fixed-d
%     closures. (200,300]: 201, 209, 249, 251, 281, 297 from public-table
%     factorizations (37 factors APR-CL, orders by descent from
%     completely factored r-1); 243 (C125) and 211 (C161) fell entirely
%     to local ECM (~35 and ~20 minutes); (300,400]: 307, 321, 347, 361,
%     387, 393 from public-table factorizations (46 factors APR-CL);
%     363 (C169) by local ECM in one round (7 factors APR-CL).
%     Open in (200,400]: d = 227, 283, 331, 379. Abstract,
%     S1.2, Table 1, Problem 1, the NCT open-problem subsection and the
%     computational appendix updated accordingly.
%   - X1 subsection and C.4 (inert enumeration) record the completed
%     directed scan at the third witness exponent: no second prime factor
%     of W_10916765939 in the progression r = 2pk+1 with k <= 10^12; the
%     witness factor itself, at k = 7, is the only hit.
%   - Empirical extensions: corner census extended to admissible
%     d <= 1000, condition-(v) p <= 1e8 (9,506,750 rectangles; four
%     hits, all known phantom rows at prime W_p; zero pairs); deep-k
%     pre-window at 30,000 exponents, k <= 1e6 (6,559 (i)-(iv)
%     candidates, zero condition-(v) passes; 623 multi-candidate
%     exponents, all pairs satisfying the separation constraints).
%   - Terminality block (gate-verified): Remark [The triple form is
%     terminal for the count axis] — 3 never divides W_p, so squarefree
%     <=> all factors non-Wieferich (via Lemma [Valuation transfer]); the
%     level-4 display E_3 <= Pi_4 + W_1 + W_2' + W_3' + W_3^sf + X_E3 is
%     stated for orientation; the squarefree omega-3 all-large stratum is
%     untaggable by every rarity tag of the paper (parity spent,
%     multiplicity definitionally absent, no Platinum factor), the
%     global-pass route adds nothing on squarefree W_p (CRT sentence
%     added to Remark [Pell-WSS for global passers]), and the stratum is
%     realized in the Carlitz calibration (explicit triples/quadruples,
%     q = 3, 5). Remark [Cross-factor cofactor congruence]: Condition
%     (II) at inert r <=> W_p/r == 1 (mod 4 d_r) (mod 8 d_r free); the
%     omega-3 closure system; perimeter, not a tag. Remark [Same-d
%     triples force alignment] (Appendix A.2): pigeonhole closes the
%     anti-aligned escape for triples; same-d triples dead at admissible
%     d <= 124 via the diagonal closures.
%
% Version 3.5 (2026-06-11). Built on v3.4.
%
% v3.5 integrates a verification-gated research round on the PPPC residual.
% The main reduction statement and all residual hypotheses are unchanged; the
% new material consists of unconditional structural constraints on the PPPC
% configuration and certified diagonal closures. Changes:
%   - New subsection "Pair separation" in Section 6 (after the small-k
%     valuation obstruction): Theorem [Pair separation] -- for two inert
%     candidates r_i = 2 p k_i + 1 = 3 mod 8 at the same p, with odd
%     d_i | (r_i+1)/4: (1) k_1 = k_2 = p (mod 4); (2) gcd(p k_1 + 1,
%     p k_2 + 1) | k_2 - k_1; (3) 4 gcd(d_1, d_2) | k_2 - k_1, equivalently
%     8 p gcd(d_1, d_2) | r_2 - r_1; (4) two same-d danger triples force
%     4 d | k_2 - k_1 with d >= 57, so k_2 - k_1 >= 228 and
%     r_2 - r_1 >= 456 p. Items (1)-(3) use neither Condition (II) nor
%     primality of p. Proposition [Pair-coupled valuation transfer]: the
%     cross-factor and half-order channels, strict strengthenings of the
%     small-k valuation obstruction (example: d = 737 at k = 1). Remarks
%     record the falsity of the rad(d) | 2k+1 strengthening (sharpness
%     witness (r,p,k,d) = (43,7,3,11)), the congruence-level consistency of
%     bounded-k off-diagonal pairs (witness class p = 166781 mod 1170324 at
%     (k_1,k_2,d_1,d_2) = (1,5,57,59); the missing input is realizability),
%     and the r+1-side gcd regularity (gcd(r_1+1, r_2+1) = 4*odd dividing
%     2(k_2-k_1)).
%   - Appendix A.2 (PPPC residual) gains the diagonal half-pinning package:
%     Lemma [Half-pinning] (every prime r > 3 dividing V_{2d} with
%     ord_r(2) = 2m, m odd, has V_d = eps 2^{(m+1)/2} mod r for a unique
%     sign), Proposition [Corner localization] (the order-divisibility locus
%     at (d,p) equals the prime divisors > 3 of four explicit gcds
%     G(eps,eps'); factorization-free, four modular exponentiations; the
%     locus over-captures the danger locus, harmlessly for exclusion),
%     Remark [Pair caps] (three of the four sign classes are product-capped;
%     only anti-aligned pairs escape; the same-eps cap improves on the
%     trivial bound exactly for p < 4 log2(1+sqrt2) d = 5.0862 d).
%     Proposition [Diagonal vanishing through d = 124]: S_diag(d) = 0 for
%     all p at every admissible d <= 124, via complete APR-CL-certified
%     factorizations of V_118, V_162, V_166, V_198, V_214, V_242 (with the
%     prior V_114, V_134); Corollary: no diagonal danger pair with both
%     k_i <= 500, unconditionally. Remark: the 29,930-rectangle
%     factorization-free corner census (admissible d in (43, 124],
%     condition-(v) p <= 10^6; single hit, the known phantom (57, 11, 683)).
%   - Directions (A.4) and Open problems: the diagonal residue algebra is
%     complete (the natural residue span at a diagonal-eligible prime
%     collapses into <2>; only the two sign bits (eps, eps') survive); the
%     remaining diagonal content is named the anti-aligned corner question;
%     honest scope note (the package transfers verbatim to F_q[t], so it
%     localizes the diagonal difficulty rather than closing it).
%   - Abstract, Section 1 and Table 1 note the two new unconditional
%     constraints (one sentence each); the conditional-reduction framing is
%     unchanged. New reference: FactorDB.
%   - Layout fix: amsart emits the abstract as one unbreakable box; at this
%     abstract length the box overflowed the page and the final lines were
%     silently clipped in the built PDF (already the case in v3.4: the
%     closing sentence and the companion-paper sentence did not appear).
%     The preamble now unvboxes the abstract so it breaks across pages;
%     formatting is unchanged.
%
% Version 3.4 (2026-06-11). Built on v3.3.
%
% v3.4 integrates a verification-gated research round. The v3.3 main
% reduction statement (PPPC + NCT-general + W = 0) is unchanged and kept as
% the base theorem; the new material extends the NCT range and sharpens the
% residuals. Changes:
%   - NCT range extension: complete fixed-d closures for d = 129, 139, 171,
%     177 via complete primitive-part factorizations, every listed factor
%     certified prime by APR-CL; the remaining parity-unblocked values in
%     (99, 200] are all closed: d=163, 179 by complete ECM factorizations of
%     Psi_652/Psi_716 and d=131 by GNFS on Psi_524 (all APR-CL-certified;
%     certificates scripts/cp353_nct_certificate_{131,163,179}.py). NCT now
%     holds unconditionally for every odd d <= 200 (updated throughout).
%   - New Mixed Pair Conjecture (MPC) on the split branch: no prime p admits
%     both a split compatible triple and an inert local pass; implied by
%     NCT-general, hence formally weaker as a hypothesis. Theorem
%     [Mixed-pair form of the main reduction]: Conjecture [suff] follows
%     from PPPC + MPC + (W = 0); the v3.3 main theorem becomes an immediate
%     corollary.
%   - W-residual trichotomy: Lemma [Valuation transfer] v_r(W_p) =
%     v_r(2^{r-1} - 1); E_3 <= Pi + W_1 + W_2, where W_1 requires a base-2
%     super-Wieferich prime = 3 mod 8 and W_2 is pinned to the Platinum
%     witness exponents; bucket-level weakening W = 0 => W_1 = W_2 = 0; SW3
%     recorded as a sufficient condition for W_1 = 0 (not comparable with
%     W = 0 as universal statements). Buckets are local-pass-defined; the
%     Pell-Wall-Sun-Sun condition is kept as a global-passer remark.
%   - X1: the finite witness-exponent leg reduced to one explicit exponent;
%     two of the three witness exponents discharged at explicit failing
%     factors (the d = 57 secondary factor; a new factor of W_1251488009).
%   - Sharpened chain: Conjecture [suff] <== PPPC + MPC^> + (W_1 = 0) + X1,
%     the conjunction implied by the v3.3 conjunction
%     {PPPC, NCT-general, W = 0}.
%   - Heuristic section: clarify the implicit inert-fraction factor in the
%     displayed E(D); add the MPC joint heuristic (ceiling 8.7e-4 / central
%     4.4e-4 / best-estimate 9.8e-10, vs NCT-general single-event 1.27e-2 /
%     2.4e-5; model-dependent, heuristic only) and the 1/(3d) split-side
%     suppression remark.
%
% Version 3.3 (2026-06-05). Built on v3.2.
%
% v3.3 is a final referee-hardening and math-salvage patch to v3.2. The main
% conditional reduction statement is unchanged. Changes:
%   - Correct the Wieferich wording: the search below 2^64 excludes base-2
%     Wieferich primes congruent to 3 mod 8, not all base-2 Wieferich primes.
%   - Add the inert primitivity lemma: ord_r(omega_3)=4d forces r to be a
%     primitive divisor of V_{2d} in the d-indexed Pell--Lucas family.
%   - Add a small-k valuation obstruction for inert candidates r=2pk+1.
%   - Add complete fixed-d NCT closures for d=107 and d=121, certified by the
%     cp350 verification script.
%   - Record a fixed-d NCT certificate corollary and soften analytic-direction
%     claims around Chebotarev / BFI.
%
% Version 3.2 (2026-06-04). Built on v3.1.
%
% v3.2 is a referee-hardening patch to v3.1. The conditional reduction
% statement is unchanged. Changes:
%   - Correct the computational-input-free triple residual: it is a separate
%     cutoff-free residual, not a consequence of above-cutoff PPPC alone.
%   - Reword the NCT d=100..200 computation as a finite-window search
%     (r < 10^11), not a full verification of that d-range.
%   - Fix the Platinum outcome table counts and distinguish Platinum local-pass
%     witnesses from secondary closures in the d<=1000 catalog.
%   - Restrict the number-field equivalent form of PPPC to prime half-orders.
%   - Correct the d=83,99 NCT closing sentence and soften several research-
%     direction claims.
%
% Version 3.1 (2026-06-04). Built on v3.0.
%
% v3.1 is a correctness patch to the lean arXiv draft. The main reduction
% statement is unchanged. Changes:
%   - Separate the true residual set T_p^> from the small-factor Platinum set
%     and from the p-indexed pre-condition-(v) search window.
%   - Correct the PPPC equivalent forms to include the diagonal d_1 = d_2 case.
%   - Replace the prime-only AP shorthand by valuation-adjusted modulus wording
%     where composite d is discussed.
%   - Correct the inert product-bound factors from q = 3 mod 8 to q in A.
%   - Clean up split/NCT logic and small rank/order proof details.
%
% v3.0 is the lean arXiv draft. The reduction statement and residual
% hypotheses are unchanged. Changes:
%   - Remove the independent split-local algebra appendix from the main paper.
%   - Compress the closed-approaches material into a short residual-obstruction
%     appendix, with the longer diagnostic log left to supplementary material.
%   - Compress the computational appendix to the certification facts needed by
%     the reduction.
%   - Move references to the end of the paper.
%
% v2.9 is an arXiv-readiness polish; the mathematics of the reduction is
% unchanged. Changes:
%   - Add a short title to fix the running-head overfull boxes.
%   - Make the abstract and main theorem framing still more explicit that this
%     is a conditional reduction, not an unconditional proof of the sufficiency
%     conjecture.
%   - Replace the ambiguous BPSW/Miller-Rabin "certified" wording in the NCT
%     finite-exhaustion note with APR-CL certification wording matching the
%     software/toolchain section.
%   - Use unnumbered tagged displays in the quartic-residuosity section to remove
%     duplicate hyperref equation anchors.
%   - Update project metadata to this draft.
%
% v2.8 (2026-05-29) is a compaction / editorial revision; the mathematics of
% the reduction is unchanged from v2.7. Section 9 (Open problems) is consolidated:
%   - The five satellite no-go Observations of S9.2 (Coincidence obstruction,
%     Geometric uniformity, Iwasawa invariants, the two unpinned-congruence
%     observations) and the two Remarks (p-adic linear forms, BCZ gcd) are
%     merged into a single "Further methods, same wall" survey, each tool's
%     specific reason preserved as a one-line entry under the unified
%     Observation [additive-coupling / multiplicative-projection obstruction].
%   - That master Observation is corrected: the three symbol-faces (R)/(C)/(A)
%     are now used with one consistent meaning throughout (a prior reuse of
%     the same letters with different referents in the Coincidence
%     observation is removed), and the analytic divisor-vs-AP counting wall
%     specific to the inert pair-count Pi is named explicitly and tied to
%     Proposition [Frobenius cannot see the order of 2] and the Finding of
%     S9.2.3.
%   - Proposition [Frobenius cannot see the order of 2] and the then-current
%     Pi < infinity diagnostic were retained as the two load-bearing results of
%     the section. Version 3.1 later weakens the latter to an AP-only limitation
%     after inserting the valuation-adjusted modulus L_d.
%   - S1.2 gains a reduction-status table (proved vs. verified vs. open for
%     each of the three residual hypotheses), consolidating figures
%     previously scattered across S1, S5, S6, S10.
%
% v2.7 added the unified obstruction and further closed leads.
% v2.4 folds in two findings from a post-v2.3 attack round; the mathematics
% of the reduction is unchanged from v2.3.
%   - Definition [Compatible triple] gains Remark [No phantom compatible
%     triples]: every compatible triple has W_p composite (r = 1 mod 8,
%     W_p = 3 mod 8), so the nonexistence of compatible triples coincides
%     with its restriction to composite W_p -- the only form the reduction
%     invokes; no phantom case arises on the split branch, in contrast to
%     the inert-branch danger triples.
%   - Section 9.1 gains a fifth closed lead: a multiplicative identity
%     coupling a split factor r to a second prime of W_p is also
%     mu_p-tautological (reduction modulo r is F_r-internal), extending the
%     single-factor barrier to cross-factor identities.
%
% v2.3 is an exposition / editorial revision; the mathematics of the
% reduction is unchanged from v2.2. Changes:
%   - Second-Moment Reduction: a computational-input-free form of the
%     bound (E_3 <= Pi_3' + W) was recorded as Remark
%     [Computational-input-free form of the bound]. Version 3.2 later
%     clarifies that Pi_3'=0 is a separate cutoff-free triple residual, not
%     a consequence of above-cutoff PPPC alone.
%   - Heuristic section gains a paragraph: Conjecture 5.2 is structural --
%     it holds by the Wagstaff Order-/Multi-Factor Pinning, not by a
%     generic scarcity of pseudoprimes (the |Q|=1 Chua test is generically
%     non-conclusive).
%   - Non-load-bearing material -- Theorem [Quartic residuosity] and the
%     mod-16 refinement of S2.3 -- explicitly flagged as context, matching
%     the existing flags on Appendix A and the d<=1000 catalog.
%   - Editorial: "Conjecture 5.2 of the project corpus" -> the named
%     conjecture; four previously-uncited references (BLS, Cohen-Lenstra,
%     Moree, Williams) cited at their natural points.
%
% v2.2 is a critical-correctness revision. An independent 3-agent adversarial
% verification of the core reduction confirmed the main result sound --
% Conjecture 5.2 follows from PPPC + NCT-general + (W=0), and the Second-
% Moment Reduction and its proof are unaffected -- but found two statement-
% level errors in Sections 5-6 supporting theorems, now corrected:
%   - Theorem [Pell-sequence exclusion for d<=43] (S5) was stated without
%     the Condition (II) hypothesis its proof actually uses, making the
%     literal statement false (inert r with ord_r(omega_3)=4d, d<=43, do
%     occur on composite W_p when Condition (II) fails). Restated with the
%     Condition (II) hypothesis; proof case 3 corrected; the r=43 / d=11
%     sub-case handled explicitly. Every downstream use was already within
%     a Condition (II) context, so the reduction is unaffected.
%   - Theorem [Exact AP density] (S6.5): the q>3 restriction (present in
%     the hypothesis) is now enforced in the conclusion, the S1.2 summary,
%     and PPPC equivalent-form (c); the prime q=3 is routed via the LTE
%     theorem (3 | W_{p-2} iff 3 | p-2).
% Editorial corrections: a truncated sentence removed from S4.3; Theorem
% [Refined SFP mod 16](c) marked conditional on Conjecture R3; a one-
% directional rewording in the [Quartic residuosity] proof; an added
% prerequisite clause in Proposition [q-th power necessity]; the compatible-
% triple index range made consistent (d >= 1); a corrected factor count.
%
% v2.1 changes (retained): Section 9 gained a new closed lead in S9.1, the
% new subsubsection S9.2.5, and one bibliography entry (Dimitrov-Gao-
% Habegger); a cycle-16 proofread fixed typographic defects in S3-4.
\documentclass[11pt]{amsart}

\usepackage{amsmath}
\usepackage{amssymb}
\usepackage{amsthm}
\usepackage[T1]{fontenc}
\usepackage{hyperref}

\theoremstyle{plain}
\newtheorem{theorem}{Theorem}[section]
\newtheorem{proposition}[theorem]{Proposition}
\newtheorem{lemma}[theorem]{Lemma}
\newtheorem{corollary}[theorem]{Corollary}

\theoremstyle{definition}
\newtheorem{definition}[theorem]{Definition}
\newtheorem{conjecture}[theorem]{Conjecture}

\theoremstyle{remark}
\newtheorem{remark}[theorem]{Remark}
\newtheorem{observation}[theorem]{Observation}
\newtheorem*{remarknonum}{Remark}

\numberwithin{equation}{section}

% amsart stores the abstract in an unbreakable \vtop; with an abstract of
% this length the box overflows the page and the overflow is clipped (this
% silently truncated the last lines of the abstract in earlier versions).
% Emit the box with \unvbox so the abstract can break across pages; all
% abstract formatting is unchanged.
\makeatletter
\def\@setabstracta{%
  \ifvoid\abstractbox
  \else
    \skip@20\p@ \advance\skip@-\lastskip \advance\skip@-\baselineskip
    \vskip\skip@
    \unvbox\abstractbox
    \prevdepth\z@
  \fi
}
\makeatother

\newcommand{\ord}{\operatorname{ord}}
\newcommand{\rad}{\operatorname{rad}}
\newcommand{\Frob}{\operatorname{Frob}}
\newcommand{\Gal}{\operatorname{Gal}}
\newcommand{\Spec}{\operatorname{Spec}}
\newcommand{\Cov}{\operatorname{Cov}}
\newcommand{\Var}{\operatorname{Var}}
\newcommand{\Prob}{\operatorname{Pr}}
\newcommand{\Z}{\mathbb{Z}}
\newcommand{\Q}{\mathbb{Q}}
\newcommand{\F}{\mathbb{F}}
\newcommand{\eC}{\mathcal{C}}

\begin{document}

\title[Wagstaff Chebyshev reduction to SPC]{Reduction of the Wagstaff Chebyshev sufficiency conjecture to the Single-Pass Conjecture}

\author{Alexey Dolotov}
\email{science@dolotov.com}
\date{2026-07-16}

\subjclass[2020]{Primary 11Y11; Secondary 11A51, 11B39, 11R45, 11N36}
\keywords{Wagstaff numbers, Chebyshev polynomials, Pell sequences, primitive divisors, Lucas sequences, multiplicative orders, Chebotarev density}

\begin{abstract}
Let $W_p = (2^p + 1)/3$ for an odd prime $p \geq 5$, and let $\omega_3 = 3 + 2\sqrt{2} \in \Z[\sqrt 2]$. The congruence $\omega_3^{(W_p + 1)/2} \equiv -1 \pmod{W_p}$ --- the $a = 3$ case of the Chua $N + 1$ test --- is known to hold for every prime $W_p$; whether it suffices to \emph{imply} primality of $W_p$ is an open question, referred to here as the \emph{Wagstaff Chebyshev sufficiency conjecture} (Conjecture~\ref{conj:suff}) --- the Chua-form counterpart of the Vrba--Reix primality test for Wagstaff numbers~\cite{ref:VrbaReix}. We reduce this conjecture to a single structural statement about local passes of the test, and, in working form, to two pair-counting fragments of that statement together with one explicit integer.

Call a prime $r$ with $\ord_r(2) = 2p$ a \emph{local passer} at $p$ if $\omega_3^{(W_p+1)/2} \equiv -1$ in $\Z[\sqrt 2]/(r)$, and write $P_p$ for the set of local passers; every member of $P_p$ divides $W_p$. The \textbf{Single-Pass Conjecture (SPC)} asserts that $|P_p| \leq 1$ for every prime $p \geq 5$. The headline reduction (Theorem~\ref{thm:spc-reduction}) is that SPC implies the sufficiency conjecture; the implication is unconditional and carries no computational input --- its only ingredient from outside this paper is a Nagell--Ljunggren result (a composite $W_p$ is never a perfect power, hence has two distinct prime factors, and under a global Condition-(II) pass both would lie in $P_p$). When $W_p$ is prime, $P_p = \{W_p\}$ and SPC holds automatically at $p$; its entire content concerns composite $W_p$, where it asserts that at most one prime factor passes locally. Across every record of this paper, exactly four composite-$W_p$ exponents with $P_p \neq \emptyset$ are known; at each of them the passing factor is inert and $|P_p| = 1$ as far as computed, and no configuration with $|P_p| \geq 2$ has ever been observed.

The members of $P_p$ at composite $W_p$ are classified by the residue $r \bmod 8$ (the only two possibilities being $r \equiv 1$ and $r \equiv 3$). Split members $r \equiv 1 \pmod 8$ are exactly the \emph{compatible triples} of the Pell rank correspondence: their nonexistence (NCT) is proved unconditionally for every odd order parameter $d \leq 200$ (Theorem~\ref{thm:nct99} and complete primitive-part certificates) and at every parity-unblocked value in $(200, 400]$ --- hence for every odd $d \leq 400$ --- and no compatible triple, equivalently no split local pass at composite $W_p$, has ever been exhibited. Inert members $r \equiv 3 \pmod 8$ are exactly the danger triples of the inert branch: each is a primitive divisor of the Pell--Lucas number $V_{2 d_r}$ lying in the pinned arithmetic progression $r \equiv 1 \pmod{2p}$. SPC accordingly splits into three fragments: the inert-pair fragment, whose above-cutoff form is the \textbf{Pell Primitive Pair Conjecture} (PPPC: no prime $p$ admits two locally passing inert factors $r > 10^{12}$); the mixed fragment, the \emph{Mixed Pair Conjecture} (MPC: no $p$ admits a compatible triple and an inert local pass simultaneously), implied by NCT for general $d$; and the split--split fragment, likewise subsumed by NCT for general $d$. The certified perimeter comprises the Platinum enumeration --- a $684{,}965{,}381$-row enumeration of all inert factors $r \leq 10^{12}$, showing that each $p$ has at most one \emph{small} local-pass inert factor --- the complete danger census through $d = 400$ (exactly two real danger triples exist with $d \leq 400$, both at exponents closed by secondary factors), the pair-separation theorem ($8 p \gcd(d_{r_1}, d_{r_2}) \mid r_2 - r_1$ for two inert passers at one $p$), and the vanishing of the diagonal pair count for every admissible $d \leq 400$.

The working refinement (Theorem~\ref{thm:sharpened-chain}) proves Conjecture~\ref{conj:suff} from PPPC + MPC$^{>}$ + X1, where MPC$^{>}$ restricts the inert datum of MPC to $r > 10^{12}$ and X1 is the single-integer statement that Condition~(II) fails at some prime factor of $W_{10916765939}$; the implication takes two stated computational inputs, the Platinum enumeration and the witness discharges. The conjunction $\{$PPPC, MPC$^{>}$, X1$\}$ is in turn implied by SPC, given the same records. This is a conditional reduction, not an unconditional proof of the sufficiency conjecture.

This paper is a companion to \cite{ref:A}, which gives three unconditional ECPP-independent primality certificates (Brillhart--Lehmer--Selfridge with a Chua $N+1$ side) for $W_{2617}, W_{10501}, W_{12391}$.
\end{abstract}

\maketitle

\section{Introduction}

\subsection{The Wagstaff Chebyshev sufficiency conjecture}

For an odd prime $p$, the Wagstaff number is $W_p = (2^p + 1)/3$. The $+1$ companion of Mersenne, $W_p$ lacks a standalone primality criterion of Lucas--Lehmer's strength \cite{ref:Williams}. The most natural candidate is the base $a = 3$ case of the Chua framework \cite{ref:Chua}: set $\omega_3 = 3 + 2\sqrt{2} \in \Z[\sqrt 2]$ and ask whether the single congruence
\[
  \omega_3^{(W_p + 1)/2} \;\equiv\; -1 \pmod{W_p} \qquad \text{(Condition (II))}
\]
implies primality of $W_p$. One half of this is straightforward: if $W_p$ is prime, then Condition~(II) holds (Proposition~\ref{prop:necessity} below, a direct consequence of Chua's framework \cite{ref:Chua}). The converse is open.

\begin{conjecture}[Wagstaff Chebyshev sufficiency conjecture]\label{conj:suff}
No composite $W_p$ satisfies Condition~(II).
\end{conjecture}

The conjecture restates an open problem of some standing. Its recurrence form is the Chua iteration $S_{i+1} = S_i^2 - 2 \pmod{W_p}$ in the base $\omega_3$, whose return test encodes Condition~(II). It is closely related to --- but, as literally stated, distinct from --- the \emph{Vrba--Reix primality test} for Wagstaff numbers~\cite{ref:VrbaReix}: the latter iterates the same Chebyshev map started at $S_0 = 1/4 \pmod{2^p + 1}$, and this iteration lives in $\Q(\sqrt{-7})$ rather than $\Q(\sqrt 2)$ --- its base is $\theta = (1 + 3\sqrt{-7})/8 = (\bar\pi/\pi)^2$, where $2 = \pi\bar\pi$ splits in $\Q(\sqrt{-7})$. The two return conditions decouple at the level of individual prime factors, in both directions, and the auxiliary no-early-return hypothesis of the Vrba--Reix prize statement holds vacuously; the exact relation is worked out in the companion note~\cite{ref:VRNote}. The forward direction (prime $\Rightarrow$ the test passes) is classical for both forms; sufficiency is open for both, has been since 2008, and for the Vrba--Reix form carries a standing 500-euro reward for a published proof~\cite{ref:VrbaReix}. Both tests are implemented in standard software, agree with Condition~(II) empirically at every tested exponent, and Conjecture~\ref{conj:suff} is the sufficiency question for the Chua/Condition-(II) form --- the canonical order-forcing statement that this paper reduces.

Every known fully resolved composite $W_p$ tested in \S\ref{sec:inert-survey} is consistent with Conjecture~\ref{conj:suff}, and no global composite pass is known. The current record/probable-prime status for Wagstaff numbers is maintained externally \cite{ref:PrimePagesWagstaff} and is not used as an input to the reduction. The purpose of this paper is to reduce Conjecture~\ref{conj:suff} to a single structural statement about the prime factors themselves, the \emph{Single-Pass Conjecture}; its working refinement routes through the \emph{Pell Primitive Pair Conjecture} (\S\ref{sec:pppc-statement} and \S\ref{sec:pppc}) and the \emph{Mixed Pair Conjecture} (\S\ref{sec:mpc}). The main theorems are conditional reduction theorems; the unconditional sufficiency conjecture remains open.

For a prime $p \geq 5$ write $P_p$ for the set of primes $r$ with $\ord_r(2) = 2p$ at which Condition~(II) holds \emph{locally}, i.e.\ $\omega_3^{(W_p+1)/2} \equiv -1$ in $\Z[\sqrt 2]/(r)$ (Definition~\ref{def:local-pass}); every member of $P_p$ divides $W_p$ (Lemma~\ref{lem:pp-containment}).

\begin{conjecture}[Single-Pass Conjecture, SPC]\label{conj:spc-intro}
For every prime $p \geq 5$: $\;|P_p| \leq 1$.
\end{conjecture}

SPC is restated, with its supporting lemmas, as Conjecture~\ref{conj:spc} in \S\ref{sec:spc}, where Theorem~\ref{thm:spc-reduction} proves that it implies Conjecture~\ref{conj:suff} --- unconditionally, with no computational input. When $W_p$ is prime, $P_p = \{W_p\}$ and SPC holds automatically at $p$; its entire content concerns composite $W_p$, where it asserts that at most one prime factor passes Condition~(II) locally.

\subsection{Reduction chain (summary)}\label{sec:reduction-chain}

The reduction proceeds in three stages. Throughout, $r$ denotes a prime divisor of a hypothetical composite $W_p$ passing Condition~(II).

\textbf{Stage A (case analysis).} The mod-8 exclusion (Lemma~\ref{lem:mod8}) forces $r \equiv 1$ or $r \equiv 3 \pmod 8$. Each class is treated separately.

\begin{itemize}
\item \emph{Split case $r \equiv 1 \pmod 8$.} A descent to the split algebra $\F_r \times \F_r$ (Section~\ref{sec:split}) identifies an odd parameter $d \mid W_{p-2}$ with $\ord(\omega_3) = 4 d$. If Condition~(II) holds at $r$, Proposition~\ref{prop:rank-equiv} produces a compatible triple $(d,r,p)$; the \textbf{nonexistence of compatible triples} (NCT) rules out this situation. NCT is proved unconditionally for every odd $d \leq 200$: for $d \leq 99$ by Theorem~\ref{thm:nct99}, and at all nine parity-unblocked values $d = 107$, $121$, $129$, $131$, $139$, $163$, $171$, $177$, $179$ of $100 \leq d \leq 200$ by complete fixed-$d$ certificates (\S\ref{sec:nct-verify}); the same certificate closes all ten parity-unblocked values of $(200, 300]$ ($d = 201$, $209$, $211$, $227$, $243$, $249$, $251$, $281$, $283$, $297$) and all nine of $(300, 400]$ ($d = 307$, $321$, $331$, $347$, $361$, $363$, $379$, $387$, $393$) (\S\ref{sec:nct-verify}); NCT therefore holds unconditionally for every odd $d \leq 400$. The related Pell-rank divisibility Conjecture R3 is a stronger sufficient split-branch obstruction and holds at all 185 known split factors (\S\ref{sec:R3-verify}).
\item \emph{Inert case $r \equiv 3 \pmod 8$.} A descent in the inert algebra $\F_{r^2}$ (Section~\ref{sec:inert}) gives an odd parameter $g = \ord(\omega_3)/4$ dividing $G_r = \gcd((r+1)/4, W_{p-2})$; Condition~(II) at $r$ holds iff $g \mid W_{p-2}$. The \textbf{Pell-sequence exclusion} (Theorem~\ref{thm:pell-exclusion}) excludes $g \leq 43$ unconditionally; the \textbf{parity obstruction} (Theorem~\ref{thm:parity-obstruction}) blocks every $g$ with a non-Class III prime factor; together these close the inert branch at 845 of 869 known inert factors (\S\ref{sec:inert-survey}).
\end{itemize}

\textbf{Stage B (cross-case pinning, Section~\ref{sec:cross-case}).} Uniformly across all candidate composite $W_p$:

\begin{enumerate}
\item \textbf{Order-Pinning Lemma} (Theorem~\ref{thm:order-pinning}): every prime factor $r > 3$ of $W_p$ has $\ord_r(2) = 2p$ \emph{exactly}.
\item \textbf{Multi-Factor Pinning} (Theorem~\ref{thm:multi-pinning}): all prime factors $r > 3$ of $W_p$ share the pinned order parameter.
\item \textbf{Platinum Lemma} (Theorem~\ref{thm:platinum}): a complete enumeration of inert prime factors $r \leq 10^{12}$ shows that every prime $p$ has at most one such factor passing Condition~(II).
\item \textbf{Exact-AP characterization} (Theorem~\ref{thm:exact-ap}): for a prime $q \in A$ with $q > 3$, the divisibility $q \mid W_{p-2}$ holds iff $p \equiv m_q + 2 \pmod{2 m_q}$, where $m_q = \ord_q(2)/2$ and $A$ is the set of Class III primes. Prime powers require the valuation-adjusted LTE modulus recorded after the theorem; the prime $q = 3$ is governed by Theorem~\ref{thm:lte}.
\item \textbf{Witness discharges} (Proposition~\ref{prop:witness-discharges}): at two of the three exponents carrying a Platinum local-pass witness, an explicit second prime factor of $W_p$ fails Condition~(II), removing the exponent from every global-pass scenario; the third exponent is the content of the single-integer statement X1 (Conjecture~\ref{conj:x1}).
\end{enumerate}

\textbf{Stage C (residual).} The residual is the \emph{Single-Pass Conjecture} (SPC, Conjecture~\ref{conj:spc-intro} above; \S\ref{sec:spc}): no prime $p$ admits two distinct locally passing primes with $\ord_r(2) = 2p$. SPC implies Conjecture~\ref{conj:suff} unconditionally and with no computational input (Theorem~\ref{thm:spc-reduction}). A hypothetical double pass splits, by the residues mod $8$ of its two members, into three fragments: the \emph{inert pair}, whose above-cutoff form is the \emph{Pell Primitive Pair Conjecture} (PPPC, Conjecture~\ref{conj:pppc-intro} below) --- the principal residual; the \emph{mixed pair}, the \emph{Mixed Pair Conjecture} (MPC, \S\ref{sec:mpc}), implied by NCT for general $d$; and the \emph{split pair}, subsumed by NCT for general $d$ (Conjecture~\ref{conj:nct-general}), which forbids even a single split member at composite $W_p$. The working chain (Theorem~\ref{thm:sharpened-chain}, \S\ref{sec:reduction}) proves Conjecture~\ref{conj:suff} from PPPC + MPC$^{>}$ + X1 --- the above-cutoff pair fragments together with one explicit integer --- taking the Platinum enumeration (item 3) and the witness discharges (item 5) as stated computational inputs; the conjunction is implied by SPC, given the same records.

Table~\ref{tab:status} summarizes the algebraic and certified inputs, finite-range evidence, and remaining open part of each of the four residual conjectures.

\begin{table}[ht]
\centering
\footnotesize
\renewcommand{\arraystretch}{1.25}
\begin{tabular}{@{}p{2.45cm}p{3.05cm}p{4.15cm}p{2.2cm}@{}}
\hline
Residual & Algebraic / certified input & Finite-range evidence & Remaining gap \\
\hline
NCT, general $d$ \newline (split; subsumes the split--split fragment of SPC) & no compatible triple for any odd $d \leq 200$, nor at any parity-unblocked value of $(200, 400]$ (Thm~\ref{thm:nct99} + twenty-eight complete fixed-$d$ closures, \S\ref{sec:nct-verify}) & all $185$ known split factors satisfy R3 & general odd $d > 400$ \\
PPPC, $\Pi = 0$ \newline (inert pair, $r \equiv 3$) & Platinum small-factor bound (Thm~\ref{thm:platinum}, certified input); pair separation (Thm~\ref{thm:pair-separation}); diagonal $S_{\mathrm{diag}}(d) = 0$ for admissible $d \leq 400$ (Props~\ref{prop:diagonal-closures},~\ref{prop:danger-census-400}); inert floor $d_r > 400$ (Cor~\ref{cor:inert-floor-400}) & at most one local-pass inert factor $r \leq 10^{12}$ per $p$; no above-cutoff pair found in the $p$-window scan & general $p$, off-diagonal and diagonal $d > 400$ \\
MPC \newline (mixed pair) & implied by NCT-general (Prop~\ref{prop:nct-implies-mpc}): inherits every unconditionally proved NCT range, so a violation needs $d_s > 400$; datum-(ii) floor $d_{r_{\mathrm{in}}} > 400$ (Rem~\ref{rem:two-slices}) & no violating pair in any record; joint heuristic bracket (\S\ref{sec:heuristic}, heuristic only) & general $p$: a compatible triple and an inert pass at one $p$ \\
X1 \newline (single integer) & the other two witness exponents discharged at explicit failing factors (Prop~\ref{prop:witness-discharges}, certified) & directed scan $k \leq 10^{12}$ at the exponent finds no second factor & one explicit integer, $W_{10916765939}$ \\
\hline
\end{tabular}
\caption{State of the residual conjectures. Rows 2--4 (PPPC, MPC, X1) sit under the Single-Pass Conjecture (Conjecture~\ref{conj:spc-intro}): SPC implies each of them --- X1 given the recorded witness divisor. Conjecture~\ref{conj:suff} follows from rows 2--4, with MPC restricted to its above-cutoff form MPC$^{>}$, via Theorem~\ref{thm:sharpened-chain} (\S\ref{sec:reduction}), taking the Platinum enumeration and the witness discharges as stated computational inputs. Row 1 (NCT-general) is the split-branch context: it implies MPC (row 3) and subsumes the split--split fragment of SPC.}
\label{tab:status}
\end{table}

\subsection{The principal fragment: PPPC}\label{sec:pppc-statement}

The inert-pair fragment of SPC, above the Platinum cutoff, is the paper's principal residual.

\begin{conjecture}[Pell Primitive Pair Conjecture, PPPC]\label{conj:pppc-intro}
For every prime $p \geq 5$, at most one prime $r$ exists satisfying all of:
\[
  r \equiv 3 \pmod 8; \qquad \ord_r(2) = 2 p; \qquad r > 10^{12};
\]
\[
  d_r := \ord_r(\omega_3)/4 > 43; \qquad d_r \mid W_{p-2}.
\]
\end{conjecture}

Equivalently: $|T_p^{>}| \leq 1$ for every prime $p \geq 5$, where $T_p^{>}$ denotes the set of primes $r$ satisfying the five bullet conditions.

PPPC is genuinely Pell-theoretic. By Theorem~\ref{thm:pell-divisibility} and Lemma~\ref{lem:inert-primitivity}, any such $r$ is a primitive divisor of the Pell--Lucas number $V_{2 d_r}$, where $V_n = (1+\sqrt 2)^n + (1-\sqrt 2)^n$; by Order-Pinning it simultaneously lies in the arithmetic progression $r \equiv 1 \pmod{2p}$; by the Exact-AP characterization and its valuation-adjusted prime-power extension, the condition $d_r \mid W_{p-2}$ is a strict congruence condition on $p$. PPPC asserts that two such primitive divisors --- not necessarily attached to different values of $d_r$ --- cannot simultaneously be pinned to the same Wagstaff exponent $p$.

The name reflects three properties of the forbidden configuration:

\begin{itemize}
\item \textbf{Pell}: $r$ divides $V_{2 d}$, the Pell--Lucas sequence associated to the fundamental unit $\alpha = 1 + \sqrt 2$ of $\Z[\sqrt 2]$.
\item \textbf{Primitive}: $r$ is primitive for the $d$-indexed Pell--Lucas family, i.e., it divides $V_{2 d_r}$ but not $V_{2 d}$ for any $d < d_r$ (Lemma~\ref{lem:inert-primitivity}); Bilu--Hanrot--Voutier \cite{ref:BHV} is used only as general background on primitive divisors of Lucas sequences.
\item \textbf{Pair}: the conjecture forbids the \emph{coexistence} of two such primitive divisors at a single $p$; a single such divisor per $p$ is entirely compatible with PPPC and occurs in three identified ``local pass'' witnesses (\S\ref{sec:secondary}).
\end{itemize}

The Platinum enumeration of \S\ref{sec:inert-survey} verifies the small-factor input used by the reduction: at most one local-pass inert factor $r \leq 10^{12}$ occurs at any pinned exponent $p$ in the enumerated range. PPPC itself concerns the above-cutoff set $T_p^{>}$, and its unconditional resolution is the main open problem of this paper. Two unconditional constraints on the forbidden configuration are established below: the Pair Separation Theorem (\S\ref{sec:pair-separation}) forces $8 p \gcd(d_{r_1}, d_{r_2}) \mid r_2 - r_1$ for any two members of $T_p^{>}$, and the diagonal ($d_{r_1} = d_{r_2}$) contribution to the pair count vanishes for every admissible order parameter $d \leq 400$ (Propositions~\ref{prop:diagonal-closures} and~\ref{prop:danger-census-400}), in particular for every pair with both indices $(r_i - 1)/2p \leq 1604$ (Corollary~\ref{cor:diagonal-k500}).

With the principal fragment named, the headline reduction can be stated in full.

\begin{theorem}[Single-pass reduction]\label{thm:spc-intro}
SPC implies Conjecture~\ref{conj:suff}, unconditionally and with no computational input. Moreover SPC implies the conjunction $\{$PPPC, MPC$^{>}$, X1$\}$ of the working chain (Theorem~\ref{thm:sharpened-chain}) --- given, for the X1 component, the recorded witness divisor of $W_{10916765939}$ --- so every result of this paper conditional on that chain is a fortiori a consequence of SPC.
\end{theorem}

The first implication is proved as Theorem~\ref{thm:spc-reduction} (\S\ref{sec:spc}); the moreover clause is proved with Theorem~\ref{thm:sharpened-chain} (\S\ref{sec:reduction}).
Theorem~\ref{thm:sharpened-chain} proves Conjecture~\ref{conj:suff} from the conjunction PPPC + MPC$^{>}$ + X1 alone, taking the Platinum enumeration (Theorem~\ref{thm:platinum}) and the discharge verifications of Proposition~\ref{prop:witness-discharges} as stated computational inputs: unlike Theorem~\ref{thm:spc-reduction}, the working chain trades the split--split fragment and the sub-cutoff legs of SPC for these certified finite records.

\textbf{Unconditional contributions.} Independently of the conditional reduction, the following are proved with no conjectural input and constitute the technical core of the paper: the single-pass reduction itself (Theorem~\ref{thm:spc-reduction}: the implication SPC $\Rightarrow$ Conjecture~\ref{conj:suff}, proved unconditionally and with no computational input); the nonexistence of compatible triples for every odd $d \leq 200$ and at every parity-unblocked value in $(200, 400]$ (Theorem~\ref{thm:nct99} and complete fixed-$d$ certificates, \S\ref{sec:nct-verify}); the Order-Pinning Lemma, the Multi-Factor Pinning Theorem, and the multiplicity/Wieferich constraints of \S\ref{sec:wieferich} (Section~\ref{sec:cross-case}); the Pair Separation Theorem and the vanishing of the diagonal pair count for every admissible $d \leq 400$ (\S\ref{sec:pair-separation}, Proposition~\ref{prop:danger-census-400}); the complete danger census through $d = 400$ with the inert order-parameter floor $d_r > 400$ (Proposition~\ref{prop:danger-census-400}, Corollary~\ref{cor:inert-floor-400}); the quartic residuosity of $2$ at primitive Pell divisors (Theorem~\ref{thm:quartic}); and the Exact-AP characterization of the divisibility $q \mid W_{p-2}$ (Theorem~\ref{thm:exact-ap}).

\subsection{Relation to the companion paper}

\begin{itemize}
\item \textbf{\cite{ref:A}} Three unconditional Brillhart--Lehmer--Selfridge \cite{ref:BLS} primality proofs for $W_{2617}, W_{10501}, W_{12391}$, independent of ECPP \cite{ref:AtkinMorain}. Uses the same Chua $N + 1$ criterion (Proposition~\ref{prop:necessity} here, Proposition 2.4 of \cite{ref:A}) on the necessity side, but is otherwise disjoint from this paper: it does not use, and does not prove, Conjecture~\ref{conj:suff}.
\end{itemize}

This paper does not depend on \cite{ref:A}; the two may be read independently.

\subsection{Notation}\label{sec:notation}

Throughout, $p \geq 5$ denotes an odd prime and $W_p = (2^p + 1)/3$; $r$ denotes a prime divisor of $W_p$. For $x$ coprime to $m$, $\ord_m(x)$ is the multiplicative order modulo $m$; $v_q$ is the $q$-adic valuation. We work in $\Z[\sqrt 2]$ with fundamental unit $\alpha = 1 + \sqrt 2$ and Chebyshev base $\omega_3 = \alpha^2 = 3 + 2\sqrt 2$. The Pell sequences $(U_n), (V_n)$ satisfy $\alpha^n = V_n/2 + U_n \sqrt 2$ with the recurrence $X_{n+1} = 2 X_n + X_{n-1}$, $U_0 = 0$, $U_1 = 1$, $V_0 = V_1 = 2$; the identity $V_n^2 - 8 U_n^2 = 4(-1)^n$ recurs throughout.

For a prime $r \neq 2$ with $(2/r) = 1$, $\rho(r)$ is the Pell rank of apparition: the least $n \geq 1$ with $r \mid U_n$. For an inert prime $r$ with $r \mid W_p$ we write
\[
  G_r \;=\; \gcd\!\left(\frac{r+1}{4},\; W_{p-2}\right), \qquad d_r \;=\; g \;:=\; \ord_r(\omega_3)/4
\]
for the order parameter; we use $d_r$ when emphasising dependence on $r$ and $g$ in local proofs. $\Phi_d(x)$ is the $d$-th cyclotomic polynomial; $A = \{q \text{ odd prime} : v_2(\ord_q(2)) = 1\}$ is the set of \emph{Class III} primes (Definition~\ref{def:classIII}). For a prime $p \geq 5$, $P_p$ denotes the passer set of Definition~\ref{def:local-pass}: the primes $r$ with $\ord_r(2) = 2p$ that pass Condition~(II) locally at $p$.

\section{Preliminaries}\label{sec:prelim}

\subsection{Chua framework and the base \texorpdfstring{$a = 3$}{a = 3}}

For an integer $a \neq 0, \pm 1$, set $\omega_a = a + \sqrt{a^2 - 1}$ so that $\omega_a \bar\omega_a = 1$ and $\omega_a^n + \omega_a^{-n} = 2 T_n(a)$. The Chua identity is:

\begin{theorem}[Chua \cite{ref:Chua}]\label{thm:chua}
Let $Q$ be an odd prime with $\gcd(a^2 - 1, Q) = 1$. With $D = a^2 - 1$, $\varepsilon = (D/Q)$, $\delta = (2(a+1)/Q)$:
\[
  \omega_a^{(Q - \varepsilon)/2} \equiv \delta \pmod Q
\]
in $\Z[\sqrt{D}]/(Q)$.
\end{theorem}

For $a = 3$: $D = 8$, $2(a+1) = 8$, so the two Legendre symbols coincide and become $(8/Q) = (2/Q)^3 = (2/Q)$. This gives the $a = 3$ specialization:

\begin{proposition}[Base $a = 3$ necessity]\label{prop:necessity}
For every Wagstaff prime $W_p$ with $p \geq 5$: $\omega_3^{(W_p + 1)/2} \equiv -1 \pmod{W_p}$.
\end{proposition}

\begin{proof}
$W_p \equiv 3 \pmod 8$ (Lemma~\ref{lem:mod8} below) gives $(2/W_p) = -1$ by the second supplement, hence $\varepsilon = \delta = (8/W_p) = -1$; apply Theorem~\ref{thm:chua}.
\end{proof}

The base $a = 3$ is uniquely distinguished among integer Chebyshev bases by $a^2 - 1$ being a power of 2; $(a-1)(a+1) = 2^k$ forces $\{a - 1, a + 1\} = \{2, 4\}$ (Mihailescu \cite{ref:Mihailescu}). This is the source of the symmetric role played by $\omega_3$ in everything that follows.

\subsection{The order of 2 and the mod-8 dichotomy}

\begin{lemma}[Order of 2 at every factor]\label{lem:ord2}
For every prime $p \geq 5$ and every prime $r \mid W_p$: $\ord_r(2) = 2p$, $r \equiv 1 \pmod{2p}$.
\end{lemma}

\begin{proof}
$2^p \equiv -1 \pmod r$, so $\ord_r(2) \mid 2p$ but $\nmid p$. Since $p$ is prime, $\ord_r(2) \in \{2, 2p\}$. If $\ord_r(2)=2$, then $r\mid 3$. But $3\nmid W_p$ for prime $p\geq5$: otherwise $9\mid 2^p+1$, while $\ord_9(2)=6$ forces $p\equiv3\pmod6$, impossible. Hence $\ord_r(2)=2p$, and $r\equiv1\pmod{2p}$.
\end{proof}

\begin{lemma}[mod-8 classification]\label{lem:mod8}
Let $p \geq 5$ be prime. Then $W_p \equiv 3 \pmod 8$; every prime factor of $W_p$ is $\equiv 1$ or $3 \pmod 8$; and the total multiplicity of prime factors $\equiv 3 \pmod 8$ is odd.
\end{lemma}

\begin{proof}
$W_p \equiv 3 \pmod 8$: for $p \geq 5$, $2^p \equiv 0 \pmod{32}$, so $2^p + 1 \equiv 1 \pmod{32}$, i.e.\ $3 W_p \equiv 1 \pmod{32}$; since $3 \cdot 11 \equiv 1 \pmod{32}$, this gives $W_p \equiv 11 \pmod{32}$, hence $W_p \equiv 3 \pmod 8$ (consistent with Theorem~\ref{thm:mod16}). Factor classes: By Lemma~\ref{lem:ord2}, $v_2(\ord_r(2)) = 1$. Suppose $r \equiv 7 \pmod 8$: then $(2/r) = +1$ by the second supplement, so $2$ is a quadratic residue and $\ord_r(2) \mid (r-1)/2$; since $v_2(r-1) = 1$ here, $(r-1)/2$ is odd, so $\ord_r(2)$ is odd and $v_2(\ord_r(2)) = 0 \neq 1$. Suppose $r \equiv 5 \pmod 8$: then $(2/r) = -1$, so $2^{(r-1)/2} \equiv -1 \pmod r$ and $\ord_r(2) \nmid (r-1)/2$, forcing $v_2(\ord_r(2)) > v_2((r-1)/2) = v_2(r-1) - 1 = 1$; with $\ord_r(2) \mid r-1$ this gives $v_2(\ord_r(2)) = v_2(r-1) = 2 \neq 1$. Both classes are excluded. Parity: modulo 8 the product $W_p \equiv 3^{\sum'} \pmod 8$ where $\sum'$ counts factors $\equiv 3 \pmod 8$ with multiplicity, and $W_p \equiv 3 \pmod 8$ forces this sum odd.
\end{proof}

\begin{corollary}\label{cor:no57}
No composite $W_p$ passing Condition~(II) has a prime factor $r \equiv 5$ or $r \equiv 7 \pmod 8$; every prime factor is split ($r \equiv 1 \pmod 8$) or inert ($r \equiv 3 \pmod 8$).
\end{corollary}

\subsection{Refined mod-16 parity}

\begin{theorem}[Refined SFP mod 16]\label{thm:mod16}
For every odd prime $p \geq 5$:

(a) $W_p \equiv 11 \pmod{16}$.

(b) Every prime factor $r > 3$ of $W_p$ with $r \equiv 3 \pmod 8$ satisfies $r \equiv 3$ or $r \equiv 11 \pmod{16}$.

(c) Assume Conjecture~\ref{conj:R3}, and suppose $W_p$ is composite and passes Condition~(II) at every factor. Then by Theorem~\ref{thm:R3-eliminates} below every factor is inert; writing $t$ for the total number of inert factors (with multiplicity) and $k_{11}$ for the number of those $\equiv 11 \pmod{16}$:
\[
  t + 2 k_{11} \;\equiv\; 3 \pmod 4.
\]
\end{theorem}

\begin{proof}
(a) For odd prime $p \geq 5$: $3 W_p = 2^p + 1 \equiv 1 \pmod{32}$ (since $p \geq 5$ forces $2^p \equiv 0 \pmod{32}$). The inverse of 3 modulo 32 is 11 (as $3 \cdot 11 = 33 \equiv 1 \pmod{32}$), so $W_p \equiv 11 \pmod{32}$, hence $W_p \equiv 11 \pmod{16}$.

(b) By Lemma~\ref{lem:mod8}, $r \equiv 1$ or $3 \pmod 8$; reduce mod 16: $r \equiv 3 \pmod 8$ splits into $3$ or $11 \pmod{16}$.

(c) Under the hypotheses every factor is $\equiv 3 \pmod 8$; so $r_i \in \{3, 11\} \pmod{16}$. In $(\Z/16)^\times$, $\ord(3) = 4$ with $\{3^0, 3^1, 3^2, 3^3\} = \{1, 3, 9, 11\}$; so $11 \equiv 3^3 \pmod{16}$. Thus $\prod r_i \equiv 3^{k_3 + 3 k_{11}} \pmod{16}$. Using (a) and $11 \equiv 3^3$: $k_3 + 3 k_{11} \equiv 3 \pmod 4$; substitute $k_3 = t - k_{11}$.
\end{proof}

\begin{corollary}[Mod-16 sub-configurations for $t = 3$]\label{cor:mod16-t3}
In the setting of Theorem~\ref{thm:mod16}(c), if $t = 3$ then $k_{11} \in \{0, 2\}$.
\end{corollary}

\begin{remarknonum}
Theorem~\ref{thm:mod16} and Corollary~\ref{cor:mod16-t3} are structural observations on the mod-16 refinement of the all-inert configuration; they are not invoked in the reduction chain of \S\S\ref{sec:split}--\ref{sec:pppc}, and a reader interested only in the main result may skip them. They are recorded because the parity datum $t + 2 k_{11} \equiv 3 \pmod 4$ constrains any future enumeration of small-$t$ inert configurations.
\end{remarknonum}

\subsection{The \texorpdfstring{$v_2$}{v2} classification of \texorpdfstring{$\ord_r(\omega_3)$}{ord r(omega 3)}}

\begin{lemma}[$v_2$ classification]\label{lem:v2-class}
Let $r \geq 5$ be prime. Then
\[
  v_2(\ord_r(\omega_3)) \;=\; \begin{cases}
    2 & r \equiv 3 \pmod 8, \\
    1 & r \equiv 5 \pmod 8, \\
    0 & r \equiv 7 \pmod 8, \\
    \leq v_2(r - 1) - 1 & r \equiv 1 \pmod 8.
  \end{cases}
\]
\end{lemma}

\begin{proof}
Set $\alpha = 1 + \sqrt 2$, so $\omega_3 = \alpha^2$ and $v_2(\ord(\omega_3)) = \max(v_2(\ord(\alpha)) - 1, 0)$.

\emph{Split cases ($r \equiv 1, 7 \pmod 8$):} $(2/r) = 1$, so $\sqrt 2 \in \F_r$ and $\alpha \in \F_r^\ast$ with $\ord(\alpha) \mid r - 1$, so $v_2(\ord(\alpha)) \leq v_2(r - 1)$. For $r \equiv 7$: $v_2(r - 1) = 1$, giving $v_2(\ord(\omega_3)) \leq 0$. For $r \equiv 1$: $v_2(r - 1) \geq 3$, giving the bound.

\emph{Inert cases ($r \equiv 3, 5 \pmod 8$):} $\sqrt 2 \notin \F_r$, $\alpha \in \F_{r^2}^\ast$. Frobenius sends $\sqrt 2 \mapsto -\sqrt 2$, so $\alpha^r = \bar\alpha$, hence $\alpha^{r+1} = \alpha \bar\alpha = -1$. Thus $\ord(\alpha) \mid 2(r+1)$ but $\nmid (r+1)$, which forces $v_2(\ord(\alpha)) = v_2(r+1) + 1$. For $r \equiv 3$: $v_2(r+1) = 2$, so $v_2(\ord(\omega_3)) = 2$; for $r \equiv 5$: $v_2(r+1) = 1$, so $v_2(\ord(\omega_3)) = 1$.
\end{proof}

The exact value $v_2(\ord(\omega_3)) = 2$ at $r \equiv 3 \pmod 8$ is used in \S\ref{sec:gap-structure} to establish $\ord(\omega_3) = 4 g$ with $g$ odd; the upper bound at $r \equiv 1$ is used in \S\ref{sec:split}.

\subsection{The fundamental identity and coprimality}

\begin{lemma}\label{lem:fundamental}
For every $p \geq 5$: $W_p + 1 = 4\,W_{p-2}$, and $\gcd(W_p, W_{p-2}) = 1$.
\end{lemma}

\begin{proof}
$W_p + 1 = (2^p + 4)/3 = 4 (2^{p-2} + 1)/3 = 4 W_{p-2}$. For coprimality: $W_p - 4 W_{p-2} = -1$.
\end{proof}

\begin{lemma}[$p \nmid W_{p-2}$]\label{lem:pndivW}
$p \nmid W_{p-2}$ for every prime $p \geq 5$.
\end{lemma}

\begin{proof}
$4 W_{p-2} = W_p + 1 = (2^p + 4)/3 \equiv (2 + 4)/3 = 2 \pmod p$ by Fermat.
\end{proof}

These two lemmas are used repeatedly: Lemma~\ref{lem:fundamental} converts Condition~(II) into a divisibility statement on $W_{p-2}$, and Lemma~\ref{lem:pndivW} prevents certain orders from accommodating $p$ in their $W_{p-2}$-divisor.

\subsection{Two distinct prime factors}

The last preliminary is a consequence of the Nagell--Ljunggren literature: a composite $W_p$ always has at least two \emph{distinct} prime factors. This is the only place where the reduction of \S\ref{sec:spc} touches results outside this paper.

\begin{lemma}[$W_p$ is never a perfect power]\label{lem:w1-empty}
For every odd prime $p \geq 5$ there are no integers $y > 1$, $q \geq 2$ with $W_p = y^q$. In particular a composite $W_p$ has at least two distinct prime factors.
\end{lemma}

\begin{proof}
Suppose $W_p = y^q$ with integers $y > 1$, $q \geq 2$; we may take $q$ prime. For $p \geq 3$ one has $2^p + 1 \equiv 1 \pmod 8$, so $W_p \equiv 3^{-1} \equiv 3 \pmod 8$; since squares are $\equiv 0, 1, 4 \pmod 8$, this rules out $q = 2$, and $q \geq 3$ is odd. For odd $p$,
\[
  W_p \;=\; \frac{2^p + 1}{3} \;=\; \frac{(-2)^p - 1}{(-2) - 1},
\]
so $W_p = y^q$ is an instance of the Nagell--Ljunggren equation $(x^n - 1)/(x - 1) = y^q$ with $x = -2$, $n = p > 2$. By the consolidated statement of Bennett and Levin \cite[Proposition~1]{ref:BennettLevin} --- which incorporates the negative-$x$ theorem of Bugeaud and Mignotte \cite{ref:BugeaudMignotte2002} --- every solution with $|x|, |y|, q > 1$ and $n > 2$, other than the four known solutions
\[
  (x, y, n, q) \;\in\; \{(3, 11, 5, 2),\ (7, 20, 4, 2),\ (18, 7, 3, 3),\ (-19, 7, 3, 3)\},
\]
satisfies $|x| \geq 10^4$. Since $|{-2}| = 2 < 10^4$ and $x = -2$ appears in none of the four, no solution with $x = -2$ exists. (Two independent walls corroborate the instance: the same proposition also gives a prime divisor of $x$ congruent to $1 \bmod q$, and the only prime divisor of $-2$ is $2$, which excludes $x = -2$ without the $10^4$ computation; and Bugeaud, Cao and Mignotte~\cite{ref:BugeaudCaoMignotte2001} had already solved $(2^n + 1)/3 = y^q$ directly for exponents greater than~$1$.)
\end{proof}

\section{Condition (I) and the sufficiency conjecture}

\subsection{Condition (I) at every factor}

\emph{Condition~(I)} is the companion Euler congruence on the $N - 1$ side,
\[
  2^{(W_p - 1)/2} \;\equiv\; -1 \pmod{W_p},
\]
which together with Condition~(II) constitutes the two-sided $(N-1)/(N+1)$ primality test. We show that Condition~(I) holds automatically at every prime factor of $W_p$, so it is non-discriminating and all arithmetic content of the test rests in Condition~(II).

\begin{theorem}[Condition (I) at every factor]\label{thm:condI}
For every odd prime $p \geq 5$ and every prime $r \mid W_p$: $2^{(W_p - 1)/2} \equiv -1 \pmod r$.
\end{theorem}

\begin{proof}
By Lemma~\ref{lem:ord2}, $\ord_r(2) = 2p$. We show $(W_p - 1)/(2p)$ is odd: $(W_p - 1)/2 = (2^{p-1} - 1)/3$, and $(W_p - 1)/(2p) = (2^{p-1} - 1)/(3 p)$ is an integer (by Fermat and $p$ odd) and odd (denominator $3p$ odd; numerator $2^{p-1} - 1$ odd for $p \geq 2$). Therefore $2^{(W_p - 1)/2} = (2^p)^{(W_p - 1)/(2p)} \equiv (-1)^{\text{odd}} = -1$.
\end{proof}

\begin{corollary}[CRT consequence]\label{cor:condI-crt}
If $W_p$ is squarefree, then $2^{(W_p - 1)/2} \equiv -1 \pmod{W_p}$, i.e., Condition~(I) holds for $W_p$.
\end{corollary}

Squarefreeness is not known for general $p$; every argument below uses only the per-prime form of Theorem~\ref{thm:condI}. We record this consolidation:

\begin{corollary}\label{cor:condI-equiv}
Condition~(I) is automatic at every prime factor of $W_p$ and is therefore non-discriminating: the two-sided $(N-1)/(N+1)$ test reduces to Condition~(II) alone, which carries all of the arithmetic content of the sufficiency question.
\end{corollary}

\subsection{The two branches}

By Lemma~\ref{lem:mod8}, every prime factor of $W_p$ is $\equiv 1$ or $3 \pmod 8$; by Corollary~\ref{cor:no57} these are the only residue classes the case analysis need consider. We call $r \equiv 1 \pmod 8$ the \emph{split branch} (since $(2/r) = 1$ splits $\Z[\sqrt 2]/(r)$ as $\F_r \times \F_r$), and $r \equiv 3 \pmod 8$ the \emph{inert branch} (since $(2/r) = -1$ makes $\Z[\sqrt 2]/(r) = \F_{r^2}$).

Sections~\ref{sec:split} and~\ref{sec:inert} treat the two branches separately. Section~\ref{sec:cross-case} gives cross-case reductions that act uniformly on any composite $W_p$ regardless of its factor structure.

\section{Local passes and the Single-Pass Conjecture}\label{sec:spc}

\subsection{The passer set \texorpdfstring{$P_p$}{Pp}}

Condition~(II) is a congruence modulo $W_p$; all information it carries
about a hypothetical composite $W_p$ is visible at the prime factors.
This section isolates the per-prime object once and for all.

\begin{definition}[Local pass; the passer set $P_p$]\label{def:local-pass}
Let $p \geq 5$ be prime. A prime $r$ with $\ord_r(2) = 2p$ \emph{passes
Condition~(II) locally at $p$} if
\[
  \omega_3^{(W_p+1)/2} \;\equiv\; -1 \qquad \text{in } \Z[\sqrt 2]/(r).
\]
Write
\[
  P_p \;=\; \{\, r \text{ prime} \;:\; \ord_r(2) = 2p \text{ and } r
  \text{ passes Condition~(II) locally at } p \,\}.
\]
\end{definition}

Since $\ord_r(2)$ determines $p$, the sets $P_p$ are pairwise disjoint.
The apparent asymmetry of the definition --- the exponent $(W_p+1)/2$ is
global, the modulus is local --- is resolved by the first lemma: the
order hypothesis already forces $r$ to divide $W_p$.

\begin{lemma}[Containment]\label{lem:pp-containment}
Every $r \in P_p$ divides $W_p$. If $W_p$ is prime, then $P_p = \{W_p\}$.
\end{lemma}

\begin{proof}
From $\ord_r(2) = 2p$, the element $2^p$ has order $2$ in $\F_r^\times$,
so $2^p \equiv -1 \pmod r$ and $r \mid 2^p + 1 = 3\,W_p$. The case
$r = 3$ is impossible ($\ord_3(2) = 2 \neq 2p$ for $p \geq 5$), so
$r \mid W_p$. If $W_p$ is prime, the containment leaves only $r = W_p$;
and $W_p \in P_p$ indeed: $\ord_{W_p}(2) = 2p$ by Lemma~\ref{lem:ord2}
applied to the factor $r = W_p$, and Condition~(II) holds at $W_p$ by
Proposition~\ref{prop:necessity}.
\end{proof}

\begin{lemma}[Restriction to prime factors]\label{lem:restriction}
Let $p \geq 5$ be prime and suppose $W_p$ satisfies Condition~(II). Then
every prime factor of $W_p$ lies in $P_p$; in particular $|P_p|$ is at
least the number of distinct prime factors of $W_p$.
\end{lemma}

\begin{proof}
The congruence $\omega_3^{(W_p+1)/2} \equiv -1 \pmod{W_p}$ restricts
along the ring surjection $\Z[\sqrt 2]/(W_p) \to \Z[\sqrt 2]/(r)$ for
every prime $r \mid W_p$, and $\ord_r(2) = 2p$ by Lemma~\ref{lem:ord2}.
\end{proof}

\subsection{The conjecture and the reduction}

\begin{conjecture}[Single-Pass Conjecture, SPC]\label{conj:spc}
For every prime $p \geq 5$: $\;|P_p| \leq 1$.
\end{conjecture}

In words: a local pass of Condition~(II) never has company --- at each
exponent, at most one prime in the pinned progression
$r \equiv 1 \pmod{2p}$ passes. By Lemma~\ref{lem:pp-containment} the
conjecture holds automatically, with $|P_p| = 1$, whenever $W_p$ is
prime; its entire content concerns composite $W_p$, where it asserts
that at most one prime factor passes locally.

\begin{theorem}[Single-pass reduction]\label{thm:spc-reduction}
SPC implies Conjecture~\ref{conj:suff}. The implication is unconditional
and carries no computational input.
\end{theorem}

\begin{proof}
Suppose a composite $W_p$ satisfies Condition~(II). By
Lemma~\ref{lem:restriction}, every prime factor of $W_p$ lies in $P_p$.
By Lemma~\ref{lem:w1-empty}, $W_p$ is not a perfect power, so, being
composite, it has two distinct prime factors $r_1 \neq r_2$. Then
$\{r_1, r_2\} \subseteq P_p$ and $|P_p| \geq 2$, contradicting SPC.
\end{proof}

\begin{remark}[What the reduction actually consumes]\label{rem:spc-strict}
SPC is deliberately slightly stronger than what
Theorem~\ref{thm:spc-reduction} needs. Every composite $W_p$ has an
inert prime factor $r \equiv 3 \pmod 8$ (Lemma~\ref{lem:mod8}: the inert
multiplicity is odd), so in the proof one of $r_1, r_2$ may always be
taken inert: Conjecture~\ref{conj:suff} already follows from the
fragments of SPC in which at least one member of the forbidden pair is
inert --- the Mixed Pair Conjecture and the inert-pair statement
developed below (\S\ref{sec:mpc}, \S\ref{sec:pppc};
Proposition~\ref{prop:branch-split}, \S\ref{sec:reduction}).
The split--split fragment ---
no two distinct split members of $P_p$ --- is included in SPC for
uniformity and falsifiability, not because the reduction uses it; it is
subsumed by NCT for general $d$ (Conjecture~\ref{conj:nct-general}),
which forbids even a single split member at composite $W_p$. In the
converse direction, Conjecture~\ref{conj:suff} does \emph{not} imply
SPC: a composite $W_p$ with two locally passing factors and one failing
factor would violate SPC while leaving Conjecture~\ref{conj:suff}
untouched. SPC is proposed as the correct single statement to test and
attack: every new factorization of a Wagstaff number tests it, and both
principal residuals of this paper are fragments of it.
\end{remark}

\begin{remark}[Known nonempty $P_p$ at composite $W_p$]\label{rem:known-passers}
Across every record of this paper --- the Platinum enumeration of all
inert factors $r \leq 10^{12}$ (Theorem~\ref{thm:platinum},
\S\ref{sec:secondary}), the complete danger census through $d = 400$
(Proposition~\ref{prop:danger-census-400}), the NCT certificates for
every odd $d \leq 400$ (\S\ref{sec:nct}, \S\ref{sec:nct-verify}), and
the $p$-indexed window scans (\S\ref{sec:empirical}) --- exactly four
composite-$W_p$ exponents with $P_p \neq \emptyset$ are known:

\begin{center}
\footnotesize
\setlength{\tabcolsep}{2pt}
\begin{tabular}{@{}llll@{}}
\hline
$p$ & passing factor $r$ & $d_r$ & second factor \\
\hline
$1{,}251{,}488{,}009$ & $2p+1$ & $41{,}716{,}267$ & fails (discharged) \\
$10{,}916{,}765{,}939$ & $152{,}834{,}723{,}147$ & $38{,}208{,}680{,}787$ & sought: X1 \\
$85{,}684{,}865{,}933$ & $2p+1$ & $57$ & fails (discharged) \\
$24{,}822{,}855{,}876{,}938{,}710{,}967$ & $148{,}937{,}135{,}261{,}632{,}265{,}803$ & $67$ & fails (v) \\
\hline
\end{tabular}
\end{center}

All four passing factors are inert; no split local pass --- equivalently,
no compatible triple (\S\ref{sec:split}) --- has ever been exhibited. In
each of the four cases $|P_p| = 1$ as far as computed: at three of the
four exponents an explicit second prime factor of $W_p$ is known and
\emph{fails} Condition~(II) (Proposition~\ref{prop:witness-discharges},
Corollary~\ref{cor:inert-floor-400}, and \S\ref{sec:secondary}), and at
the fourth a directed scan finds no
second factor to date (Conjecture~\ref{conj:x1}). No configuration with
$|P_p| \geq 2$ --- at any size, on any branch --- has ever been observed.
\end{remark}

\begin{remark}[Plan of the paper, restated]\label{rem:spc-plan}
The remainder of the paper is the structure theory and the perimeter of
SPC. Sections~\ref{sec:split} and~\ref{sec:inert} classify the members
of $P_p$ at composite $W_p$ by the residue $r \bmod 8$: split members
are exactly the compatible triples, inert members are exactly the danger
triples, and both inherit rigid order data. Section~\ref{sec:cross-case}
proves the pinning and separation constraints and states the certified
computational records (the Platinum enumeration, the danger census ---
stated on the inert side, Proposition~\ref{prop:danger-census-400} ---
and the witness discharges). Section~\ref{sec:pppc} develops the
principal fragment PPPC with its empirical and heuristic support.
Section~\ref{sec:reduction} proves the working refinement
(Theorem~\ref{thm:sharpened-chain}): the certified records reduce SPC's
role to its above-cutoff fragments, and Conjecture~\ref{conj:suff}
follows from PPPC + MPC$^{>}$ + X1 alone.
\end{remark}

\section{The split branch: \texorpdfstring{$r \equiv 1 \pmod 8$}{r = 1 mod 8}}\label{sec:split}

\subsection{Descent in the split algebra}

For $r \equiv 1 \pmod 8$ with $r \mid W_p$: $(2/r) = 1$, so $\Z[\sqrt 2]/(r) \cong \F_r \times \F_r$ via $\sqrt 2 \mapsto (s, -s)$ where $s^2 = 2$ in $\F_r$. The element $\omega_3$ maps to $(\alpha_0, \alpha_0^{-1})$ with $\alpha_0 = 3 + 2 s$.

\textbf{Descent.} Condition~(II) at $r$ becomes $\alpha_0^{(N+1)/2} = -1$ in $\F_r^\ast$. Therefore $\ord(\alpha_0) \mid N + 1 = 4 W_{p-2}$ and $v_2(\ord(\alpha_0)) = v_2(N + 1) = 2$. Write $\ord(\alpha_0) = 4 d$ with $d$ odd and $d \mid W_{p-2}$.

\begin{theorem}[Order of $\sqrt 2$]\label{thm:ord-sqrt2}
Let $r \equiv 1 \pmod 8$, $r \mid W_p$. Then $\ord_{\F_r^\ast}(s) = 4p$.
\end{theorem}

\begin{proof}
$s^{2p} = 2^p \equiv -1 \pmod r$, so $\ord(s) \nmid 2p$; $s^{4p} = 1$, so $\ord(s) \mid 4p$. For prime $p$ the only divisor not dividing $2p$ is $4$ or $4p$; $\ord(s) = 4$ gives $4 \equiv 1 \pmod r$, impossible for $r \geq 5$.
\end{proof}

\begin{proposition}[Self-supporting exponents]\label{prop:ladder}
Let $p \geq 5$ be prime and let $W_{p-2}$ be prime. If $W_p$ satisfies Condition~(II), then $W_p$ is prime.
\end{proposition}

\begin{proof}
Suppose $W_p$ is composite and let $r < W_p$ be a prime factor; Condition~(II) restricts along $\Z[\sqrt2]/(W_p) \to \Z[\sqrt2]/(r)$, so it holds at $r$. If $r \equiv 1 \pmod 8$, the descent above gives $\ord(\alpha_0) = 4d$ with $d$ odd and $d \mid W_{p-2}$; primality of $W_{p-2}$ leaves $d \in \{1,\, W_{p-2}\}$. For $d = 1$: $\alpha_0^2 = -1$ forces $r \mid \mathrm N(\omega_3^2 + 1) = \mathrm N(18 + 12\sqrt 2) = 36$, impossible for $r \equiv 1 \pmod 8$. For $d = W_{p-2}$: $\ord(\alpha_0) = 4 W_{p-2} = W_p + 1$ divides $r - 1$, so $r > W_p$ --- impossible. If $r \equiv 3 \pmod 8$, the inert descent (Section~\ref{sec:inert}) gives $\ord(\omega_3) = 4g$ at $r$ with $g$ odd and $g \mid W_{p-2}$, and the Pell-sequence exclusion (Theorem~\ref{thm:pell-exclusion}) gives $g > 43$; so $g = W_{p-2}$ and $W_p + 1 = 4 W_{p-2} = \ord(\omega_3) \leq r + 1$, i.e.\ $r \geq W_p$ --- again impossible for a proper factor.
\end{proof}

\begin{corollary}[The certified ladder]\label{cor:ladder}
Condition~(II) is an unconditional primality test for $W_p$ at every prime $p$ such that $p - 2$ is also prime and $W_{p-2}$ carries a primality certificate: presently
\[
  p \;\in\; \{\,5,\ 7,\ 13,\ 19,\ 103,\ 193,\ 349,\ 3541,\ 83{,}341,\ 117{,}241,\ 127{,}033\,\},
\]
the complete list arising from the certified Wagstaff primes: the three largest rungs rest on the ECPP certificates of $W_{83{,}339}$, $W_{117{,}239}$ and $W_{127{,}031}$ (completed in 2014, 2022 and 2023 respectively) \cite{ref:PrimePagesWagstaff}, and for every other certified Wagstaff exponent $q$ the shift $q + 2$ is composite.
\end{corollary}

\begin{remark}[BLS specialization; the ladder is self-supporting]\label{rem:ladder}
Proposition~\ref{prop:ladder} is the fully-factored-$(N+1)$ specialization of Brillhart--Lehmer--Selfridge \cite{ref:BLS}, via $W_p + 1 = 4 W_{p-2}$. What the present machinery adds is that Condition~(II) \emph{alone} completes it: the $d = 1$ branch dies on the norm $\mathrm N(\omega_3^2 + 1) = 36$ and the inert branch on the Pell-sequence exclusion --- no auxiliary base, no extra congruence. Under a twin-prime-type heuristic the ladder $\{p : p,\ p-2,\ W_{p-2} \text{ all prime}\}$ is infinite, and each new Wagstaff primality certificate at an exponent $q$ with $q + 2$ prime adds a rung (enumeration: \texttt{scripts/cp410\_twin\_ladder.py}). In the self-supporting reading, sufficiency of Condition~(II) --- Conjecture~\ref{conj:suff} --- would propagate primality \emph{up} the ladder without any factoring at all.
\end{remark}

\subsection{Conjecture R3 and its conditional reduction}

\begin{remark}\label{rem:alpha0-square}
Since $\alpha_0 = (1 + s)^2$: $\ord(\alpha_0) = \ord(1 + s) / \gcd(\ord(1+s), 2)$. The key divisibility is $p \mid \ord(1 + s)$.
\end{remark}

\begin{conjecture}[Pell-rank divisibility, R3]\label{conj:R3}
For every prime $r \equiv 1 \pmod 8$ with $r \mid W_p$: $p \mid \ord_{\F_r^\ast}(1 + s)$, where $s = \sqrt 2 \bmod r$.
\end{conjecture}

\begin{theorem}[Conditional on R3]\label{thm:cond-R3}
Let $W_p$ be composite with a prime factor $r \equiv 1 \pmod 8$. If Conjecture R3 holds at $r$, then $r$ cannot satisfy Condition~(II).
\end{theorem}

\begin{proof}
By R3: $p \mid \ord(1 + s)$; since $\alpha_0 = (1+s)^2$, $p \mid \ord(\alpha_0) = 4d$, so $p \mid d \mid W_{p-2}$, contradicting Lemma~\ref{lem:pndivW}.
\end{proof}

\begin{theorem}[Split branch eliminates under R3]\label{thm:R3-eliminates}
If Conjecture R3 holds at every split factor of every composite $W_p$, then every composite $W_p$ passing Condition~(II) has all factors inert ($\equiv 3 \pmod 8$).
\end{theorem}

\begin{proof}
Immediate from Theorem~\ref{thm:cond-R3} and Corollary~\ref{cor:no57}.
\end{proof}

\subsection{Proved cases of Conjecture R3}

\begin{proposition}[Algebraic identity]\label{prop:alg-identity}
In $\Z[\sqrt 2]$: $\Phi_{2p}(\omega_3) = U_p \cdot (1 + \sqrt 2)^{p-1}$. Consequently $\mathrm{N}(\Phi_{2p}(\omega_3)) = U_p^2$.
\end{proposition}

\begin{proof}
Since $p$ is an odd prime, $\Phi_{2p}(x) = (x^p + 1)/(x + 1)$; hence
\[
  \Phi_{2p}(\omega_3) \;=\; \frac{\omega_3^p + 1}{\omega_3 + 1}.
\]
Using $\omega_3 = \alpha^2$ with $\alpha = 1 + \sqrt 2$ and the Pell formula $\alpha^p = V_p/2 + U_p \sqrt 2$: $\omega_3^p = \alpha^{2p} = (V_p/2 + U_p \sqrt 2)^2 = (V_p^2/4 + 2 U_p^2) + V_p U_p \sqrt 2$. Using $V_p^2 - 8 U_p^2 = -4$ (so $V_p^2/4 = 2 U_p^2 - 1$), the rational part simplifies:
\[
  \begin{aligned}
  \omega_3^p &= (4 U_p^2 - 1) + V_p U_p \sqrt 2,\\
  \omega_3^p + 1 &= 4 U_p^2 + V_p U_p \sqrt 2
                 = U_p (4 U_p + V_p \sqrt 2).
  \end{aligned}
\]
Also $\omega_3 + 1 = 4 + 2\sqrt 2 = 2 \alpha \cdot \alpha^{-1} \cdot (2 + \sqrt 2) = 2 (2 + \sqrt 2)$. Dividing:
\[
  \Phi_{2p}(\omega_3) \;=\; \frac{U_p (4 U_p + V_p \sqrt 2)}{2(2 + \sqrt 2)} \;=\; U_p \cdot f_p, \qquad f_p := \frac{4 U_p + V_p \sqrt 2}{4 + 2 \sqrt 2}.
\]
The Pell recurrences $U_{n+1} = 2 U_n + U_{n-1}$, $V_{n+1} = 2 V_n + V_{n-1}$ give $f_{n+1} = 2 f_n + f_{n-1}$, and by direct computation $f_0 = 2\sqrt 2/(4 + 2\sqrt 2) = \sqrt 2 - 1 = \alpha^{-1}$, $f_1 = (4 + 2\sqrt 2)/(4 + 2\sqrt 2) = 1 = \alpha^0$. The same recurrence is satisfied by $\alpha^{n-1}$ (since $\alpha^2 = 2\alpha + 1$), so $f_n = \alpha^{n-1}$ by induction. Hence $\Phi_{2p}(\omega_3) = U_p \alpha^{p-1} = U_p (1+\sqrt 2)^{p-1}$. Norm: $\mathrm{N}((1+\sqrt 2)^{p-1}) = (-1)^{p-1} = 1$ since $p$ is odd; so $\mathrm{N}(\Phi_{2p}(\omega_3)) = U_p^2$.
\end{proof}

\begin{theorem}[Case A]\label{thm:caseA}
If $r \equiv 1 \pmod 8$, $r \mid W_p$, and $r \mid U_p$, then $\ord(\alpha) = 4p$ in the split algebra.
\end{theorem}

\begin{proof}
Reducing Proposition~\ref{prop:alg-identity} modulo $r$: $r \mid U_p$ gives $\Phi_{2p}(\omega_3) \equiv 0 \pmod r$ (since $\Phi_{2p}(\omega_3) = U_p \cdot (1 + \sqrt 2)^{p-1}$ and $r \nmid (1 + \sqrt 2)^{p-1}$ in the split algebra, as $1 + \sqrt 2$ is a unit). So $\omega_3 \bmod r$ is a primitive $2p$-th root of unity in $\F_r^\ast$: $\ord(\omega_3) = 2 p$.

Now $\omega_3 = \alpha^2$, so $\alpha^{4p} = \omega_3^{2p} = 1$ and $\ord(\alpha) \mid 4p$. Since $\ord(\omega_3) = 2p$, we have $\omega_3^p = -1$ (the unique element of order 2 in the cyclic group $\langle\omega_3\rangle$), so $\alpha^{2p} = \omega_3^p = -1$, hence $\ord(\alpha) \nmid 2p$. The divisors of $4p$ not dividing $2p$ are $4$ and $4p$, and $\ord(\alpha) = 4$ is impossible since $\alpha^4 = \omega_3^2$ has order $p \geq 5$. Therefore $\ord(\alpha) = 4p$.
\end{proof}

\begin{theorem}[Case B]\label{thm:caseB}
If $r \equiv 1 \pmod 8$, $r \mid W_p$, and $r \mid V_p$, then $\ord(\alpha) = 2p$ in the split algebra.
\end{theorem}

\begin{proof}
With the split embedding $\sqrt 2 \mapsto (s, -s)$ and $r \mid V_p$, the identity $\alpha^p = V_p/2 + U_p \sqrt 2$ gives
\[
  \alpha^p \;\equiv\; (U_p s,\; -U_p s) \quad\text{in } \F_r \times \F_r.
\]
Squaring coordinatewise: $\alpha^{2p} \equiv (2 U_p^2,\; 2 U_p^2)$. From the Pell identity $V_p^2 - 8 U_p^2 = -4$ with $r \mid V_p$: $8 U_p^2 \equiv 4 \pmod r$, so $2 U_p^2 \equiv 1 \pmod r$. Hence $\alpha^{2p} \equiv (1, 1) = 1$.

Meanwhile $\alpha^p \equiv (U_p s, -U_p s)$ has coordinates of opposite signs, so $\alpha^p \not\equiv (1, 1)$. Thus $\ord(\alpha) \mid 2p$ and $\ord(\alpha) \nmid p$, forcing $\ord(\alpha) \in \{2, 2p\}$. If $\ord(\alpha) = 2$: $\omega_3 = \alpha^2 \equiv (1, 1)$, i.e., $3 + 2 s \equiv 1$ and $3 - 2 s \equiv 1$ in $\F_r$; adding: $6 \equiv 2$, i.e., $4 \equiv 0 \pmod r$, contradicting $r \geq 17$. Therefore $\ord(\alpha) = 2p$.
\end{proof}

\begin{theorem}[Case $4 U_p^2 \equiv 1$]\label{thm:caseC}
If $r \equiv 1 \pmod 8$, $r \mid W_p$, and $4 U_p^2 \equiv 1 \pmod r$, then $\ord(\alpha) = 8p$ in the split algebra.
\end{theorem}

\begin{proof}
With $4 U_p^2 \equiv 1 \pmod r$: compute $\alpha^{2p}$ using the identity $\alpha^p = V_p/2 + U_p \sqrt 2$. Squaring:
\[
  \alpha^{2p} \;=\; V_p^2/4 + V_p U_p \sqrt 2 + 2 U_p^2 \;=\; (V_p^2/4 + 2 U_p^2) + V_p U_p \sqrt 2.
\]
From $V_p^2 - 8 U_p^2 = -4$: $V_p^2/4 + 2 U_p^2 = V_p^2/4 + (V_p^2 + 4)/4 = V_p^2/2 + 1$. Substituting: $\alpha^{2p} = (V_p^2/2 + 1) + V_p U_p \sqrt 2 \equiv V_p^2/2 + 1 + V_p U_p s$ in the split algebra.

Now use $4 U_p^2 \equiv 1$: $U_p^2 \equiv 1/4$; and $V_p^2 \equiv 8 U_p^2 - 4 \equiv 2 - 4 = -2$; so $V_p^2/2 + 1 \equiv -1 + 1 = 0$. Also $V_p U_p$ has square $V_p^2 U_p^2 \equiv (-2)(1/4) = -1/2$. Hence $\alpha^{2p} \equiv V_p U_p s$, with $(\alpha^{2p})^2 = (V_p U_p)^2 \cdot s^2 = (-1/2) \cdot 2 = -1$. So $\alpha^{4p} \equiv -1 \pmod r$.

Therefore $\ord(\alpha) \mid 8p$ but $\ord(\alpha) \nmid 4p$; the divisors of $8p$ not dividing $4p$ are $8$ and $8p$. If $\ord(\alpha) = 8$: $\alpha^8 = (1 + \sqrt 2)^8 = 577 + 408\sqrt 2 \equiv 1 \pmod r$ would force $r \mid \gcd(576, 408) = 24$, contradicting $r \geq 17$. Hence $\ord(\alpha) = 8p$.
\end{proof}

In each case $p \mid \ord(\alpha)$, confirming Conjecture R3 (via Proposition~\ref{prop:rank-equiv} below). Among the 185 known split factors, these three cases cover 29\% (1 + 6 + 46 factors respectively); see \S\ref{sec:R3-verify}.

\subsection{Pell rank and equivalence with Conjecture R3}

\begin{definition}[Pell rank of apparition]\label{def:pell-rank}
For a prime $r$ with $(2/r) = 1$, the Pell rank $\rho(r)$ is the least $n \geq 1$ with $r \mid U_n$. We say $r$ is a primitive divisor of $U_n$ if $\rho(r) = n$. Bilu--Hanrot--Voutier \cite{ref:BHV} gives primitive divisors for all sufficiently large Lucas--Lehmer indices; the small indices used in the finite checks below are handled directly.
\end{definition}

\begin{proposition}[Rank equivalence]\label{prop:rank-equiv}
Let $r \equiv 1 \pmod 8$, $r \mid W_p$. Suppose Condition~(II) holds at $r$ in a composite $W_p$; so $\ord(\alpha_0) = 4 d$ with $d$ odd. Then:

(a) $\rho(r) = 4 d$ and $r$ is a primitive divisor of $U_{4 d}$.

(b) Conjecture R3 at $r$ holds iff $p \mid \rho(r)$.
\end{proposition}

\begin{proof}
Set $\beta=1+s$, so $\beta^2=\alpha_0$. We first establish $\ord(\beta) = 8 d$. Since $\beta^2 = \alpha_0$,
\[
  \ord(\alpha_0) = \frac{\ord(\beta)}{\gcd(\ord(\beta), 2)}.
\]
An odd $\ord(\beta)$ would force $\ord(\alpha_0)$ odd, contradicting $\ord(\alpha_0) = 4 d$, which has $v_2 = 2$. Hence $\ord(\beta) = 2\,\ord(\alpha_0) = 8 d$.

\emph{(a)} From $\ord(\beta) = 8 d$: $\beta^{4 d} = -1$ (the unique order-2 element in $\langle\beta\rangle$, since $4 d = \ord(\beta)/2$). Now $\beta^{2n} = \alpha_0^n = (3 + 2 s)^n$; using the identification of Pell polynomials with Chebyshev, one obtains
\[
  \beta^{2n} \;=\; (4 U_n^2 + (-1)^n) + V_n U_n s \qquad\text{in } \F_r.
\]
(Direct verification: for $n = 1$, $\beta^2 = 1 + 2s + s^2 = 3 + 2s = 4 U_1^2 - 1 + V_1 U_1 s = 4 - 1 + 2 s = 3 + 2 s$ \checkmark. Induction on $n$ then propagates the formula, using the recurrences $U_{n+1} = 2 U_n + U_{n-1}$, $V_{n+1} = 2 V_n + V_{n-1}$ together with the Pell cross relation
\[
  V_n \;=\; 2\,(U_n + U_{n-1}),
\]
which the inductive step uses. This relation is obtained by comparing the $\sqrt 2$-coefficients on the two sides of $\alpha^{n} = \alpha \cdot \alpha^{n-1}$: from $\alpha^{n} = V_{n}/2 + U_{n}\sqrt 2$ and $(1+\sqrt 2)(V_{n-1}/2 + U_{n-1}\sqrt 2)$ one reads off $U_{n} = V_{n-1}/2 + U_{n-1}$, i.e. $V_{n-1} = 2(U_{n} - U_{n-1})$; shifting the index and using $U_{n+1} = 2 U_n + U_{n-1}$ gives $V_{n} = 2(U_{n+1} - U_{n}) = 2(U_n + U_{n-1})$. It is this identity that converts the cross term in $\beta^{2(n+1)}$ into the required $V_{n+1} U_{n+1}$ coefficient.)

So $\beta^{2 n} = (-1)^n$ iff the constant term satisfies $4 U_n^2 + (-1)^n \equiv (-1)^n \pmod r$ and the $s$-coefficient $V_n U_n$ vanishes mod $r$. The constant-term condition reads $4 U_n^2 \equiv 0 \pmod r$, hence $r \mid U_n$ (as $r$ is odd); the $s$-coefficient condition $V_n U_n \equiv 0$ is then automatic.

For $n = 4 d$: $\beta^{8 d} = 1 = (-1)^{4 d}$, so $r \mid U_{4 d}$. For any integer $1 \leq n < 4 d$: if $n$ is even, $\beta^{2 n} = 1$ requires $8 d \mid 2 n$, so $4 d \mid n$, contradicting $n < 4 d$. If $n$ is odd, $\beta^{2 n} = -1$ requires $\beta^{2 n} = \beta^{4 d}$, i.e., $8 d \mid 2 n - 4 d$, i.e., $2 d \mid n$, impossible for odd $n$. So $\rho(r) = 4 d$, and $r$ is a primitive divisor of $U_{4 d}$.

\emph{(b)} Conjecture R3 at $r$ says $p \mid \ord(1 + s) = \ord(\beta) = 8 d$. Since $p$ is odd, $p \mid 8 d$ iff $p \mid d$, iff $p \mid 4 d = \rho(r)$.
\end{proof}

\begin{definition}[Compatible triple]\label{def:compat-triple}
A compatible triple $(d, r, p)$ is: an odd integer $d \geq 1$; a prime $r \equiv 1 \pmod 8$ that is a primitive divisor of $U_{4d}$; a prime $p \geq 5$ with $\ord_r(2) = 2p$ and $d \mid W_{p-2}$.
\end{definition}

By Proposition~\ref{prop:rank-equiv}, if Condition~(II) passes at some split factor of a composite $W_p$, a compatible triple exists.

\begin{remark}[No phantom compatible triples]\label{rem:no-phantom-split}
Every compatible triple has $W_p$ composite: $r \mid W_p$ with $r \equiv 1 \pmod 8$, while $W_p \equiv 3 \pmod 8$ (Lemma~\ref{lem:mod8}), so $r \neq W_p$ and $r$ is a \emph{proper} divisor. Hence the nonexistence of compatible triples coincides with its restriction to composite $W_p$ --- the only form the reduction of \S\ref{sec:reduction} invokes; no separate ``$W_p$ prime'' case arises on the split branch. This is in contrast to the inert branch, where $r \equiv 3 \equiv W_p \pmod 8$ permits $r = W_p$ and hence \emph{phantom} danger triples (Definition~\ref{def:danger-triple}, Remark~\ref{rem:known-danger}).
\end{remark}

\subsection{Quartic residuosity of 2}

\begin{theorem}[Quartic residuosity at primitive Pell divisors]\label{thm:quartic}
Let $n \geq 1$ be odd. If $r \equiv 1 \pmod 8$ is a prime primitive divisor of $U_{4n}$, then $2^{(r-1)/4} \equiv 1 \pmod r$, i.e., 2 is a fourth-power residue modulo $r$.
\end{theorem}

\begin{proof}
In $\Q(\zeta_8) = \Q(i, \sqrt 2)$ we have the identity
\begin{equation*}
  \sqrt 2 \;=\; (1 - \zeta_8)^2 \cdot \zeta_8^3 \cdot (1 + \sqrt 2), \tag{$\ast$}
\end{equation*}
verified as follows: $(1 - \zeta_8)^2 = 1 - 2 \zeta_8 + \zeta_8^2 = 1 - 2 \zeta_8 + i$; multiplying by $\zeta_8^3$ and using $\zeta_8^2 = i$, a direct multiply-out gives $(1 - \zeta_8)^2 \cdot \zeta_8^3 = 2 - \sqrt 2$, then $(2 - \sqrt 2)(1 + \sqrt 2) = 2 + 2\sqrt 2 - \sqrt 2 - 2 = \sqrt 2$. ($\ast$) is confirmed.

Since $r \equiv 1 \pmod 8$, $r$ splits completely in $\Z[\zeta_8]$; fix a prime ideal $\mathfrak{P} \mid r$. Reduction mod $\mathfrak{P}$ embeds $\zeta_8 \mapsto \omega$ (a primitive 8-th root of unity in $\F_r^\times$), $\sqrt 2 \mapsto s$, and $1 + \sqrt 2 \mapsto \alpha$, giving
\begin{equation*}
  s \;\equiv\; (1 - \omega)^2 \cdot \omega^3 \cdot \alpha \qquad\text{in } \F_r. \tag{$\ast\ast$}
\end{equation*}

Raise $(\ast\ast)$ to $(r-1)/2$. We first record, \emph{unconditionally}, that a prime $r \equiv 1 \pmod 8$ with $\rho(r) = 4 n$ (i.e. a primitive divisor of $U_{4n}$) and $n$ odd has $\ord(\alpha) = 8 n$ and $\alpha^{4n} = -1$. Indeed, in the split algebra $\alpha^m$ has the two components $V_m/2 \pm U_m s$, so $\alpha^m \in \{\pm 1\}$ --- an identity of \emph{both} components, whose difference gives $2 U_m s = 0$ --- forces $r \mid U_m$, i.e. $\rho(r) \mid m$ (for odd $m$ the norm $\mathrm N(\alpha)^m = -1 \neq \mathrm N(\pm 1)$ already rules $\alpha^m \in \{\pm 1\}$ out); the converse holds when $m$ is even, the only case used below. At $m = 4 n$ the identity $V_{4n}^2 - 8 U_{4n}^2 = 4$ with $r \mid U_{4n}$ gives $V_{4n}/2 \equiv \pm 1 \pmod r$, so $\alpha^{4n} \in \{\pm 1\}$. If $\alpha^{4n} = +1$, then $e := \ord(\alpha)$ divides $4 n$; but $\alpha^e = 1$ forces $r \mid U_e$, hence $\rho(r) = 4 n \mid e$, so $e = 4 n$, whence $\alpha^{2n}$ has order $2$, $\alpha^{2n} = -1$, and $r \mid U_{2n}$ --- contradicting $\rho(r) = 4 n \nmid 2 n$. Therefore $\alpha^{4n} = -1$ and $\ord(\alpha) = 8 n$. (This is the converse direction of the rank--order correspondence; Proposition~\ref{prop:rank-equiv}(a), which assumes Condition~(II) at $r$, is not invoked here, and Theorem~\ref{thm:quartic} makes no Condition-(II) hypothesis.) Writing $r - 1 = 8 n M$ for a positive integer $M$:

\begin{itemize}
\item $[(1 - \omega)^2]^{(r-1)/2} \equiv 1 \pmod{\mathfrak{P}}$: $(1 - \omega)^{r-1} = 1$ by Fermat's little theorem in $\F_r$.
\item $[\omega^3]^{(r-1)/2} = \omega^{3(r-1)/2}$: since $\ord(\omega) = 8$ and $(r-1)/2 = 4 n M$, $\omega^{3 \cdot 4 n M} = \omega^{12 n M}$. Reducing mod 8: $12 n M \equiv 4 n M \pmod 8$. Since $n$ is odd, $4 n M \equiv 4 M \pmod 8$, so $\omega^{4 n M} = \omega^{4 M} = (\omega^4)^M = (-1)^M$.
\item $\alpha^{(r-1)/2} = \alpha^{4 n M} = (\alpha^{4n})^M = (-1)^M$.
\end{itemize}

Multiplying the three factors gives $s^{(r-1)/2} = 1 \cdot (-1)^M \cdot (-1)^M = (-1)^{2 M} = 1$.

Hence $2^{(r-1)/2} = s^{r-1} = (s^{(r-1)/2})^2 = 1$ trivially, but the sharper conclusion $s^{(r-1)/2} = 1$ means $s \in (\F_r^\times)^2$, so $\sqrt s \in \F_r$. Now $s^2 = 2$, so $(\sqrt s)^4 = 2$, and $\sqrt s$ has order dividing $r - 1$: writing $\sqrt s = t$ with $t^4 = 2$: $2^{(r-1)/4} = t^{r-1} = 1$.
\end{proof}

\begin{remarknonum}
Theorem~\ref{thm:quartic} is recorded for context. It is cited once, in the obstruction analysis of \S\ref{sec:nct-problem}, but is not used in the reduction chain itself: the split branch is closed through the Pell-rank reformulation (Proposition~\ref{prop:rank-equiv}) and the nonexistence of compatible triples (Theorem~\ref{thm:nct99} and Conjecture~\ref{conj:nct-general}). Further split-local character refinements are relegated to the supplementary material.
\end{remarknonum}

\subsection{No Compatible Triples}\label{sec:nct}

\begin{theorem}[NCT for $d \leq 99$]\label{thm:nct99}
For each odd $d$ with $1 \leq d \leq 99$, no compatible triple $(d, r, p)$ exists.
\end{theorem}

\begin{proof}
Call $d$ \emph{admissible} if every prime factor $q$ of $d$ satisfies $v_2(\ord_q(2)) = 1$; otherwise $q \mid W_{p-2}$ is impossible (if $\ord_q(2)$ is odd, no half-order $m_q$ exists; if $4 \mid \ord_q(2)$, then $m_q$ is even but $p - 2$ is odd, contradicting $m_q \mid p - 2$). The parity obstruction eliminates ${\sim}80\%$ of odd $d \leq 500$; within $[1, 99]$ it leaves 14 values --- the 12 in $[1, 81]$ together with $\{83, 99\}$. The seven excluded values $\{85, 87, 89, 91, 93, 95, 97\}$ in $(81, 99]$ are parity-blocked: each has a prime factor $q$ with $q \notin A$ (e.g., $85 = 5 \cdot 17$ with $5 \notin A$; $91 = 7 \cdot 13$ with $7 \notin A$; etc.).

\emph{Cases $d = 1, 3$.} $U_4 = 12$ and $U_{12} = 13{,}860$ have no prime factor $\equiv 1 \pmod 8$, so no split $r$ exists.

\emph{Cases $d \in \{9, 11, 19, 27, 33, 43, 57, 59, 67, 81\}$.} $U_{4 d}$ is completely factored (largest at $d = 81$, $124$ digits); each of the 18 primitive divisors $r \equiv 1 \pmod 8$ with $\rho(r) = 4 d$ is checked and found to have $\ord_r(2)$ either odd or $\ord_r(2)/2$ composite, so no Wagstaff prime $p$ can correspond.

\emph{Cases $d \in \{83, 99\}$.} Rather than factor the full $U_{4 d}$ (sizes $\approx \alpha^{4 d}$, prohibitive at $d = 99$ for general-purpose ECM), we use the cyclotomic decomposition
\[
  V_{2 d} \;=\; \prod_{\substack{m \mid d \\ m\text{ odd}}} \Psi_{4 m}, \qquad U_n \;=\; \prod_{e \mid n} \Psi_e \quad\text{(M\"obius-inverted),}
\]
where $\Psi_n$ is the $n$-th \emph{primitive part} of the Pell $U$-sequence. The decomposition is derived from $U_{2 n} = V_n U_n$ (specialise $n = 2 d$, so $V_{2 d} = U_{4 d}/U_{2 d}$; the surviving primitive parts are exactly $\Psi_e$ with $e \mid 4 d$, $e \nmid 2 d$, i.e., $e = 4 m$ with $m \mid d$ odd). Each piece $\Psi_{4 m}$ has size $\approx \alpha^{2 \varphi(m)}$, far smaller than $V_{2 d}$ itself. Every primitive split divisor of $V_{2 d}$ --- and hence every potential primitive Wagstaff factor at this $d$ --- lies in some $\Psi_{4 m}$.

For $d = 83$ (prime, $\ord_{83}(2) = 82 = 2 \cdot 41$, so $83 \in A$): $\Psi_{4 \cdot 83}$ has 63 decimal digits and five prime factors; four are $\equiv 3 \pmod 8$ (inert, not relevant to the split-branch search), and the single split factor
\[
  r = 7{,}370{,}827{,}889{,}154{,}320{,}981{,}521 \quad (\equiv 1 \pmod 8)
\]
satisfies
\[
  \begin{aligned}
  \ord_r(2) &= 1{,}842{,}706{,}972{,}288{,}580{,}245{,}380,\\
  \ord_r(2)/2 &= 921{,}353{,}486{,}144{,}290{,}122{,}690.
  \end{aligned}
\]
The half-order is even, hence composite. Not Wagstaff.

For $d = 99 = 9 \cdot 11$ (with $\ord_{11}(2) = 10 = 2 \cdot 5$, $11 \in A$, and $3 \in A$): $\Psi_{4 \cdot 99}$ has 46 decimal digits, a balanced semiprime of two split factors. One factor satisfies $\ord_r(2)$ odd; the other has $\ord_r(2)/2$ divisible by $5^2$, hence composite. Neither is Wagstaff.

No primitive split divisor at $d \in \{83, 99\}$ yields a compatible triple. Combined with the cases above, this closes NCT for $d \leq 99$ unconditionally.
\end{proof}

(Detailed per-$d$ obstruction table, factorizations of $U_{4 d}$ for $d \leq 81$, and factorizations of $\Psi_{4 m}$ for odd $m \leq 99$, are archived in the supplementary Zenodo deposit; each asserted prime factor used in the finite exhaustion is independently certified by APR-CL in PARI/GP, with BPSW and deterministic Miller--Rabin used only as preliminary screens.)

\begin{proposition}[$q$-th power necessity]\label{prop:qth-power}
In any compatible triple $(d, r, p)$, for every odd prime $q \mid d$: 2 must be a $q$-th power residue modulo $r$.
\end{proposition}

\begin{proof}
Suppose not. Since $r$ is a primitive divisor of $U_{4d}$ with $r \equiv 1 \pmod 8$, its rank of apparition $\rho(r) = 4 d$ divides $r - 1$; hence $q \mid d \mid r - 1$, and the failure of 2 to be a $q$-th-power residue modulo $r$ forces $q \mid \ord_r(2) = 2 p$. So $q = p$, whence $p \mid d \mid W_{p-2}$, contradicting Lemma~\ref{lem:pndivW}.
\end{proof}

\begin{conjecture}[NCT for general $d$]\label{conj:nct-general}
For every odd $d \geq 1$, no compatible triple $(d, r, p)$ exists.
\end{conjecture}

\begin{remark}\label{rem:strengthening-false}
A natural strengthening of Proposition~\ref{prop:qth-power} --- that $v_2(\ord_r(2)) \geq 2$ whenever 2 is a $q$-th power modulo a primitive divisor $r$ of $U_{4 q}$ --- is \emph{false}. For $q = 163$, the primitive divisor $r = 12{,}036{,}824{,}328{,}615{,}957{,}881$ of $U_{652}$ satisfies $2^{(r-1)/163} \equiv 1 \pmod r$ yet $v_2(\ord_r(2)) = 1$. Conjecture~\ref{conj:nct-general} is unaffected: $\ord_r(2)/2$ is composite (\S\ref{sec:primdiv-survey}). The take-away: the split-branch gap is structural, and must be attacked directly via compatible triples.
\end{remark}

\subsection{The Mixed Pair Conjecture}\label{sec:mpc}

The reduction invokes NCT only inside a global Condition-(II) pass, and a global pass always supplies a second, inert, locally passing factor (Lemma~\ref{lem:mod8}: the inert multiplicity is odd, hence at least one inert factor exists). The split-branch residual can therefore be weakened from the single-coincidence statement NCT to the following pair-type statement.

\begin{conjecture}[Mixed Pair Conjecture, MPC]\label{conj:mpc}
No prime $p \geq 5$ admits simultaneously
\begin{itemize}
\item[(i)] a compatible triple $(d_s, r_s, p)$: an odd $d_s \geq 1$, a prime $r_s \equiv 1 \pmod 8$ that is a primitive divisor of $U_{4 d_s}$, with $\ord_{r_s}(2) = 2p$ and $d_s \mid W_{p-2}$ (Definition~\ref{def:compat-triple});
\item[(ii)] an inert prime $r_{\mathrm{in}} \equiv 3 \pmod 8$ with $\ord_{r_{\mathrm{in}}}(2) = 2p$ and $d_{r_{\mathrm{in}}} \mid W_{p-2}$, where $d_{r_{\mathrm{in}}} = \ord_{r_{\mathrm{in}}}(\omega_3)/4$.
\end{itemize}
\end{conjecture}

\begin{remark}[Reading the two data]\label{rem:mpc-data}
The divisibilities $r_s \mid W_p$ and $r_{\mathrm{in}} \mid W_p$ are automatic: $\ord_r(2) = 2p$ forces $r \mid 2^p + 1$ with $r \neq 3$, hence $r \mid W_p$. The two data are not symmetric, and datum (ii) depends on the conjunction. \emph{Given} (i), $W_p$ is composite --- $r_s \equiv 1 \pmod 8$ is a proper divisor of $W_p \equiv 3 \pmod 8$ (Remark~\ref{rem:no-phantom-split}) --- and \emph{then} the inert datum (ii) automatically satisfies $d_{r_{\mathrm{in}}} > 43$ and admissibility (every prime factor of $d_{r_{\mathrm{in}}}$ in $A$), by the Pell-sequence exclusion (Theorem~\ref{thm:pell-exclusion}) and the parity obstruction (Theorem~\ref{thm:parity-obstruction}); so (ii) is then exactly an inert local Condition-(II) pass at a composite $W_p$. Standalone, datum (ii) has \emph{phantom} instances at prime $W_p$: for example $(p, r, d) = (5, 11, 3)$ satisfies (ii) with $W_5 = 11$ prime. The conjunction with (i) excludes them.
\end{remark}

\begin{proposition}[NCT-general implies MPC]\label{prop:nct-implies-mpc}
Conjecture~\ref{conj:nct-general} implies Conjecture~\ref{conj:mpc}.
\end{proposition}

\begin{proof}
A violation of MPC contains, in datum (i), a compatible triple.
\end{proof}

No converse implication is known, so MPC is \emph{formally weaker as a hypothesis} than NCT-general; this is the only comparison we assert. (Establishing that MPC is properly weaker would require exhibiting a compatible triple without an inert co-passer --- a configuration that would itself falsify NCT-general; and if NCT-general is true, the material implication MPC $\Rightarrow$ NCT-general holds trivially.) MPC forbids a coincidence of two pinned local passes at one $p$ --- the same shape as PPPC --- so the two principal residuals of the working chain (\S\ref{sec:reduction}) are structurally uniform: both are pair-type statements. MPC inherits every unconditionally proved NCT range: a violation requires $d_s > 400$ (Theorem~\ref{thm:nct99} and the complete fixed-$d$ closures, \S\ref{sec:nct-verify}); on the inert side, the order-alignment in datum (ii) is independently rare (\S\ref{sec:empirical}: $0$ of $369{,}876$ tested candidates satisfy it). A joint heuristic for MPC is recorded in \S\ref{sec:heuristic}. Proving MPC meets the same structural obstruction as NCT --- Observation~\ref{obs:additive-multiplicative} applies unchanged to its split datum --- so MPC is offered as a weaker residual, not as a more tractable theorem target.

For the cutoff-split form of the reduction (Theorem~\ref{thm:sharpened-chain}) write \textbf{MPC$^{>}$} for the restriction of MPC in which datum (ii) is additionally required to satisfy $r_{\mathrm{in}} > 10^{12}$. The complementary leg $r_{\mathrm{in}} \leq 10^{12}$ is absorbed by the Platinum enumeration restricted to composite $W_p$ and by the witness-exponent statement X1 (Conjecture~\ref{conj:x1}); see \S\ref{sec:reduction}.

\section{The inert branch: \texorpdfstring{$r \equiv 3 \pmod 8$}{r = 3 mod 8}}\label{sec:inert}

\subsection{The gap structure}\label{sec:gap-structure}

For $r \equiv 3 \pmod 8$ with $r \mid W_p$ composite: $(2/r) = -1$, so $\Z[\sqrt 2]/(r) \cong \F_{r^2}$. The element $\omega_3$ has norm $+1$, so lies in the norm-1 subgroup $\mu_{r+1}$; by Lemma~\ref{lem:v2-class}, $v_2(\ord(\omega_3)) = 2$, so $\ord(\omega_3) = 4 g$ for $g$ odd. Combined with $\ord(\omega_3) \mid r + 1$ and $v_2(r + 1) = 2$: $g \mid (r+1)/4$.

\textbf{Descent.} Condition~(II) at $r$ becomes $\omega_3^{(N+1)/2} = -1$, i.e., $\omega_3^{2 W_{p-2}} = -1$. With $\ord(\omega_3) = 4 g$: this holds iff $g \mid W_{p-2}$.

\textbf{Notation.} For inert $r \mid W_p$:
\[
  G_r \;=\; \gcd\!\left(\frac{r+1}{4},\; W_{p-2}\right).
\]
If Condition~(II) passes, then $g \mid G_r$.

\begin{remark}[Why the norm argument doesn't suffice]\label{rem:norm-insufficient}
Condition~(II) already implies $\omega_3^N \equiv \omega_3^{-1} \pmod N$, so quadratic Frobenius strengthenings \cite{ref:Grantham} add nothing. The norm $\mathrm{N}(\alpha) = -1$ constrains only the 2-part of orders; the odd part $g$ is not constrained by the norm argument alone.
\end{remark}

\subsection{Parity and valuation obstructions}

\begin{definition}[Class III primes]\label{def:classIII}
An odd prime $q$ is Class III if $v_2(\ord_q(2)) = 1$. Write $A = \{q : v_2(\ord_q(2)) = 1\}$.
\end{definition}

Class-III primes are exactly the odd primes that can divide $W_n = (2^n + 1)/3$ for odd $n \geq 3$: $q \mid W_n$ requires $\ord_q(2) = 2 m$ with $m$ odd. Every $q \equiv 3 \pmod 8$ lies in $A$; $q \equiv 5$ or $7 \pmod 8$ do not.

\begin{theorem}[Parity obstruction for inert factors]\label{thm:parity-obstruction}
Let $r \equiv 3 \pmod 8$, $r \mid W_p$ composite, $g = \ord_r(\omega_3)/4$. If $g$ has a prime factor $q \notin A$, then $g \nmid W_{p-2}$ and Condition~(II) fails at $r$.
\end{theorem}

\begin{proof}
$q \notin A$ and $q \mid g$: then $q \nmid W_n$ for any odd $n$, so $q \nmid W_{p-2}$.
\end{proof}

\begin{theorem}[Generalized LTE for Class III primes]\label{thm:lte}
Let $q \in A$ with $\ord_q(2) = 2 m$, $m$ odd. For odd $n$ with $m \mid n$: $v_q(2^n + 1) = v_q(2^m + 1) + v_q(n / m)$. For $q \geq 5$: $v_q(W_n) = v_q(2^m + 1) + v_q(n/m)$; for $q = 3$: $v_3(W_n) = v_3(n)$.
\end{theorem}

\begin{proof}
Apply the lifting-the-exponent lemma to $(2^m)^\ell + 1$ with $q \mid 2^m + 1$.
\end{proof}

Computationally, $v_q(2^m + 1) = 1$ for all 194 Class III primes $q < 5000$, where $m = \ord_q(2)/2$; no ``Wieferich-at-minus-1'' behaviour is observed.

\textbf{Valuation obstruction.} If $v_q(g) > v_q(W_{p-2})$ for some $q \in A$ with $q \mid g$, then $g \nmid W_{p-2}$ and Condition~(II) fails.

Among the 869 inert factors with $r < 2 \times 10^5$ tested in \S\ref{sec:inert-survey}, the parity and valuation obstructions account for the bulk of the exclusions; together with the Pell-sequence exclusion (Theorem~\ref{thm:pell-exclusion}) they block 845 of the 869 (Corollary~\ref{cor:small-Gr}).

\subsection{Pell-sequence exclusion}

\begin{theorem}[Pell-sequence divisibility]\label{thm:pell-divisibility}
Let $r \equiv 3 \pmod 8$, $r \mid W_p$, and $\ord_r(\omega_3) = 4 d$. Then
\[
  r \;\;\Big|\;\; \gcd\!\left(\frac{V_{4d} + 2}{2},\; U_{4d}\right).
\]
\end{theorem}

\begin{proof}
$\ord_r(\omega_3) = 4 d$ gives $\omega_3^{2d} = -1$, hence $\alpha^{4d} = -1$. Expanding $\alpha^{4d} = V_{4d}/2 + U_{4d} \sqrt 2$ and comparing with $-1 = -1 + 0 \cdot \sqrt 2$: $r \mid U_{4d}$ and $r \mid (V_{4d}/2 + 1) = (V_{4d} + 2)/2$.
\end{proof}

\begin{definition}[Pell exclusion number]\label{def:pell-exclusion-number}
For odd $d \geq 1$: $N_d = \gcd((V_{4d} + 2)/2,\; U_{4d})$.
\end{definition}

\begin{proposition}[Closed form]\label{prop:pell-closed-form}
For odd $d \geq 1$: $N_d = V_{2d} = V_d^2 + 2$.
\end{proposition}

\begin{proof}
Three Pell identities: $V_{4d} = V_{2d}^2 - 2$, so $V_{4d} + 2 = V_{2d}^2$; $U_{4d} = U_{2d} V_{2d}$; $(V_{2d}/2)^2 - 2 U_{2d}^2 = 1$, so $\gcd(V_{2d}/2, U_{2d}) = 1$. Since $V_{2d} = 2k$ even, $N_d = \gcd(2k^2, 2 k U_{2d}) = 2k = V_{2d}$. Finally $V_{2d} = V_d^2 - 2(-1)^d = V_d^2 + 2$ for odd $d$.
\end{proof}

\begin{lemma}[Inert primitivity]\label{lem:inert-primitivity}
Let $r \equiv 3 \pmod 8$ be prime and suppose $\ord_r(\omega_3)=4d$. Then $r\mid V_{2d}$ and $r\nmid V_{2e}$ for every integer $0<e<d$. Thus $r$ is primitive for the $d$-indexed Pell--Lucas family $\{V_{2e}\}_{e\geq1}$ at index $d$.
\end{lemma}

\begin{proof}
Write $V_{2e}=\omega_3^e+\omega_3^{-e}$. Since $\ord_r(\omega_3)=4d$, the element $\omega_3^d$ has order $4$, so $(\omega_3^d)^2=-1$ and $\omega_3^d+\omega_3^{-d}=0$; hence $r\mid V_{2d}$.

Conversely, if $r\mid V_{2e}$, then $\omega_3^e+\omega_3^{-e}=0$ in $\F_{r^2}$, so $\omega_3^{2e}=-1$. Also $\omega_3^{2d}=-1$. Therefore $\omega_3^{2(e-d)}=1$, and $4d\mid 2(e-d)$, equivalently $e\equiv d\pmod{2d}$. No integer $e$ with $0<e<d$ satisfies this congruence.
\end{proof}

\begin{remark}\label{rem:zsigmondy}
Since $N_d = V_d^2 + 2 = \mathrm{N}_{\Z[\sqrt{-2}]}(V_d + \sqrt{-2})$, every prime $r \equiv 3 \pmod 8$ dividing $N_d$ splits in $\Z[\sqrt{-2}]$ (since $(-2/r) = 1$). This Zsigmondy-type realization is the structural basis of attack angle N2 (Stewart--Tijdeman bounds on $\omega(N_d)$); see \S\ref{sec:pppc-directions}.
\end{remark}

\begin{theorem}[Pell-sequence exclusion for $d \leq 43$]\label{thm:pell-exclusion}
For each odd $d$ in the admissible set $\{1, 3, 9, 11, 19, 27, 33, 43\}$ (the $d \leq 43$ with all prime factors in $A$): no prime $r \equiv 3 \pmod 8$ that divides a composite $W_p$ and passes Condition~(II) can have $\ord_r(\omega_3) = 4 d$.
\end{theorem}

\begin{proof}
Suppose, for contradiction, that $r \equiv 3 \pmod 8$ divides a composite $W_p$, passes Condition~(II), and has $\ord_r(\omega_3) = 4 d$ for some $d$ in the admissible set. By the descent of \S\ref{sec:gap-structure}, Condition~(II) at $r$ is equivalent to $d \mid W_{p-2}$; and by Theorem~\ref{thm:pell-divisibility}, $r \mid N_d = V_d^2 + 2$ (Proposition~\ref{prop:pell-closed-form}). We exhaust the admissible $d$:

\emph{$d = 1$:} $N_1 = V_1^2 + 2 = 2^2 + 2 = 6 = 2 \cdot 3$. The only odd prime factor is 3, and $3 \mid W_p$ is impossible by Lemma~\ref{lem:ord2}. No valid $r$.

\emph{$d = 3$:} $N_3 = V_3^2 + 2 = 14^2 + 2 = 198 = 2 \cdot 3^2 \cdot 11$. The odd prime factors $\equiv 3 \pmod 8$ are 3 (excluded) and 11. For $r = 11$: $\ord_{11}(2) = 10 = 2 \cdot 5$, so the unique prime $p$ with $\ord_r(2) = 2 p$ is $p = 5$. But $W_5 = 11$ is itself prime, so $r = 11$ does not divide any \emph{composite} $W_p$. No valid $r$.

\emph{$d \in \{9, 11, 19, 27, 33, 43\}$:} We use the complete factorizations of $N_d = V_d^2 + 2$ (computed with SymPy \texttt{factorint}, results in the supplementary archive; the largest factorization is $N_{43}$, a 33-digit integer whose prime decomposition has six prime factors, all under 20 digits). For each prime factor $r$ of $N_d$ with $r \equiv 3 \pmod 8$ and $r > 11$, we compute $\ord_r(2)$ and check:
\begin{enumerate}
\item If $v_2(\ord_r(2)) \neq 1$, then $r \nmid W_p$ for any prime $p$ (by Lemma~\ref{lem:mod8}); $r$ is excluded.
\item If $\ord_r(2) = 2 m$ with $m$ composite or $m$ not prime, then no prime $p$ satisfies $\ord_r(2) = 2 p$; $r$ is excluded.
\item Otherwise, $\ord_r(2) = 2 p$ for a unique prime $p$, and Condition~(II) at $r$ requires $d \mid W_{p-2}$. One survivor must be set aside first: $r = 43$, with $\ord_{43}(2) = 14$ and hence $p = 7$, divides only $W_7 = 43$, which is itself prime --- so $r = 43$ divides no \emph{composite} $W_p$ (the exclusion already used for $r = 11$ at $d = 3$). For every other case-3 survivor the divisibility $d \mid W_{p-2}$ fails --- either the arithmetic progression of Theorem~\ref{thm:exact-ap} ($p \equiv a_d \pmod{2 m_d}$) is incompatible with $p = \ord_r(2)/2$, or $v_3(d) > v_3(p-2)$ (Theorem~\ref{thm:lte}, valuation excess for the prime 3), or an analogous $q$-adic mismatch occurs for some $q \mid d$ with $q > 3$ --- so $d \nmid W_{p-2}$, contradicting Condition~(II).
\end{enumerate}
Condition 1 or 2 alone disposes of the overwhelming majority of the primitive prime factors of $N_d$ for $d \in \{9, 11, 19, 27, 33, 43\}$; condition 3 is needed only for a handful of survivors. The complete per-$r$ verification table is archived with the paper (supplementary data). No admissible $d \leq 43$ yields a valid $r$.

For $d$ not in the admissible set (i.e. $d \leq 43$ with some prime factor $q \notin A$): Theorem~\ref{thm:parity-obstruction} gives $g \nmid W_{p-2}$, so $\ord_r(\omega_3)/4 \neq d$ in any Condition-(II)-passing scenario, regardless of the value of $N_d$.
\end{proof}

\begin{corollary}[Inert factors with small $G_r$]\label{cor:small-Gr}
If $r \equiv 3 \pmod 8$ divides a composite $W_p$ and $G_r \leq 43$, then Condition~(II) fails at $r$.
\end{corollary}

\begin{proof}
Condition~(II) gives $g \leq G_r \leq 43$, $g$ odd. The prime factors of $g$ are all in $A$ (else Theorem~\ref{thm:parity-obstruction} applies); so $g$ is in the admissible set of Theorem~\ref{thm:pell-exclusion}, which excludes it.
\end{proof}

\subsection{Lower bound on \texorpdfstring{$g$}{g}}

\begin{lemma}[Class-III division]\label{lem:classIII-division}
For $q \in A$ with $q > 3$ and $m_q = \ord_q(2)/2$: $q \mid W_{p-2}$ iff $m_q \mid p - 2$.
\end{lemma}

\begin{proof}
Let $n=p-2$, so $n$ is odd. If $q \mid W_n$, then $2^n \equiv -1 \pmod q$; since $2^{m_q}\equiv -1 \pmod q$ and $\ord_q(2)=2m_q$, this is equivalent to $n \equiv m_q \pmod {2m_q}$, in particular $m_q\mid n$. Conversely, if $m_q\mid n$, then $n/m_q$ is odd (both $m_q$ and $n$ are odd), so $2^n=(2^{m_q})^{n/m_q}\equiv -1 \pmod q$, whence $q\mid 3W_n$ and, as $q>3$, $q\mid W_n$.
\end{proof}

\begin{lemma}[Lower bound]\label{lem:lower-bound}
For inert $r \mid W_p$ composite: $g > \log(2p)/\log(3 + 2\sqrt 2) \approx 0.393 \cdot \log_2 p$.
\end{lemma}

\begin{proof}
From $\ord_r(\omega_3) = 4g$: $\omega_3^{2g} = -1 \pmod r$, hence $V_{2g} \equiv 0 \pmod r$. Since $r \geq 2p + 1$: $2p + 1 \leq V_{2g} < (3 + 2\sqrt 2)^g + 1$.
\end{proof}

\textbf{Product bound.} For a given $p$:
\[
  G_{\max}(p) \;=\; 3^{v_3(p - 2)} \cdot
  \prod_{\substack{q > 3,\; q \in A,\\ m_q \mid (p-2)}}
  q^{v_q(W_{p-2})}
\]
is an upper bound on any compatible $g$. The product ranges over all Class III primes, not only over primes $q \equiv 3 \pmod 8$: Class III primes with $q \equiv 1 \pmod 8$ can divide $W_{p-2}$ and can occur in $(r+1)/4$ for inert $r$. When $G_{\max}(p) < 3$, no inert factor can pass Condition~(II).

\subsection{Density limits}

Corollary~\ref{cor:small-Gr} handles 845 of 869 known inert factors unconditionally (\S\ref{sec:inert-survey}). A natural subclass of the residual configurations has density zero:

\begin{proposition}[Density zero for pure Class III inert factors]\label{prop:density-zero}
The set of primes $r \equiv 3 \pmod 8$ for which every odd prime factor of $(r + 1)/4$ belongs to Class III has density zero within the primes $r \equiv 3 \pmod 8$.
\end{proposition}

\begin{proof}[Proof sketch]
Let $S = \{q : v_2(\ord_q(2)) \geq 2\}$. $S$ contains all $q \equiv 5 \pmod 8$, so $\sum_{q \in S} 1/q = \infty$. The hypothesis forces $q \nmid (r + 1)$ for every $q \in S$, equivalently $r \not\equiv -1 \pmod q$. For each finite $S' \subset S$, the primes satisfying these conditions have density at most $\prod_{q \in S'} (q-2)/(q-1)$ in the set of primes $\equiv 3 \pmod 8$; since $\sum 1/(q-1) = \infty$ this product has infimum zero.
\end{proof}

The full residual set is slightly larger than this subclass (it requires only $g$'s prime factors to be Class III, not all of $(r+1)/4$'s), and Proposition~\ref{prop:density-zero} is a density-zero (not pointwise) statement; a naive GRH/Chebotarev argument in the style of Hooley \cite{ref:Hooley} and Heath-Brown \cite{ref:HeathBrown} fails to produce a pointwise closure. See \S\ref{sec:pppc-problem} for the structural barrier.

\subsection{Multi-factor reduction}\label{sec:multi-factor}

\begin{definition}[Danger triple]\label{def:danger-triple}
A danger triple $(d, r, p)$ consists of: an odd integer $d > 43$ with every prime factor in $A$; a prime $r \equiv 3 \pmod 8$ that is primitive for the $d$-indexed Pell--Lucas family at $V_{2 d}$ and has $\ord_r(\omega_3) = 4 d$ in $\F_{r^2}$; and the unique prime $p$ with $\ord_r(2) = 2 p$, satisfying $d \mid W_{p-2}$. A danger triple is phantom if $W_p$ is prime; otherwise real.
\end{definition}

By Theorem~\ref{thm:pell-divisibility} and Lemma~\ref{lem:inert-primitivity}, any inert factor $r$ of a composite $W_p$ with $d > 43$ gives rise to a danger triple.

\begin{theorem}[Multi-factor reduction]\label{thm:multi-factor}
Let $W_p$ be composite with every prime factor $r \equiv 3 \pmod 8$. If Condition~(II) holds at every factor, then $W_p$ has $t \geq 3$ prime factors counted with multiplicity, and each factor $r_i$ has $d_i = \ord_{r_i}(\omega_3)/4 > 43$ with every prime factor of $d_i$ in $A$.
\end{theorem}

\begin{proof}
\emph{Three-or-more factors.} Let $t$ be the number of prime factors of $W_p$ counted with multiplicity. Each prime factor $r_j \equiv 3 \pmod 8$, so the product $\prod r_j \equiv 3^t \pmod 8$, the product running over the $t$ prime factors with multiplicity. Combined with $W_p \equiv 3 \pmod 8$ (Lemma~\ref{lem:mod8}): $3^t \equiv 3 \pmod 8$, forcing $t$ odd. Since $W_p$ is composite, $t \geq 3$.

\emph{Per-factor $d_i > 43$.} Each factor $r_i$ passes Condition~(II), so $d_{r_i} \mid W_{p-2}$ (\S\ref{sec:gap-structure}); hence every prime factor of $d_{r_i}$ lies in $A$ by Theorem~\ref{thm:parity-obstruction} (the parity obstruction). Pell-sequence exclusion (Theorem~\ref{thm:pell-exclusion}) rules out $d_{r_i} \leq 43$ for an inert factor of a composite $W_p$. Therefore $d_{r_i} > 43$.
\end{proof}

\begin{remark}[Known danger triples, $d \leq 1000$]\label{rem:known-danger}
A primitive-divisor survey of the admissible $d \leq 1000$ has identified five danger triples:
\[
  \begin{aligned}
  (d, p) \in \{& (57, 11),\; (57, 85{,}684{,}865{,}933),\\
               & (67, 24{,}822{,}855{,}876{,}938{,}710{,}967),\; (331, 17),\; (683, 13)\}.
  \end{aligned}
\]
Three are \emph{phantom} (with $W_p$ prime): $(57, 11)$ since $W_{11} = 683$, $(331, 17)$ since $W_{17} = 43{,}691$, $(683, 13)$ since $W_{13} = 2731$; the two real triples are disposed of at secondary factors (\S\ref{sec:secondary}). The survey is complete --- full factorization of the Pell--Lucas number $V_{2d}$ --- only for the $d$ whose $V_{2d}$ was factored completely (notably $d \in \{57, 67\}$, via $V_{114}$ and $V_{134}$); for the remaining admissible $d$ the factorization is partial (e.g. $V_{1366}$ at $d = 683$ retains composite half-cofactors totalling more than five hundred digits), so the catalog was recorded as a lower bound; Proposition~\ref{prop:danger-census-400} below makes it exhaustive for every admissible $d \leq 400$, and it remains a lower bound only on $(400, 1000]$. The catalog itself is \emph{not} used in the reduction: the cross-case argument of Section~\ref{sec:cross-case} routes through the Platinum Lemma --- an $r$-indexed enumeration (Theorem~\ref{thm:platinum}) --- and not through any $d$-indexed catalog; the census enters only through Corollary~\ref{cor:inert-floor-400} and the diagonal closures.
\end{remark}

\begin{proposition}[Complete danger census through $d = 400$]\label{prop:danger-census-400}
For the admissible odd $d$ with $43 < d \leq 400$ --- the $34$ values
\[
  \begin{gathered}
  57,\, 59,\, 67,\, 81,\, 83,\, 99,\, 107,\, 121,\, 129,\, 131,\, 139,\, 163,\, 171,\, 177,\, 179,\, 201,\, 209,\\
  211,\, 227,\, 243,\, 249,\, 251,\, 281,\, 283,\, 297,\, 307,\, 321,\, 331,\, 347,\, 361,\, 363,\, 379,\, 387,\, 393
  \end{gathered}
\]
--- the danger triples $(d, r, p)$ of Definition~\ref{def:danger-triple} are exactly the four rows of Remark~\ref{rem:known-danger} with $d \leq 400$: the phantoms $(57, 683, 11)$ and $(331, 43{,}691, 17)$, and the real triples
\[
  (57,\ 171{,}369{,}731{,}867,\ 85{,}684{,}865{,}933), \qquad
  (67,\ 148{,}937{,}135{,}261{,}632{,}265{,}803,\ 24{,}822{,}855{,}876{,}938{,}710{,}967).
\]
In particular the $d \leq 400$ part of that catalog is exhaustive, and no two danger triples with the same $d \leq 400$ share the pinned exponent.
\end{proposition}

\begin{proof}[Verification]
A danger triple at $d$ has $r$ primitive at $V_{2 d}$ (Definition~\ref{def:danger-triple}, Lemma~\ref{lem:inert-primitivity}); non-admissible $d$ carry none, by the Class III requirement of the definition. For each of the $34$ admissible rungs the complete factorization of $V_{2 d}$ was assembled --- from the fixed-$d$ certificate bundle of \S\ref{sec:nct-verify}, whose pinned $\Psi_{4 d}$ factorizations cover the $28$ rungs in $(100, 400]$; from the diagonal certificates of Proposition~\ref{prop:diagonal-closures}; and from the public factor database \cite{ref:FactorDB} with local factoring for the remaining rungs --- and re-verified by the exact product identity, with multiplicities, against the recomputed $V_{2 d}$; the $226$ distinct prime factors occurring across the $34$ rungs are each certified prime by APR-CL. For each of the $340$ (rung, prime) pairs the danger prerequisites are then decided exactly as in Proposition~\ref{prop:diagonal-closures} --- residue $r \equiv 3 \pmod 8$; primitivity $\ord_r(\omega_3) = 4 d$, by order descent in $\F_{r^2}$; $v_2(\ord_r(2)) = 1$, by the square chain; primality of the pinned exponent $p = \ord_r(2)/2$ --- together with condition (v), $d \mid W_{p-2}$, by the exact criterion $2^{\,p-2} \equiv -1 \pmod{3 d}$. Exactly the four listed rows survive. Every verdict was reproduced by an independent re-implementation, and the survivor rungs' factorizations were additionally re-derived from scratch. Scripts (\S\ref{sec:toolchain}): \texttt{scripts/}\allowbreak\texttt{cp412\_danger\_}\allowbreak\texttt{census\_}\allowbreak\texttt{d400.py} and \texttt{scripts/}\allowbreak\texttt{cp412b\_t1\_}\allowbreak\texttt{aprcl\_}\allowbreak\texttt{sweep.py}; the per-rung certificate is archived with the supplementary data.
\end{proof}

\begin{corollary}[Inert order-parameter floor $d_r > 400$]\label{cor:inert-floor-400}
Let $W_p$ be composite and suppose Condition~(II) holds at every prime factor. Then every inert prime factor $r \equiv 3 \pmod 8$ of $W_p$ has $d_r = \ord_r(\omega_3)/4 > 400$. This subsumes the per-factor conclusion $d_r > 43$ of Theorems~\ref{thm:pell-exclusion} and~\ref{thm:multi-factor}.
\end{corollary}

\begin{proof}
Suppose $d_r \leq 400$ for some inert factor $r$. Condition~(II) at $r$ gives $d_r \mid W_{p-2}$ with $d_r$ odd (descent of \S\ref{sec:gap-structure}); every prime factor of $d_r$ is Class III (Theorem~\ref{thm:parity-obstruction}); $d_r > 43$ (Theorem~\ref{thm:pell-exclusion}, applicable as $W_p$ is composite); $r$ is primitive at $V_{2 d_r}$ (Lemma~\ref{lem:inert-primitivity}); and $\ord_r(2) = 2 p$ (Theorem~\ref{thm:order-pinning}). So $(d_r, r, p)$ is a real danger triple with $43 < d_r \leq 400$, and Proposition~\ref{prop:danger-census-400} forces $p = 85{,}684{,}865{,}933$ or $p = 24{,}822{,}855{,}876{,}938{,}710{,}967$. At either exponent a secondary prime factor of $W_p$ fails Condition~(II) (\S\ref{sec:secondary}): at the first, the inert $r_2 = 7{,}675{,}646{,}348{,}774{,}307{,}083$, by the parity obstruction (Proposition~\ref{prop:witness-discharges}); at the second, the split $r_2 = 5{,}474{,}333{,}343{,}676{,}555{,}561{,}818{,}313$, where $\ord_{r_2}(\omega_3) = 4 d_2$ with $d_2 = 3^3 \cdot 1021 \cdot p$ divisible by $p$ itself, so $d_2 \nmid W_{p-2}$ ($p \nmid W_{p-2}$: $2^{\,p-2} \equiv 2^{-1} \not\equiv -1 \pmod p$ for $p > 3$) and Condition~(II) fails at $r_2$. Either way Condition~(II) fails at a prime factor of $W_p$, contradicting the hypothesis.
\end{proof}

\begin{remark}[Two complementary slices]\label{rem:two-slices}
The census is the $d$-indexed complement of the $r$-indexed Platinum enumeration (Theorem~\ref{thm:platinum}): Platinum enumerates every locally passing inert factor with $r \leq 10^{12}$ and unrestricted $d_r$; the census covers every $d_r \leq 400$ with no bound on $r$. Between them, a locally passing inert factor of a composite $W_p$ has $d_r > 400$ or is one of the two census rows (both at exponents closed by secondary factors), and has $r > 10^{12}$ or is one of the three Platinum rows (two closed; the third is X1). In the mixed-pair notation of \S\ref{sec:mpc}, datum~(ii) of any MPC-violating pair therefore carries $d_{r_{\mathrm{in}}} > 400$; on the all-inert branch, every member of $T_p^{>}$ at an undischarged exponent carries $d_r > 400$. The reduction chain of Theorem~\ref{thm:sharpened-chain} is unchanged; the census enters as a floor, not as a new hypothesis.
\end{remark}

\section{Cross-case structure: pinning, Platinum, separation}\label{sec:cross-case}

This section records the cross-case structure that acts uniformly across all candidate composite $W_p$, beyond the branch-by-branch analysis of Sections~\ref{sec:split} and~\ref{sec:inert}. None of it depends on Conjecture R3 or on computational input at a single factor; each constraint is a structural consequence of $r \mid W_p$ with Condition~(II). The section also states the certified computational records --- the Platinum enumeration (Theorem~\ref{thm:platinum}) and the witness discharges (Proposition~\ref{prop:witness-discharges}) --- that the reduction theorems of \S\ref{sec:reduction} take as stated inputs.

\subsection{Order-Pinning Lemma}

\begin{theorem}[Order-Pinning]\label{thm:order-pinning}
Let $r$ be a prime factor of $W_p$ with $r > 3$. Then $\ord_r(2) = 2 p$.
\end{theorem}

\begin{proof}
$r \mid W_p$ gives $2^p \equiv -1 \pmod r$, so $\ord_r(2) \mid 2 p$ and $\ord_r(2) \nmid p$. Divisors of $2 p$ not dividing $p$: $\{2, 2 p\}$. If $\ord_r(2) = 2$ then $r \mid 3$, contradicting $r > 3$.
\end{proof}

Order-Pinning is the single most load-bearing structural identity in the paper. It converts the Wagstaff condition ``$r \mid W_p$'' into the sharp equation $\ord_r(2) = 2 p$, which means every prime factor of $W_p$ lies in the \emph{exactly one} arithmetic progression $r \equiv 1 \pmod{2 p}$ and shares a single prime $p$.

\subsection{Multi-Factor Pinning}

\begin{theorem}[Multi-Factor Pinning]\label{thm:multi-pinning}
Let $W_p$ be composite. All prime factors $r > 3$ of $W_p$ share the common value $\ord_r(2)/2 = p$. Consequently, two composite Wagstaff numbers $W_{p_1}, W_{p_2}$ with $p_1 \neq p_2$ share no prime factor $r > 3$.
\end{theorem}

\begin{proof}
Immediate from Theorem~\ref{thm:order-pinning}.
\end{proof}

Multi-Factor Pinning is a \emph{global} constraint relating all prime factors of one $W_p$ to a common ``pinned'' parameter $p$; the same principle blocks any prime from being a factor of two different $W_p$.

\subsection{The finite danger catalog}

\begin{definition}\label{def:danger-at-d}
A danger triple at $d$ is a triple $(r, p, d)$ with $r$ an inert prime factor of $W_p$, $d = \ord_r(\omega_3)/4$, and $(r, p, d)$ surviving every constraint of Section~\ref{sec:inert} (parity, Pell exclusion, Condition~(II) pass at $r$); cf.\ Definition~\ref{def:danger-triple}.
\end{definition}

The primitive-divisor survey of \S\ref{sec:multi-factor} identifies the known danger triples (Remark~\ref{rem:known-danger}).

\begin{remark}[Phantom exclusion of the known catalog]\label{rem:phantom-exclusion}
The five known danger triples (Remark~\ref{rem:known-danger}) have pairwise distinct pinned exponents, $p \in \{11,\; 85\,684\,865\,933,\; 24\,822\,855\,876\,938\,710\,967,\; 17,\; 13\}$. Since all prime factors of one $W_p$ share a single pinned exponent (Theorem~\ref{thm:multi-pinning}), no composite $W_p$ carries two factors drawn from this list.
\end{remark}

\begin{remark}[The catalog is not load-bearing]\label{rem:catalog-not-loadbearing}
Were the $d \leq 1000$ danger catalog known to be exhaustive, Remark~\ref{rem:phantom-exclusion} together with Theorem~\ref{thm:multi-factor} would force any composite all-inert $W_p$ passing Condition~(II) to have at least two factors with $d_i > 1000$. The catalog is \emph{not} known to be exhaustive (Remark~\ref{rem:known-danger}: $V_{2d}$ is only partially factored for most admissible $d$), so the reduction does not take this route. The reduction instead invokes the Platinum Lemma (Theorem~\ref{thm:platinum}) --- an $r$-indexed enumeration requiring no $d$-indexed completeness (\S\ref{sec:reduction}).
\end{remark}

\subsection{Platinum Lemma}

\begin{theorem}[Platinum Lemma]\label{thm:platinum}
For every prime $p$, at most one inert prime factor $r \leq 10^{12}$ of $W_p$ passes Condition~(II) at $r$.
\end{theorem}

\begin{proof}[Verification]
Every prime $r \equiv 3 \pmod 8$ with $r \leq 10^{12}$ is enumerated (\S\ref{sec:inert-survey}); by Order-Pinning each has a unique pinned exponent $p = \ord_r(2)/2 < 5 \times 10^{11}$. Pell-sequence exclusion (Theorem~\ref{thm:pell-exclusion}) disposes of $d_r \leq 43$; for $d_r > 43$, Condition~(II) at $r$ --- equivalently $d_r \mid W_{p-2}$ --- is tested by direct evaluation of $\omega_3^{2 W_{p-2}} \pmod r$. The enumeration exhibits exactly three inert factors $r \leq 10^{12}$ of \emph{composite} $W_p$ passing Condition~(II), and these lie at three \emph{distinct} exponents $p$ (the local-pass witnesses of \S\ref{sec:secondary}). Hence every prime $p$ has at most one passing inert factor $r \leq 10^{12}$. Full scripts and certificates are archived.
\end{proof}

The Platinum Lemma is the only computational input in the cross-case reductions; every downstream consequence including Case B of Theorem~\ref{thm:sharpened-chain} is deterministic once the enumeration is accepted.

\subsection{The witness exponents and X1}\label{sec:x1}

The Platinum enumeration records exactly three locally passing inert factors $r \leq 10^{12}$, at three distinct exponents (Theorem~\ref{thm:platinum}, \S\ref{sec:secondary}); we call the three exponents
\[
  p \;\in\; \{\,1{,}251{,}488{,}009,\;\; 10{,}916{,}765{,}939,\;\; 85{,}684{,}865{,}933\,\}
\]
the \emph{witness exponents}. (The enumeration ranges over prime factors of \emph{composite} $W_p$: at prime $W_p$ the pass of $r = W_p$ itself is automatic and structural --- $P_p = \{W_p\}$, Lemma~\ref{lem:pp-containment} --- and such phantom rows, e.g.\ $W_{11} = 683$, are not part of the record.) Two of the three exponents are discharged by explicit failing factors; the third is the content of a single-integer statement.

\begin{conjecture}[X1, at-a-factor form]\label{conj:x1}
Condition~(II) fails at some prime factor of $W_{10916765939}$.
\end{conjecture}

$W_{10916765939}$ \emph{is} composite: the Platinum witness $r = 152{,}834{,}723{,}147 = 14 p + 1$ divides it and passes Condition~(II) locally (\S\ref{sec:secondary}). A directed scan of the progression $r = 2 p k + 1$ at this exponent, complete for $k \leq 10^{12}$, finds no second prime factor: the witness factor itself, at $k = 7$, is the only hit. The at-a-factor form is deliberate: a single failing factor removes the exponent from every local-pass-defined set --- in particular it bounds $|P_{10916765939}|$ without resolving the full cofactor --- and a fortiori excludes a global pass. The global form (``$W_{p}$ fails Condition~(II)'') would \emph{not} by itself remove the exponent from every local-pass-defined set, since local passes at every factor do not imply a global pass when $W_p$ is not squarefree. The at-a-factor form is also exactly what the discharges below establish.

\begin{proposition}[Discharged witness exponents]\label{prop:witness-discharges}
Condition~(II) fails at an explicit prime factor of $W_p$ for two of the three witness exponents:

(a) $p = 85{,}684{,}865{,}933$: the prime
\[
  r_2 = 7{,}675{,}646{,}348{,}774{,}307{,}083 \equiv 3 \pmod 8
\]
divides $W_p$ with $\ord_{r_2}(2) = 2p$ and
\[
  d_{r_2} = \ord_{r_2}(\omega_3)/4 = 3 \cdot 487 \cdot 1{,}313{,}423{,}399{,}858{,}711;
\]
the prime factors $487$ and $1{,}313{,}423{,}399{,}858{,}711$ of $d_{r_2}$ are not Class III, so Condition~(II) fails at $r_2$ by the parity obstruction (Theorem~\ref{thm:parity-obstruction}). This is the secondary factor of the $d = 57$ witness already recorded in \S\ref{sec:secondary}.

(b) $p = 1{,}251{,}488{,}009$: the prime
\[
  r = 33{,}829{,}735{,}778{,}964{,}491 = 2 p k + 1, \quad k = 13{,}515{,}805, \quad r \equiv 3 \pmod 8,
\]
divides $W_p$ with $\ord_r(2) = 2p$ and
\[
  d_r = \ord_r(\omega_3)/4 = 3 \cdot 41 \cdot 68{,}759{,}625{,}567{,}001;
\]
the prime factors $41$ and $68{,}759{,}625{,}567{,}001$ of $d_r$ are not Class III, so Condition~(II) fails at $r$ by the parity obstruction.

Consequently neither exponent admits a pass of Condition~(II) at every prime factor, and neither $W_p$ passes Condition~(II).
\end{proposition}

\begin{proof}[Verification]
Computational, in the at-a-factor logic above: a factor failing Condition~(II) excludes the exponent from every local-pass-defined set and from any global pass. For each of the two factors, primality is certified by APR-CL, the divisibilities $2^p \equiv -1 \pmod r$, the pinned orders $\ord_r(2) = 2p$, the orders $\ord_r(\omega_3) = 4 d_r$, the factorizations of $d_r$, and the non-Class-III property of the listed prime factors ($\ord_{487}(2) = 243$ odd; $\ord_{41}(2) = 20$ with $v_2 = 2$; analogously for the two large factors) are verified by the archived scripts (\S\ref{sec:toolchain}), with an independent end-to-end recomputation in \texttt{scripts/cp352\_verify\_discharges\_recompute.py} that also re-tests Condition~(II) at each factor by direct congruence evaluation, bypassing the descent.
\end{proof}

\subsection{Multiplicity and Wieferich constraints}\label{sec:wieferich}

A repeated prime factor of $W_p$ is strongly constrained: the engine is an elementary valuation identity.

\begin{lemma}[Valuation transfer]\label{lem:w1-valuation}
For every prime $p \geq 5$ and every prime factor $r > 3$ of $W_p$:
\[
  v_r(W_p) \;=\; v_r\!\left(2^{\,r-1} - 1\right).
\]
\end{lemma}

\begin{proof}
Since $r > 3$, $v_r(W_p) = v_r(2^p + 1)$. Order-Pinning (Theorem~\ref{thm:order-pinning}) gives $\ord_r(2) = 2p$, so $r \mid 2^p + 1$ and $r \nmid 2^p - 1$; as $2^{2p} - 1 = (2^p - 1)(2^p + 1)$, this yields $v_r(2^{2p} - 1) = v_r(2^p + 1)$. By Lemma~\ref{lem:ord2}, $r \equiv 1 \pmod{2p}$; write $r - 1 = 2 p k$ with $1 \leq k < r$. The lifting-the-exponent lemma applied to $x^k - 1$ with $x = 2^{2p} \equiv 1 \pmod r$ gives
\[
  v_r\!\left(2^{\,r-1} - 1\right) = v_r\!\left((2^{2p})^k - 1\right) = v_r(2^{2p} - 1) + v_r(k) = v_r(2^{2p} - 1),
\]
since $0 < k < r$ forces $v_r(k) = 0$. Combining the two displays: $v_r(2^{\,r-1} - 1) = v_r(2^p + 1) = v_r(W_p)$.
\end{proof}

\begin{definition}[Super-Wieferich prime]\label{def:super-wieferich}
A prime $r$ is \emph{base-$2$ super-Wieferich} if $r^3 \mid 2^{\,r-1} - 1$, i.e.\ $v_r(2^{\,r-1} - 1) \geq 3$. No base-$2$ super-Wieferich prime is known in \emph{any} congruence class: the only base-$2$ Wieferich primes below $2^{64}$, namely $1093$ and $3511$, both have $v_r(2^{\,r-1} - 1) = 2$ exactly \cite{ref:OEISWieferich}.
\end{definition}

\begin{corollary}[Multiplicity forces Wieferich behaviour]\label{cor:multiplicity-wieferich}
Let $r > 3$ be a prime factor of $W_p$. If $v_r(W_p) \geq 2$, then $r$ is a base-$2$ Wieferich prime; if $v_r(W_p) \geq 3$, then $r$ is base-$2$ super-Wieferich.
\end{corollary}

In particular the prime-power configuration was doubly walled: a hypothetical $W_p = r^t$ has $r \equiv 3 \pmod 8$ and $t$ odd (Lemma~\ref{lem:mod8}), hence $t \geq 3$, so $r$ would be a base-$2$ super-Wieferich prime $\equiv 3 \pmod 8$ --- no such prime is known in any congruence class --- independently of the Nagell--Ljunggren exclusion of Lemma~\ref{lem:w1-empty}.

\begin{remark}[Repeated factors are Wieferich--Wagstaff]\label{rem:wieferich}
Consider a composite $W_p$ passing Condition~(II) at a repeated prime factor $r$, $v_r(W_p) \geq 2$ (by Lemma~\ref{lem:w1-empty} such a $W_p$ still has a second distinct prime factor). A prime $r$ with $v_r(W_p) \geq 2$ satisfies $r^2 \mid 3 W_p = 2^p + 1$; since $\ord_r(2) = 2 p$ (Order-Pinning, Theorem~\ref{thm:order-pinning}), the congruence $2^p \equiv -1 \pmod{r^2}$ forces $\ord_{r^2}(2) = 2 p$ as well, hence $\ord_{r^2}(2) = \ord_r(2)$, which is exactly the condition $2^{r-1} \equiv 1 \pmod{r^2}$ --- i.e.\ $r$ is a base-$2$ \emph{Wieferich prime}. This is the \emph{Wieferich--Wagstaff condition}. Passing Condition~(II) at the repeated factor in fact requires \emph{more} than that $r$ be base-$2$ Wieferich, and the extra requirement is half-automatic. Condition~(II) concerns $\omega_3 = \alpha^2$, and the inert factor already gives $\alpha^{N+1} \equiv -1 \pmod r$; write the modulo-$r^2$ defect as
\[
  \alpha^{N+1} \;=\; -1 + r\,(u + v\sqrt 2) \quad\text{in } \Z[\sqrt 2]/r^2, \qquad (u, v) \in (\Z/r)^2,
\]
so the factor passes Condition~(II) modulo $r^2$ iff $(u, v) = (0, 0)$. The \emph{rational} part vanishes unconditionally: since $N(\alpha) = -1$ and $N + 1 = 4 W_{p-2}$ is even, $N(\alpha^{N+1}) = (-1)^{N+1} = 1$, while $N(-1 + r(u + v\sqrt 2)) \equiv 1 - 2 r u \pmod{r^2}$, forcing $u \equiv 0 \pmod r$. The residual requirement is thus the single congruence $v \equiv 0 \pmod r$ --- equivalently $\ord_{r^2}(\omega_3) = \ord_r(\omega_3)$, the order not growing (the alternative $\ord_{r^2}(\omega_3) = 4 d_r\, r$ would force $4 d_r r \mid 4 W_{p-2}$, hence $r \mid W_{p-2}$, impossible as $\gcd(W_p, W_{p-2}) = 1$) --- a base-$(1+\sqrt 2)$, or \emph{Pell--Wall--Sun--Sun-type}, condition independent of the rational base-$2$ one. A repeated factor of a Condition-(II)-passing $W_p$ therefore carries two Wieferich-type conditions at once, of joint heuristic density $r^{-2}$ rather than $r^{-1}$; this makes the non-existence of such repeated factors still more strongly expected, though no unconditional proof follows --- each is a Wieferich-type non-existence statement, open in general. The only base-$2$ Wieferich primes below $2^{64}$ are $1093$ and $3511$, congruent to $5$ and $7 \pmod 8$ respectively, so no base-$2$ Wieferich prime $\equiv 3 \pmod 8$ occurs below $2^{64}$ \cite{ref:OEISWieferich}; consequently no inert repeated prime factor of any $W_p$ lies below $2^{64}$; and that a Condition-(II)-passing $W_p$ never carries a repeated factor at \emph{any} size is precisely the doubly-Wieferich non-existence expectation above --- conjectural, like every Wieferich-type non-existence statement.
\end{remark}

\begin{remark}[Pell--Wall--Sun--Sun condition for global passers]\label{rem:pell-wss-global}
The constraints above are local-pass statements, and Condition~(II) modulo $r^2$ does \emph{not} follow from local passes at the individual factors. If, however, $W_p$ moreover passes Condition~(II) \emph{globally} --- the only case Conjecture~\ref{conj:suff} needs --- then the congruence holds modulo $r^2$ for any repeated factor $r$, and the analysis of Remark~\ref{rem:wieferich} applies: the repeated factor additionally satisfies the Pell--Wall--Sun--Sun-type condition $\ord_{r^2}(\omega_3) = \ord_r(\omega_3)$, a second Wieferich-type event on top of the base-$2$ one. The Wieferich and super-Wieferich conclusions of Corollary~\ref{cor:multiplicity-wieferich} themselves require no global hypothesis. For \emph{squarefree} $W_p$ the situation is rigid in the opposite direction: by the Chinese Remainder isomorphism $\Z[\sqrt 2]/(W_p) \cong \prod_i \Z[\sqrt 2]/(r_i)$, the global Condition~(II) is \emph{exactly} the conjunction of the local passes --- under the isomorphism $-1 \bmod W_p$ is $(-1, \dots, -1)$ componentwise, so $\omega_3^{2 W_{p-2}} = -1$ holds modulo $W_p$ iff it holds in every component --- so a global pass adds nothing beyond the conjunction of local passes, and the local/global gap is confined to repeated factors.
\end{remark}

\subsection{Exact-AP characterization}

\begin{theorem}[Exact AP characterization]\label{thm:exact-ap}
Let $q>3$ be prime. If $q\notin A$, then $q\nmid W_{p-2}$ for every prime $p\geq5$. If $q\in A$, set $m_q=\ord_q(2)/2$ (odd); then for every prime $p\geq5$:
\[
  q \mid W_{p-2} \;\Longleftrightarrow\; p \equiv m_q + 2 \pmod{2 m_q}.
\]
In particular, for each $q \in A$ with $q > 3$ the set of admissible $p$ forms a single AP of period $2 m_q$; for $q \notin A$ the divisibility $q \mid W_{p-2}$ is impossible. The prime $q = 3 \in A$ is the lone exception to the period-$2 m_q$ form: by Theorem~\ref{thm:lte}, $3 \mid W_{p-2}$ iff $3 \mid p - 2$.
\end{theorem}

\begin{proof}
Since $q > 3$, $\gcd(q, 3) = 1$, so $q \mid W_{p-2} = (2^{p-2}+1)/3$ iff $q \mid 2^{p-2} + 1$ iff $2^{p-2} \equiv -1 \pmod q$. The latter forces $\ord_q(2) = 2 m_q$ with $m_q \mid p - 2$ and $m_q \nmid (p - 2)/2$, i.e., $v_2(\ord_q(2)) = 1$; $2^{m_q}$ has order 2 so equals $-1$. The condition $2^{p-2} \equiv -1 \equiv 2^{m_q} \pmod q$ is equivalent to $p - 2 \equiv m_q \pmod{2 m_q}$.
\end{proof}

\begin{remark}[Prime powers and the valuation-adjusted modulus]\label{rem:valuation-adjusted-Ld}
For composite $d$, the prime-level periods in Theorem~\ref{thm:exact-ap} must be lifted with valuations. If $q>3$, $q\in A$, $m_q=\ord_q(2)/2$, and $a_q=v_q(2^{m_q}+1)$, then for $e\geq1$ and odd $n$,
\[
  q^e \mid W_n
  \quad\Longleftrightarrow\quad
  n \equiv m_q q^{\max(0,e-a_q)}
  \pmod{2m_q q^{\max(0,e-a_q)}}.
\]
For $q=3$, Theorem~\ref{thm:lte} gives $v_3(W_n)=v_3(n)$, so $3^e\mid W_n$ iff $3^e\mid n$. Thus if
\[
  d=3^{e_3}\prod_{q>3}q^{e_q}
\]
is admissible, the exact condition $d\mid W_{p-2}$ is a single compatible system of congruences modulo the valuation-adjusted modulus
\[
  L_d=\operatorname{lcm}\!\left(
    3^{e_3},
    \;2m_q q^{\max(0,e_q-a_q)} \;:\; q>3,\ q^{e_q}\Vert d
  \right),
\]
with the corresponding residue class determined by the displayed congruences. The shorthand $M_d=\operatorname{lcm}(m_q:q\mid d)$ is useful only for squarefree or heuristic prime-level discussions; it is not the exact modulus for general $d$.
\end{remark}

\begin{proposition}[Small-$k$ valuation obstruction]\label{prop:small-k-valuation}
Let $r=2pk+1$ be an inert local-pass candidate, with $p\geq5$ prime, $r\equiv3\pmod8$, $d=d_r=\ord_r(\omega_3)/4$, and
\[
  d\mid \frac{r+1}{4}=\frac{pk+1}{2},
  \qquad
  d\mid W_{p-2}.
\]
Let $q^e\Vert d$.

If $q=3$, then
\[
  3^e\mid 2k+1.
\]
If $q>3$, then $q\in A$; with $m_q=\ord_q(2)/2$, $a_q=v_q(2^{m_q}+1)$, and $b=\max(0,e-a_q)$, one has
\[
  q^b\mid 2k+1.
\]
In particular, for $k=1$ one has $v_3(d)\leq1$ and $e\leq a_q$ for every prime power $q^e\Vert d$ with $q>3$.
\end{proposition}

\begin{proof}
Because $q$ is odd and $q^e\mid d\mid(pk+1)/2$, we have
\[
  pk\equiv -1 \pmod {q^e}.
\]
In particular $q\nmid k$.

For $q=3$, the condition $3^e\mid W_{p-2}$ is equivalent by Theorem~\ref{thm:lte} to $p\equiv2\pmod{3^e}$. Combining with $pk\equiv-1\pmod{3^e}$ gives $2k\equiv-1\pmod{3^e}$, hence $3^e\mid2k+1$.

For $q>3$, the divisibility $q^e\mid W_{p-2}$ forces $q\in A$ by Theorem~\ref{thm:exact-ap}. The valuation-adjusted congruence of Remark~\ref{rem:valuation-adjusted-Ld} gives
\[
  p-2\equiv m_q q^b \pmod {2m_q q^b},\qquad b=\max(0,e-a_q).
\]
Modulo $q^b$ this is $p\equiv2\pmod {q^b}$ (with nothing to show when $b=0$). Reducing $pk\equiv-1\pmod {q^e}$ modulo $q^b$ and substituting $p\equiv2$ yields $2k\equiv-1\pmod {q^b}$.

If $k=1$, the displayed divisibilities become $3^e\mid3$ and $q^b\mid3$ for every $q>3$, so $e\leq1$ for $q=3$ and $b=0$ for $q>3$, equivalently $e\leq a_q$.
\end{proof}

\subsection{Pair separation}\label{sec:pair-separation}

Proposition~\ref{prop:small-k-valuation} constrains each inert candidate separately. The same-$p$ pinning also couples two candidates at one exponent; until now the only pair-level constraint was the immediate consequence $2p \mid \gcd(r_1 - 1, r_2 - 1)$ of Order-Pinning. The next theorem is the first pair constraint beyond it. The mechanism is elementary: everything two candidates \emph{share} must divide their $k$-gap.

\begin{theorem}[Pair separation]\label{thm:pair-separation}
Let $p \geq 5$ be prime and let $r_i = 2 p k_i + 1$ ($i = 1, 2$) be distinct primes with $r_i \equiv 3 \pmod 8$ and $k_1 < k_2$. Let $d_1, d_2$ be odd positive integers with $d_i \mid (r_i + 1)/4$. Then:
\begin{enumerate}
\item $k_1 \equiv k_2 \equiv p \pmod 4$;
\item $\gcd(p k_1 + 1,\, p k_2 + 1) \mid k_2 - k_1$;
\item $4 \gcd(d_1, d_2) \mid k_2 - k_1$; equivalently, $8 p \gcd(d_1, d_2) \mid r_2 - r_1$;
\item if moreover $(d, r_1, p)$ and $(d, r_2, p)$ are danger triples with the same parameter $d = d_1 = d_2$ (Definition~\ref{def:danger-triple}; here $d_i = d_{r_i} = \ord_{r_i}(\omega_3)/4$, an odd divisor of $(r_i+1)/4$ by \S\ref{sec:gap-structure}), then $4 d \mid k_2 - k_1$ with $d \geq 57$; hence $k_2 - k_1 \geq 228$ and $r_2 - r_1 \geq 456\, p$.
\end{enumerate}
Items (1)--(3) use neither Condition~(II) nor the primality of $p$: they hold verbatim for any odd $p \geq 3$ and any odd $d_i \mid (r_i + 1)/4$. Item (4) is where the Condition-(II) context enters, through $d > 43$ (Pell-sequence exclusion) and admissibility of $d$ (every prime factor in $A$, forced by $d \mid W_{p-2}$ and Theorem~\ref{thm:exact-ap}).
\end{theorem}

\begin{proof}
(1) $r_i \equiv 3 \pmod 8$ means $2 p k_i \equiv 2 \pmod 8$, i.e.\ $p k_i \equiv 1 \pmod 4$. Since $p$ is odd, $p^2 \equiv 1 \pmod 4$, so multiplying by $p$ gives $k_i \equiv p \pmod 4$ for both $i$; in particular $4 \mid k_2 - k_1$.

(2) The integer identity $k_2 (p k_1 + 1) - k_1 (p k_2 + 1) = k_2 - k_1$ shows that every common divisor of $p k_1 + 1$ and $p k_2 + 1$ divides $k_2 - k_1$.

(3) Set $g = \gcd(d_1, d_2)$, an odd integer. First, $\gcd(g, p) = 1$: a common prime divisor $q$ of $g$ and $p$ would divide $d_1 \mid (r_1 + 1)/4 = (p k_1 + 1)/2$, hence $p k_1 + 1$, and also $p k_1$ (as $q \mid p$), forcing $q \mid 1$. Next, $g$ divides both $(p k_i + 1)/2$, hence their difference $p (k_2 - k_1)/2$; by (1) we have $4 \mid k_2 - k_1$, so $(k_2 - k_1)/2$ is an (even) integer, and coprimality to $p$ gives $g \mid (k_2 - k_1)/2$, i.e.\ $2 g \mid k_2 - k_1$. Combining $2 g \mid k_2 - k_1$ with $4 \mid k_2 - k_1$ and $g$ odd yields $4 g \mid k_2 - k_1$. Multiplying by $2 p$: $8 p g \mid 2 p (k_2 - k_1) = r_2 - r_1$.

(4) Here $d_1 = d_2 = d$, so (3) gives $4 d \mid k_2 - k_1$, and $k_2 > k_1$ gives $k_2 - k_1 \geq 4 d$. In a danger triple, $d > 43$ and every prime factor of $d$ lies in $A$; the least odd integer $> 43$ with all prime factors in $A$ is $57 = 3 \cdot 19$: among $45, \dots, 55$ the odd values $45 = 3^2 \cdot 5$, $49 = 7^2$, $51 = 3 \cdot 17$, $55 = 5 \cdot 11$ contain the non-Class-III primes $5, 7, 17$, while $47 \notin A$ ($\ord_{47}(2) = 23$, odd) and $53 \notin A$ ($v_2(\ord_{53}(2)) = 2$). Hence $d \geq 57$, $k_2 - k_1 \geq 228$, and $r_2 - r_1 = 2 p (k_2 - k_1) \geq 456\, p$.
\end{proof}

The next proposition transfers valuations \emph{across} the pair. It strictly strengthens Proposition~\ref{prop:small-k-valuation} in two ways: part (a) lets the factor with the higher $q$-valuation constrain the other factor's $k$, and part (b) routes a constraint through the half-order of one prime divisor of $d$ acting on another.

\begin{proposition}[Pair-coupled valuation transfer]\label{prop:pair-couplings}
In the setting of Theorem~\ref{thm:pair-separation}, assume additionally that $d_i \mid W_{p-2}$ for both $i$ (condition (v) at both candidates). Then every prime $q > 3$ dividing $d_1 d_2$ lies in $A$; write $m_q = \ord_q(2)/2$ and $a_q = v_q(2^{m_q} + 1)$ as in Remark~\ref{rem:valuation-adjusted-Ld}. For every such $q$:

(a) (\emph{cross-factor channel}) with $e_i = v_q(d_i)$, $E = \max(e_1, e_2)$ and $b = \max(0,\, E - a_q)$:
\[
  q^{\min(e_i,\, b)} \;\Big|\; 2 k_i + 1 \qquad (i = 1, 2);
\]

(b) (\emph{half-order channel}) for every prime $\ell$ with $v = v_\ell(m_q) \geq 1$ and $e = v_\ell(d_i) \geq 1$ for some $i$:
\[
  \ell^{\min(v,\, e)} \;\Big|\; 2 k_i + 1.
\]
\end{proposition}

\begin{proof}
Both $d_i$ divide $W_{p-2}$, hence so does $\operatorname{lcm}(d_1, d_2)$; in particular every prime $q > 3$ dividing $d_1 d_2$ divides $W_{p-2}$ and so lies in $A$ (Theorem~\ref{thm:exact-ap}).

(a) $q^E \mid \operatorname{lcm}(d_1, d_2) \mid W_{p-2}$, so the valuation-adjusted congruence of Remark~\ref{rem:valuation-adjusted-Ld} gives $p - 2 \equiv m_q q^{b} \pmod{2 m_q q^{b}}$; since $q \nmid m_q$, reducing modulo $q^b$ gives $p \equiv 2 \pmod{q^b}$. For each $i$, $q^{e_i} \mid d_i \mid (p k_i + 1)/2$ gives $p k_i \equiv -1 \pmod{q^{e_i}}$. Reducing both congruences modulo $q^{\min(e_i, b)}$ and substituting $p \equiv 2$ yields $2 k_i \equiv -1$.

(b) $q \mid W_{p-2}$ gives $m_q \mid p - 2$ (Lemma~\ref{lem:classIII-division}), hence $p \equiv 2 \pmod{\ell^{v}}$. Since $d_i$ is odd, $\ell$ is odd, and $\ell^{e} \mid d_i \mid (p k_i + 1)/2$ gives $p k_i \equiv -1 \pmod{\ell^{e}}$. Reduce modulo $\ell^{\min(v, e)}$ and substitute $p \equiv 2$.
\end{proof}

\begin{remark}[Strict strengthenings of the small-$k$ obstruction]\label{rem:couplings-strengthen}
Proposition~\ref{prop:small-k-valuation} is the single-candidate shadow of Proposition~\ref{prop:pair-couplings}: its exponent $\max(0, e_i - a_q)$ never exceeds the $\min(e_i, b)$ of part (a), and part (a) is strictly stronger whenever the \emph{other} factor carries the higher valuation ($e_i < E$). Part (b) needs no pair at all --- its proof uses only $q \mid W_{p-2}$ --- so it already strengthens Proposition~\ref{prop:small-k-valuation} for a single candidate $(r, p, k, d)$ with $q \mid d$. Example: $d = 737 = 11 \cdot 67$ is admissible with $a_{11}, a_{67} \geq 1$, so Proposition~\ref{prop:small-k-valuation} imposes no constraint at $k = 1$; but $m_{67} = 33 = 3 \cdot 11$, and part (b) with $q = 67$, $\ell = 11$ forces $11 \mid 2k + 1$, killing $k = 1$ (and every $k \not\equiv 5 \pmod{11}$).
\end{remark}

\begin{remark}[No off-diagonal kill at the congruence level]\label{rem:no-congruence-kill}
Two honest limits, both witnessed constructively. (i) The natural strengthening of Proposition~\ref{prop:small-k-valuation} to $\rad(d) \mid 2k + 1$ is \emph{false}: the inert prime $r = 43 = 2 \cdot 7 \cdot 3 + 1$ has $\ord_{43}(\omega_3) = 44$, giving $(r, p, k, d) = (43, 7, 3, 11)$ with $11 \mid W_5$, yet $11 \nmid 2k + 1 = 7$. The reason is structural: the progression condition on $p$ (modulus $2 m_q q^{b}$) and the divisor congruence $p k \equiv -1 \pmod q$ live on coprime moduli unless a channel of Proposition~\ref{prop:pair-couplings} happens to link them. (ii) No combination of the constraints of this subsection rules out bounded-$k$ \emph{off-diagonal} pairs: at $(k_1, k_2) = (1, 5)$, $(d_1, d_2) = (57, 59)$, every prime $p \equiv 166{,}781 \pmod{1{,}170{,}324}$ (the displayed residue is itself prime) satisfies simultaneously every congruence imposed by items (1)--(3), by Proposition~\ref{prop:small-k-valuation}, by Proposition~\ref{prop:pair-couplings}, and by condition (v) for both $d_i$ (Theorem~\ref{thm:exact-ap}: $3 \mid p - 2$, $p - 2 \equiv 9 \pmod{18}$, $p - 2 \equiv 29 \pmod{58}$). These congruences force only the modulus $2^2 \cdot 3^2 \cdot 29 = 1{,}044$; the finer modulus $1{,}170{,}324 = 1{,}044 \cdot 19 \cdot 59$ of the displayed witness carries the realizability factors that fix the prescribed $\omega_3$-orders ($d_1 = 57 = 3 \cdot 19$, $d_2 = 59$), not additional forced congruences. What the witness class does not supply is \emph{realizability}: a prime $p$ in the class for which $2p + 1$ and $10 p + 1$ are both prime, both inert, with the prescribed $\omega_3$-orders. Separating congruence-consistency from realizability at a fixed divisor is exactly the analytic wall of \S\ref{sec:weaker-blocked}: the separation theorem yields repulsion for free, while attraction --- what PPPC must exclude --- is untouched at this level.
\end{remark}

\begin{remark}[The $r + 1$ side of pair gcds]\label{rem:gcd-rplus1}
For any pair as in Theorem~\ref{thm:pair-separation}, $v_2(r_i + 1) = 2$ and $(r_i + 1)/4 = (p k_i + 1)/2$ is odd, so
\[
  \gcd(r_1 + 1,\, r_2 + 1) \;=\; 2 \gcd(p k_1 + 1,\, p k_2 + 1),
\]
which is $4$ times an odd integer and divides $2 (k_2 - k_1)$ by item (2). This explains a regularity of the candidate-pair scans (\S\ref{sec:empirical}): while the $r - 1$ side carries the forced large gcd $2 p \mid \gcd(r_1 - 1, r_2 - 1)$, the observed values of $\gcd(r_1 + 1, r_2 + 1)$ on real candidate pairs are tiny --- $4$ and $12$ in the inert pairs examined --- being capped by twice the $k$-gap rather than scaled by $p$.
\end{remark}

\begin{remark}[Cross-factor cofactor congruence]\label{rem:cofactor-congruence}
At an inert prime factor $r$ of a composite $W_p$, with $d_r = \ord_r(\omega_3)/4$ as in \S\ref{sec:gap-structure}, Condition~(II) at $r$ is \emph{equivalent} to a congruence on the complementary cofactor:
\[
  \text{Condition~(II) holds at } r
  \iff
  \frac{W_p}{r} \equiv 1 \pmod{4 d_r}.
\]
Forward: the pass gives $d_r \mid W_{p-2}$, hence $4 d_r \mid 4 W_{p-2} = W_p + 1$ (Lemma~\ref{lem:fundamental}), while the gap structure gives $4 d_r \mid r + 1$; thus $W_p \equiv r \equiv -1 \pmod{4 d_r}$, $r$ is invertible with $r^{-1} \equiv -1$, and the cofactor is $\equiv (-1)(-1) = 1$. Backward: $W_p/r \equiv 1$ and $r \equiv -1$ force $W_p \equiv -1 \pmod{4 d_r}$, i.e.\ $d_r \mid W_{p-2}$, which is the pass. The congruence holds modulo $8 d_r$ for free ($W_p \equiv r \equiv 3 \pmod 8$, Lemma~\ref{lem:mod8}). At an all-inert passer with $\omega = 3$ this becomes a closure system: for each $i$, the \emph{other two} factors are mutual inverses, $r_j r_l \equiv 1 \pmod{4 d_{r_i}}$. The congruence is verified at all three Platinum witness factors and at the $d = 67$ danger factor of Remark~\ref{rem:known-danger}. Like the pair-separation constraints of Theorem~\ref{thm:pair-separation}, these are perimeter, not a tag: the system has nonempty congruence-consistent classes (it transfers to the function-field calibration verbatim), and for $\omega \geq 3$ its use as an exclusion is circular without new input.
\end{remark}

\section{The Pell Primitive Pair Conjecture}\label{sec:pppc}

\subsection{Statement}

The principal residual of the paper --- the above-cutoff inert-pair fragment of SPC --- is the Pell Primitive Pair Conjecture, already stated informally in \S\ref{sec:pppc-statement}. For completeness:

\begin{conjecture}[Pell Primitive Pair Conjecture, PPPC]\label{conj:pppc}
For every prime $p \geq 5$: $|T_p^{>}| \leq 1$, where $T_p^{>}$ is the set of primes $r$ with
\[
  r \equiv 3 \pmod 8,\quad \ord_r(2) = 2 p,\quad r > 10^{12},\quad d_r > 43,\quad d_r \mid W_{p-2}.
\]
\end{conjecture}

Equivalently, no two distinct primes $r_1, r_2$ of this form exist for the same $p$.

\subsection{Equivalent forms}\label{sec:equiv-forms}

PPPC admits several equivalent formulations, each useful in different contexts. The diagonal case $d_1=d_2$ is part of the statement: two distinct primes in $T_p^{>}$ may have the same order parameter $d_r$.

\textbf{(a) Pair-count form.} $\Pi := \sum_p \binom{|T_p^{>}|}{2} = 0$.

\textbf{(b) Primitive-divisor form.} No prime $p$ admits two distinct primitive divisors $r_1,r_2$ (of $V_{2d_1}$ and $V_{2d_2}$, with $d_1,d_2>43$ not necessarily distinct) such that
\[
  r_i \equiv 3 \pmod 8,\qquad \ord_{r_i}(2)=2p,\qquad
  r_i>10^{12},\qquad d_i\mid W_{p-2}.
\]
Equivalently, both the off-diagonal coincidences $d_1<d_2$ and the diagonal multiplicities at a fixed $d$ vanish.

\textbf{(c) AP-coincidence form.} Let $S_d$ be the set of primes $p$ satisfying the valuation-adjusted congruence condition equivalent to $d\mid W_{p-2}$ (prime moduli are described in Theorem~\ref{thm:exact-ap}; prime powers are governed by the LTE refinement below it). Then PPPC asserts that the pinned-order map attached to the primitive inert divisors of $V_{2d}$ has no common value across $S_{d_1}$ and $S_{d_2}$ for $d_1<d_2$, and is injective on the relevant primitive-divisor set when $d_1=d_2$.

\textbf{(d) Number-field form.} Let $\mathcal{P}_d^{\mathrm{prime}} \subset \Spec\,\Z$ be the set of primes $r$ that are primitive divisors of $V_{2 d}$ with $r \equiv 3 \pmod 8$, $r > 10^{12}$, $p_r:=\ord_r(2)/2$ prime, and $d\mid W_{p_r-2}$. The maps
\[
  \mathcal{P}_d^{\mathrm{prime}} \longrightarrow \Z,\qquad r \longmapsto p_r
\]
have pairwise disjoint images across distinct admissible $d>43$, and each map is injective on each fixed $d$.

Both coordinates of a hypothetical pair are rigid: by Multi-Factor Pinning both members share the pinned exponent, and each condition $d_{r_i} \mid W_{p-2}$ is the valuation-adjusted congruence of Theorem~\ref{thm:exact-ap} and Remark~\ref{rem:valuation-adjusted-Ld}. Both coordinates also carry unconditional pair-level constraints. Any two members of $T_p^{>}$ at one $p$ are separated, by the Pair Separation Theorem (Theorem~\ref{thm:pair-separation}), at scale $8 p \gcd(d_{r_1}, d_{r_2})$ in $r$ --- at scale $3208\,p$ on the diagonal, where the census floor $d > 400$ applies; and the diagonal multiplicities of form (c)/(d) vanish entirely for every admissible $d \leq 400$, for all $p$ and all $k$, by the certified closures of Propositions~\ref{prop:diagonal-closures} and~\ref{prop:danger-census-400} --- in particular, no diagonal pair with both $k_i = (r_i - 1)/2p \leq 1604$ exists (Corollary~\ref{cor:diagonal-k500}). The residue-level structure of the diagonal sub-case, and the question that remains open there, are developed in Appendix~\ref{sec:pppc-problem} and \S\ref{sec:pppc-directions}.

\subsection{Empirical evidence}\label{sec:empirical}

\textbf{Small-factor enumeration.} The Platinum enumeration of \S\ref{sec:inert-survey} records every inert factor $r \leq 10^{12}$ in the stated range and tests Condition~(II). It verifies that, for each $p$, at most one such \emph{small} inert factor passes locally. This is not a verification of PPPC, whose set $T_p^{>}$ is defined by the opposite cutoff $r>10^{12}$. Its role in the reduction (Theorem~\ref{thm:sharpened-chain}) is sharper and more limited: in Case B it caps the number of sub-cutoff passing factors at one, pinning any remaining double pass above the cutoff.

\textbf{Extended $p$-indexed enumeration.} Beyond the Platinum range $r \leq 10^{12}$, a $p$-indexed scan covered $p^\ast \in [5 \cdot 10^{11},\, 5 \cdot 10^{11} + 10^9]$ ($37{,}123{,}274$ primes). For each $p^\ast$ and each $k = 1, 2, \dots, k_{\max}$, the candidate $r = 2 p^\ast k + 1$ was tested for primality, residue $r \equiv 3 \pmod 8$, and order $\ord_r(2) = 2 p^\ast$; the pre-condition-(v) window
\[
  \begin{aligned}
  T_{p^\ast}^{\leq k_{\max},\,\mathrm{pre}} := \{\,r = 2 p^\ast k + 1 :\;& 1 \leq k \leq k_{\max},\; r \text{ prime},\\
  & r \equiv 3 \pmod 8,\; \ord_r(2) = 2 p^\ast\,\}.
  \end{aligned}
\]
contains the candidates before the PPPC condition $d_r \mid W_{p^\ast-2}$ is applied. Intersecting it with $r>10^{12}$, $d_r>43$, and condition (v) gives the portion of $T_{p^\ast}^{>}$ visible in the window; the inclusion is therefore from $T_{p^\ast}^{>}$ into this pre-window, not conversely. Distribution at $k_{\max} = 500$ on the $37$M-prime sample:

\begin{center}
\begin{tabular}{r r r}
\hline
$|T_{p^\ast}^{\leq 500,\mathrm{pre}}|$ & count & fraction \\
\hline
0 & $33{,}052{,}256$ & $89.03\%$ \\
1 & $3{,}886{,}080$ & $10.47\%$ \\
2 & $180{,}116$ & $0.485\%$ \\
3 & $4{,}750$ & $0.0128\%$ \\
4 & $72$ & $0.000194\%$ \\
\hline
\end{tabular}
\end{center}

Sample mean $\lambda = 0.1148$, sample variance $0.1121$, ratio $\Var/\text{mean} = 0.977$ (compared to Poisson's $1.0$); $\chi^2$ against $\mathrm{Poisson}(\lambda)$ is approximately $10^4$, indicating significantly sub-Poisson behaviour. Extending to $k_{\max} = 5000$ on a 1M-prime subsample multiplies the pair fraction by $1.75 \times$ and the triple fraction by $2.5 \times$, consistent with a non-degenerate but slowly-growing tail.

\textbf{Selection-bias note: condition (v) filter.} At $184{,}938$ exponents of the $37$M-prime sample --- the rows $|T_{p^\ast}^{\leq 500,\mathrm{pre}}| \geq 2$ of the table above, $180{,}116 + 4{,}750 + 72$ --- the $k_{\max} = 500$ window contains at least two candidates satisfying conditions (i)--(iv) of PPPC ($r_i \equiv 3 \pmod 8$, $r_i > 10^{12}$, $\ord_{r_i}(2) = 2 p^\ast$, primality). At each such exponent a candidate pair $(r_1, r_2)$ was recorded, and independent verification computed condition (v) --- $d_{r_i} \mid W_{p^\ast - 2}$ --- for both members of every recorded pair via $\F_{r_i^2} = \F_{r_i}[\sqrt 2]$ arithmetic: \textbf{$0$ of the $369{,}876$ tested $r_i$ satisfy (v).} The conditional density of (v) given (i)--(iv) is thus $\ll 10^{-5}$ empirically, consistent with the heuristic $E \approx 10^{-3}$ of \S\ref{sec:heuristic} and with $\Pi = 0$. In every tested case the failure is of condition (v) itself --- the order-alignment $d_{r_i} \mid W_{p^\ast - 2}$, equivalently $2^{p^\ast - 2} \equiv -1 \pmod{3 d_{r_i}}$ --- with conditions (i)--(iv) holding by construction of the candidate list. A deep-$k$ complement to this record was computed at $30{,}000$ exponents $p^\ast$ from $5 \times 10^{11}$ with the window extended two thousandfold to $k \leq 10^6$: $6{,}559$ candidates satisfy (i)--(iv), and again \emph{zero} satisfy condition (v); $623$ of these exponents carry two or more (i)--(iv) candidates, and every such pair satisfies the unconditional separation constraints of Theorem~\ref{thm:pair-separation} (script \texttt{scripts/}\allowbreak\texttt{cp358\_falsification\_prewindow.py}).

\textbf{Methodological remark.} The $d$-indexed enumeration (\S\ref{sec:inert-survey} Platinum, \S\ref{sec:multi-factor} primitive-divisor survey, \S\ref{sec:nct-verify} sieve) and the $p$-indexed enumeration above probe complementary slices of the PPPC search space: $d$-indexed scans capture pairs where both factors have small order parameter, while $p$-indexed scans capture the arithmetic progression $r \equiv 1 \pmod{2 p^\ast}$ before condition (v) is applied. The $d \leq 1000$ catalog of the primitive-divisor survey covers a small fraction of the $p$-indexed pre-window at $p^\ast = 5 \times 10^{11}$ --- most of the $r$'s appearing in the $T_{p^\ast}^{\leq 500,\mathrm{pre}}$ distribution above have $d_r$ well above $1000$. Absence of pairs in the $d \leq 1000$ catalog is rigorous within scope but is not directly comparable with the $p$-indexed pre-window density; both lines of empirical evidence are needed.

\subsection{Heuristic support}\label{sec:heuristic}

\textbf{Heuristic count.} Modelling each primitive inert divisor of $V_{2 d}$ as a random coincidence with $\ord_r(2)/2 = p$ among its admissible residues, the expected total number of single danger triples $(d, r, p)$ with $d > D$ is
\[
  E(D) \;\approx\; \sum_{d > D,\; d \text{ admissible}} \frac{(\ln \ln |V_{2 d}|)^2}{2 \ln |V_{2 d}| \cdot \varphi(L_d)},
\]
where $L_d$ is the valuation-adjusted modulus of Remark~\ref{rem:valuation-adjusted-Ld} and $|V_{2 d}| \sim (1 + \sqrt 2)^{2 d}$. For admissible prime $d = q$ with the generic valuation $a_q=1$, $L_d = 2m_q$ and typically $\varphi(L_d)$ is of order $q$, giving a rapidly convergent prime-level model. The displayed expression carries an implicit \emph{inert-fraction} factor of $\approx 1/2$ --- only the primitive divisors lying in the inert class $r \equiv 3 \pmod 8$ are relevant, and the displayed sum itself evaluates to $\approx 0.0076$ at $D = 1000$ before that factor is applied. Numerical evaluation at $D = 1000$ in the squarefree prime-level model, with the inert fraction included, gives $E \approx 0.003$, consistent with zero observed.

\textbf{A formalized pair-count model for PPPC (heuristic only).} The
pair heuristic can be made an explicit expected-count computation. The
model has three assumptions, stated here once and numbered; \emph{none is
a theorem}, and nothing in this paragraph proves anything about PPPC.

\begin{itemize}
\item[\textbf{M1}] (\emph{Candidate production.}) At a prime exponent
$p$, the candidates satisfying conditions (i)--(iv) of
Conjecture~\ref{conj:pppc} --- the primes $r = 2 p k + 1 > 10^{12}$ with
$k \equiv p \pmod 4$ (Theorem~\ref{thm:pair-separation}(1)) --- arrive
independently with intensity $\Prob[\text{(i)--(iv) at } k] \approx
C_{\mathrm{cal}} / (2 k \ln(2 p k))$: the factor $1/(2k) = p/(r-1)$ is
the order-alignment probability ($\ord_r(2) = 2p$), and the single
calibration constant $C_{\mathrm{cal}}$ absorbs the AP-primality and
mod-$8$ corrections. The measured window mean of \S\ref{sec:empirical},
$\lambda(500) = 0.1148$ at $p^\ast \approx 5 \times 10^{11}$, gives
$C_{\mathrm{cal}} = 3.69$; the model then predicts
$\lambda(10^6) = 0.208$ against the measured $6559/30000 = 0.219$
(ratio $0.95$). Integrated over the full factor spectrum, the intensity
gives $\mu_p \approx \tfrac12(\ln \ln W_p - \ln \ln \max(2p, 10^{12}))$
expected above-cutoff candidates per exponent ($\mu_{p^\ast} \approx
11.6$), the normal-order count of inert factors.
\item[\textbf{M2}] (\emph{Condition-(v) density.}) Given (i)--(iv), the
alignment $d_r \mid W_{p-2}$ holds with probability $\delta$, anchored on
the paper's two measured records in three readings:
\[
  \delta_{\mathrm{ceil}} = \frac{3}{376{,}435} = 8.0 \times 10^{-6},
  \qquad
  \delta_{\mathrm{cen}} = \frac{3}{15{,}587{,}021} = 1.9 \times 10^{-7},
\]
and $\delta_{\mathrm{best}}(r) = \delta_{\mathrm{cen}} \cdot
\min\!\big(1,\, 10^{12}/r\big)$. The ceiling is the $95\%$ rule-of-three
bound from the above-cutoff zero-pass record ($369{,}876 + 6{,}559$
tested candidates, $0$ passes, \S\ref{sec:empirical}); the central value
is the only nonzero measured density --- the Platinum enumeration's $3$
passes among the $15{,}587{,}021$ rows surviving the $G_r \leq 43$
exclusion at $r \leq 10^{12}$ (\S\ref{sec:inert-survey}); the best
estimate decays the measured density like $1/r$ beyond the measurement
scale (the alignment mass of divisors $d_r \mid (r+1)/4$ against the
fixed integer $W_{p-2}$), the same structural-decay device as the MPC
heuristic's best column.
\item[\textbf{M3}] (\emph{Independence.}) Candidate production and the
(v)-events at distinct candidates are independent, so the pass count
$X_p = |T_p^{>}|$ is Poisson with mean $\nu_p = \sum_k
\Prob[\text{(i)--(iv)}]\,\delta$. The measured variance ratio $0.977$
(sub-Poisson, \S\ref{sec:empirical}) and the unconditional repulsions of
Theorem~\ref{thm:pair-separation} and of the diagonal closures
(Proposition~\ref{prop:diagonal-closures}, modulo their stated
certificates) --- neither credited by the model --- push the true pair
count below the model's.
\end{itemize}

\emph{Per-exponent violation probability.} Under M1--M3,
$\Prob[|T_p^{>}| \geq 2] \leq \mathbb{E}\binom{X_p}{2} = \nu_p^2/2$. At
the Platinum frontier $p^\ast = 5 \times 10^{11}$:
$\nu_{p^\ast} = 9.3 \times 10^{-5}$ $/$ $2.2 \times 10^{-6}$ $/$
$7.8 \times 10^{-9}$ (ceiling / central / best), giving per-exponent
pair probabilities $\nu_{p^\ast}^2/2 = 4.3 \times 10^{-9}$ $/$
$2.5 \times 10^{-12}$ $/$ $3.1 \times 10^{-17}$.

\emph{Summed expectation.} Under the flat readings the $p$-indexed sums
diverge --- with constant $\delta$ and $\mu_p \sim \tfrac12 \ln\ln W_p$,
single passes recur infinitely often, as expected: the enumeration
already exhibits three witness exponents --- so those readings are
summed through the $d$-indexed form of the count, each order parameter
$d$ contributing its convergent single-triple summand from the display
above (inert fraction included) times a partner factor
$\min(1,\, \omega_{\mathrm{in}}(\bar p_d)\, \delta)$,
$\omega_{\mathrm{in}}(x) = \tfrac12(\ln x - \ln\ln x)$, at the pinned
scale $\bar p_d = V_{2 d}/4$ --- the same partner device as the MPC
heuristic. The best reading needs no such cap: its $1/r$ decay makes
$\sum_p \nu_p^2/2$ converge outright, dominated within a decade of the
anchor scale $p^\ast$ --- exactly where the calibration is measured. The
results:
\[
  \mathbb{E}[\Pi] \;\approx\;
  1.1 \times 10^{-3} \;\text{(ceiling)}, \qquad
  \sim 10^{-4} \;\text{(central)}, \qquad
  6.0 \times 10^{-7} \;\text{(best)}.
\]
For self-consistency: the central reading predicts $0.07$ expected
(v)-passes among the $369{,}876$ tested candidates (observed: $0$), and
the ceiling reading predicts $3$ (it \emph{is} the $95\%$ bound built
from that record).
The three-orders spread between ceiling and best is genuine model
uncertainty --- the cost of extrapolating a condition-(v) density
measured at $r \leq 10^{12}$ to factors of unbounded size --- and is
reported bracketed rather than resolved, exactly as for the MPC joint
heuristic below. Constants and calibration gates:
\texttt{scripts/cp395\_pppc\_heuristic\_constants.py} (the script
reproduces the $E(1000)$ figures and window means above before computing
the bracket). These figures are heuristics only; they prove nothing
about PPPC, which remains an open conjecture, and
Conjecture~\ref{conj:suff} is expected to hold for the structural reason
described below, not because of these numbers. The rigorous proof of
$\Pi = 0$ is open.

\textbf{Joint heuristic for MPC (heuristic only).} The Mixed Pair Conjecture (Conjecture~\ref{conj:mpc}) admits an analogous joint count: a violation requires a split compatible triple \emph{and} an inert local pass at the same pinned $p$. Summing the joint expectation over admissible $d_s > 199$, with the then-open trio $\{131, 163, 179\}$ carried as separate line items (all three have since been closed, \S\ref{sec:nct-verify}; their retained line items contribute negligibly to the totals), gives an expected violation count $E_{\mathrm{MPC}}$ with ceiling $8.7 \times 10^{-4}$, central value $4.4 \times 10^{-4}$, and best estimate $9.8 \times 10^{-10}$, against an NCT-general \emph{single-event} heuristic of $1.27 \times 10^{-2}$ (ceiling) and $2.4 \times 10^{-5}$ (structural) computed in the same framework. The joint heuristic is thus 1--5 orders of magnitude below the corresponding NCT-general single-event heuristic (model-dependent), approaching PPPC's pair scale $O(10^{-6})$ only in the structural-decay reading. The ceiling/best-estimate spread of six orders is genuine model uncertainty --- it extrapolates the empirically measured condition-(v) density from $r \sim 10^{12}$ to factors of size $e^{200}$ and beyond --- and is reported bracketed rather than resolved. These figures are heuristics only; they prove nothing about MPC, which remains an open conjecture.

A structural ingredient of the split branch deserves separate mention: in any compatible triple, the unconditionally proved band $8 p d \mid r_s - 1$ (Proposition~\ref{prop:nct-band}) forces the index of $2$ in $\F_{r_s}^\times$ --- namely $(r_s - 1)/\ord_{r_s}(2) = (r_s - 1)/2p$ --- to absorb the entire $d$-part of $r_s - 1$, the order of $2$ being coprime to $d$. In the index model this is an extra alignment of heuristic cost $\sim 1/(3d)$ per split factor, a suppression with no inert analogue (the inert order datum lives on the $r + 1$ side, where no corresponding index constraint arises).

Assembled in the SPC framing (\S\ref{sec:spc}): the expected number of SPC violations is the sum of the three fragment expectations --- inert pair, mixed pair, split--split --- and is dominated by the computed PPPC and MPC brackets; the split--split term carries no observed candidates at all: no compatible triple is known.

\textbf{The conjecture is structural, not probabilistic.} The rapid convergence above is an artefact of the pinning constraints and would not survive their removal. The Chua congruence is a strong Lucas test with parameter $Q = \mathrm{N}(\omega_3) = +1$; for a Lucas or Frobenius test with $|Q| = 1$ the standard pseudoprime heuristic predicts \emph{infinitely many} composite numbers passing the single congruence, absent further structure (cf.\ \cite{ref:Grantham}). What forces $\Pi$ to be small is not a generic scarcity of pseudoprimes but the Wagstaff-specific Order-Pinning (Theorem~\ref{thm:order-pinning}) and Multi-Factor Pinning (Theorem~\ref{thm:multi-pinning}): every prime factor of one $W_p$ shares a single pinned exponent $p$ \emph{exactly}, and each divisibility $d_r \mid W_{p-2}$ is the rigid valuation-adjusted congruence condition of Theorem~\ref{thm:exact-ap} and Remark~\ref{rem:valuation-adjusted-Ld}. Conjecture~\ref{conj:suff} is therefore expected to hold for a structural reason rather than a probabilistic one --- which is also why a purely analytic count treating the factors as independent random primes cannot reach it, the precise content of \S\ref{sec:weaker-blocked}.

\subsection{Connection to classical problems}

PPPC is structurally related to several classical problems on Pell--Lucas sequences:

\begin{itemize}
\item \textbf{Primitive-divisor structure of $V_{2 d}$.} Bilu--Hanrot--Voutier \cite{ref:BHV} guarantees a primitive divisor for $V_n$ for all $n \geq 30$. Stewart--Tijdeman \cite{ref:Stewart} gives $\omega(V_n) \ll n / \log n$. A sharpening of the Stewart bound at small $n \cdot \log \log n$ scales could contribute to bounding components of the pair count $\Pi$ (attack angle N2).

\item \textbf{Prime divisors of $a^2 + 2 b^2$.} By Remark~\ref{rem:zsigmondy}, any $r \in T_p^{>}$ divides $N_d = V_d^2 + 2$, so $r$ is a prime of form $a^2 + 2 b^2$. Iwaniec's half-dimensional sieve \cite{ref:Iwaniec} bounds the count of such primes in arithmetic progressions (attack angle N3).

\item \textbf{Effective Chebotarev density.} The Artin-like question ``for a fixed $\alpha \in \Z[\sqrt 2]^\times$, how often does a primitive divisor of $V_n$ fall in a prescribed AP?'' is related to Chebotarev on the tower $\Q(\alpha^{1/m})$; effective forms (Lagarias--Odlyzko; under GRH, Hooley--Heath-Brown \cite{ref:Hooley, ref:HeathBrown}; see the survey \cite{ref:Moree}) could supply inputs relevant to AP components of $\Pi$, but do not by themselves control the fixed-divisor count.

\item \textbf{Bombieri--Friedlander--Iwaniec-type levels of distribution.} An analogue of the Bombieri--Friedlander--Iwaniec level-of-distribution technology for the non-standard Lucas sequence $\{V_{2 d} : d \leq X\}$ would be relevant to components of the pair count (attack angle N4). No such direct application is used here.
\end{itemize}

Section~\ref{sec:pppc-problem} elaborates these connections.

\section{The reduction theorems}\label{sec:reduction}

\subsection{Branch decomposition}

The members of $P_p$ at composite $W_p$ are exactly the objects
classified in Sections~\ref{sec:split} and~\ref{sec:inert}; the
following proposition packages the classification and is the bridge
between SPC and the named conjectures.

\begin{proposition}[Branch decomposition of a double pass]\label{prop:branch-split}
Let $p \geq 5$ be prime and suppose $|P_p| \geq 2$; pick distinct
$r_1, r_2 \in P_p$. Then $W_p$ is composite, $r_1, r_2$ are distinct
prime factors of $W_p$, each $\equiv 1$ or $3 \pmod 8$, and, after
ordering, exactly one of the following holds.
\begin{itemize}
\item[(a)] \emph{Inert pair:} $r_1 \equiv r_2 \equiv 3 \pmod 8$. Each
$r_i$ has $d_{r_i}$ odd, $d_{r_i} \mid W_{p-2}$, $d_{r_i} > 43$, and
every prime factor of $d_{r_i}$ in $A$. If both $r_i > 10^{12}$, then
$\{r_1, r_2\} \subseteq T_p^{>}$ and PPPC (Conjecture~\ref{conj:pppc})
fails at $p$. Otherwise some $r_i \leq 10^{12}$ and $p$ is a witness
exponent of the Platinum enumeration (\S\ref{sec:x1}).
\item[(b)] \emph{Mixed pair:} $r_1 \equiv 1$, $r_2 \equiv 3 \pmod 8$.
Then $(d_{r_1}, r_1, p)$ is a compatible triple and $r_2$ carries datum
(ii) of the Mixed Pair Conjecture: MPC (Conjecture~\ref{conj:mpc})
fails at $p$.
\item[(c)] \emph{Split pair:} $r_1 \equiv r_2 \equiv 1 \pmod 8$. Then
$(d_{r_1}, r_1, p)$ and $(d_{r_2}, r_2, p)$ are two distinct compatible
triples; in particular NCT for general $d$
(Conjecture~\ref{conj:nct-general}) fails twice at $p$.
\end{itemize}
The witness-exponent identification in case (a) takes the Platinum
enumeration (Theorem~\ref{thm:platinum}) as a stated computational
input; cases (b) and (c) carry none.
\end{proposition}

\begin{proof}
By Lemma~\ref{lem:pp-containment} both $r_i$ divide $W_p$; being
distinct primes, they force $W_p$ composite, and each is a proper
divisor. Lemma~\ref{lem:mod8} restricts each to $r \equiv 1$ or
$3 \pmod 8$.

\emph{Inert members.} Let $r \in P_p$ with $r \equiv 3 \pmod 8$. By
Lemma~\ref{lem:v2-class}, $v_2(\ord_r(\omega_3)) = 2$, so
$d_r = \ord_r(\omega_3)/4$ is odd, and the local pass is equivalent to
$d_r \mid W_{p-2}$ (\S\ref{sec:gap-structure}). Since $W_p$ is composite
and Condition~(II) holds at $r$, Pell-sequence exclusion
(Theorem~\ref{thm:pell-exclusion}) gives $d_r > 43$ and the parity
obstruction (Theorem~\ref{thm:parity-obstruction}) forces every prime
factor of $d_r$ into $A$. If moreover $r > 10^{12}$, then $r$ satisfies
all five defining conditions of $T_p^{>}$. If instead $r \leq 10^{12}$,
then $p \leq (r-1)/2 < 5 \times 10^{11}$ (Lemma~\ref{lem:ord2}), so
$(p, r)$ is a locally passing row of the Platinum enumeration at
composite $W_p$, and $p$ is one of the three witness exponents
(\S\ref{sec:x1}).

\emph{Split members.} Let $r \in P_p$ with $r \equiv 1 \pmod 8$. $W_p$
is composite and Condition~(II) holds at $r$, so
Proposition~\ref{prop:rank-equiv} produces a compatible triple
$(d_r, r, p)$.

The three cases enumerate the possible residue patterns of
$\{r_1, r_2\}$; the dichotomy inside case (a) is exhaustive on its own,
with the Platinum Lemma (Theorem~\ref{thm:platinum}) supplying the
witness-exponent identification and showing that at most one $r_i$ is
$\leq 10^{12}$.
\end{proof}

\begin{remark}[Converse classification]\label{rem:branch-converse}
The correspondences are exact, not one-way. A compatible triple
$(d, r, p)$ has $r \in P_p$: by the rank--order correspondence
(the unconditional paragraph in the proof of Theorem~\ref{thm:quartic}),
$\ord_r(\alpha) = 8d$, hence $\ord_r(\omega_3) = 4d$ with $d$ odd, and
$d \mid W_{p-2}$ makes $\omega_3^{2 W_{p-2}} \equiv -1 \pmod r$, i.e.\
$r$ passes locally. Likewise an inert prime $r \equiv 3 \pmod 8$ with
$\ord_r(2) = 2p$ and $d_r \mid W_{p-2}$ lies in $P_p$. Consequently
$T_p^{>} \subseteq P_p$ for every $p$, and SPC implies PPPC and MPC
outright.
\end{remark}

\subsection{The working chain}

\begin{theorem}[Refined reduction chain]\label{thm:sharpened-chain}
Assume
\begin{itemize}
\item PPPC (Conjecture~\ref{conj:pppc});
\item MPC$^{>}$ (Conjecture~\ref{conj:mpc} with datum (ii) restricted to $r_{\mathrm{in}} > 10^{12}$, \S\ref{sec:mpc});
\item X1 (Conjecture~\ref{conj:x1}).
\end{itemize}
Then Conjecture~\ref{conj:suff} holds. The implication uses, as stated computational inputs, the Platinum enumeration (Theorem~\ref{thm:platinum}) and the discharge verifications of Proposition~\ref{prop:witness-discharges}. Moreover, the conjunction $\{$PPPC, MPC$^{>}$, X1$\}$ is implied by SPC (Conjecture~\ref{conj:spc}), given the same records.
\end{theorem}

\begin{proof}
Let $W_p$ be composite and pass Condition~(II). By
Lemma~\ref{lem:restriction} every prime factor of $W_p$ lies in $P_p$,
and by Lemma~\ref{lem:w1-empty} there are at least two distinct prime
factors; every factor is $\equiv 1$ or $3 \pmod 8$
(Corollary~\ref{cor:no57}).

\emph{Case A: some factor $r_s \equiv 1 \pmod 8$.} By
Lemma~\ref{lem:mod8} the inert multiplicity of $W_p$ is odd, so an inert
prime factor $r_{\mathrm{in}}$ exists, and $r_{\mathrm{in}} \in P_p$.
By Proposition~\ref{prop:branch-split}(b) the pair
$(r_s, r_{\mathrm{in}})$ violates MPC. If
$r_{\mathrm{in}} > 10^{12}$, MPC$^{>}$ is violated --- contradiction. If
$r_{\mathrm{in}} \leq 10^{12}$, then $p$ is a witness exponent
(Proposition~\ref{prop:branch-split}(a), inert-member analysis): for
$p \in \{85{,}684{,}865{,}933,\ 1{,}251{,}488{,}009\}$,
Proposition~\ref{prop:witness-discharges} exhibits a prime factor of
$W_p$ at which Condition~(II) fails, contradicting the pass at every
factor; for $p = 10{,}916{,}765{,}939$, X1 supplies the failing factor.

\emph{Case B: every factor inert.} Pick distinct prime factors
$r_1 \neq r_2$, both in $P_p$. If both exceed $10^{12}$, then
$\{r_1, r_2\} \subseteq T_p^{>}$ by
Proposition~\ref{prop:branch-split}(a) and PPPC is violated ---
contradiction. Otherwise some factor is $\leq 10^{12}$, $p$ is a witness
exponent, and the discharges and X1 close the three exponents exactly as
in Case A.

\emph{Moreover.} Assume SPC. PPPC follows since
$T_p^{>} \subseteq P_p$ (Remark~\ref{rem:branch-converse}), and MPC
since a violating pair $(r_s, r_{\mathrm{in}})$ has two distinct members
of $P_p$ (Remark~\ref{rem:branch-converse} again); MPC implies
MPC$^{>}$. For X1: $W_{10916765939}$ is composite, with the recorded
proper divisor $r = 152{,}834{,}723{,}147$ (\S\ref{sec:secondary}; the
divisibility is verifiable by a single modular exponentiation). If
Condition~(II) held at every prime factor, then every prime factor
would lie in $P_p$ ($\ord$ pinned by Lemma~\ref{lem:ord2}, the local
pass given), while Lemma~\ref{lem:w1-empty} supplies two distinct prime
factors --- so $|P_p| \geq 2$, contradicting SPC. Hence some factor
fails, which is X1.
\end{proof}

\begin{remark}[Status]\label{rem:sharpened-status}
Theorem~\ref{thm:sharpened-chain} is the working form of the reduction:
PPPC and MPC$^{>}$ are fragments of SPC
(Remark~\ref{rem:branch-converse}) and pair statements of the same
shape, and X1 --- a consequence of SPC given the recorded witness
divisor --- concerns one explicit integer, discharged by exhibiting any
single failing factor of $W_{10916765939}$. None of the three inputs is proved
here; the chain is a conditional reduction. Unlike
Theorem~\ref{thm:spc-reduction}, it carries the two stated computational
inputs --- that is the price of trading the split--split fragment and
the sub-cutoff legs for certified finite records.
\end{remark}

\subsection{Counting and erosion variants}\label{sec:variants}

The proof of Theorem~\ref{thm:sharpened-chain} is an all-or-nothing
argument, but it also yields a quantitative form, an input-free form,
and a formally weaker triple-count form.

\begin{remark}[Counting forms and partial progress]\label{rem:counting-form}
Write $N_{\mathrm{in}}$ for the number of primes $p \geq 5$ such that $W_p$ is
composite, every prime factor of $W_p$ is inert, and $W_p$ passes
Condition~(II) --- the all-inert branch of Conjecture~\ref{conj:suff}.

(i) \emph{Quantitative.} Given the Platinum enumeration and the
discharges, Case B of the proof shows that every $p$ counted by $N_{\mathrm{in}}$
either contributes two distinct members to $T_p^{>}$ --- hence at least
$1$ to the pair count $\Pi$ of \S\ref{sec:equiv-forms}(a) --- or equals
the X1 exponent $10{,}916{,}765{,}939$. Therefore
\[
  N_{\mathrm{in}} \;\leq\; \Pi + 1,
\]
and any finite bound on $\Pi$ bounds the all-inert branch: partial
progress on PPPC transfers quantitatively.

(ii) \emph{Input-free.} With no computational input at all, the same
argument gives: every $p$ counted by $N_{\mathrm{in}}$ admits two \emph{distinct}
inert primes $r_1 \neq r_2$, both passing Condition~(II) locally, of
unrestricted size (Lemmas~\ref{lem:restriction} and~\ref{lem:w1-empty})
--- a violation of the cutoff-free inert-pair fragment of SPC. The
Platinum enumeration's only role in Theorem~\ref{thm:sharpened-chain}
is to push this pair above the cutoff $10^{12}$ --- converting the
fragment into PPPC --- at the price of the three witness exponents.
\end{remark}

\begin{corollary}[Erosion to a triple count]\label{cor:triple-erosion}
Assume MPC$^{>}$, X1, and the two statements
\begin{itemize}
\item[(T)] $|T_p^{>}| \leq 2$ for every prime $p \geq 5$ --- the triple
form of PPPC: no exponent carries three distinct above-cutoff passers;
\item[(C)] no composite $W_p$ with every prime factor inert and passing
Condition~(II) at every factor has exactly two distinct prime factors,
both $> 10^{12}$ --- the \emph{Wieferich-paired corner}.
\end{itemize}
Then Conjecture~\ref{conj:suff} holds, with the same two stated
computational inputs as Theorem~\ref{thm:sharpened-chain} and no new
ones. Both (T) and (C) are implied by PPPC, so the conjunction
$\{$MPC$^{>}$, (T), (C), X1$\}$ is formally weaker than that of
Theorem~\ref{thm:sharpened-chain}.
\end{corollary}

\begin{proof}
Case A of Theorem~\ref{thm:sharpened-chain} uses only MPC$^{>}$, the
Platinum enumeration, the discharges and X1, and closes unchanged. In
Case B every prime factor of the composite $W_p$ is inert and lies in
$P_p$. If some factor is $\leq 10^{12}$, then $p$ is a witness exponent
(Proposition~\ref{prop:branch-split}(a)) and the discharges and X1
close the three exponents as before. Otherwise every distinct prime
factor exceeds $10^{12}$ and lies in $T_p^{>}$
(Proposition~\ref{prop:branch-split}(a) again). By
Lemma~\ref{lem:w1-empty} there are at least two distinct factors: at
least three would give $|T_p^{>}| \geq 3$, contradicting (T); exactly
two is the configuration excluded by (C). For the implications from
PPPC: $|T_p^{>}| \leq 1$ gives (T) outright, and a corner instance
exhibits two distinct members of $T_p^{>}$, contradicting PPPC.
\end{proof}

\begin{remark}[The corner is doubly conditioned; the erosion is terminal]\label{rem:erosion-terminal}
A corner instance has exactly two distinct inert prime factors and
total multiplicity odd (Lemma~\ref{lem:mod8}), hence $\geq 3$, so some
factor is repeated: by Corollary~\ref{cor:multiplicity-wieferich} that
factor is a base-$2$ Wieferich prime $\equiv 3 \pmod 8$, necessarily
$> 2^{64}$ (\S\ref{sec:wieferich}), coexisting with a second locally passing factor $> 10^{12}$
at the same exponent --- a doubly-conditioned configuration.
Heuristically the erosion replaces a squared coincidence by a cubed one
plus this corner; empirically $|T_p^{>}| \leq 1$ wherever computed
(\S\ref{sec:empirical}), so the pair, triple and corner counts all
vanish as far as computed. The step does not iterate usefully: at the
next level a squarefree stratum appears --- three distinct large
factors, all of multiplicity one --- that carries no rarity tag: the
parity dividend is spent, squarefreeness means precisely that every
factor is base-$2$ non-Wieferich (Lemma~\ref{lem:w1-valuation}), and no
factor lies in the Platinum range; the function-field calibration of
\S\ref{sec:weaker-blocked} realizes exactly this configuration, so no
technique transferable to $\F_q[t]$ can empty it. Structurally, (T)
meets the obstruction of Observation~\ref{obs:additive-multiplicative}
unchanged: the erosion shrinks the footprint of the residual, not the
wall.
\end{remark}

\begin{remark}[Legacy forms of the chain]\label{rem:legacy-chains}
Earlier versions of this paper phrased the reduction through NCT for
general $d$ and a blanket Wieferich-type exclusion, and through a
second-moment counting bound with a bucket taxonomy for the
non-squarefree configurations. Both are subsumed: NCT-general implies
MPC (Proposition~\ref{prop:nct-implies-mpc}), so
$\{$PPPC, NCT-general, X1$\}$ also closes the chain; and the two-distinct-
factor input that the counting machinery laboured to extract from
squarefreeness now comes unconditionally from
Lemma~\ref{lem:w1-empty}, making the multiplicity buckets unnecessary.
The Wieferich constraints on repeated factors retain independent
interest and are collected in \S\ref{sec:wieferich}.
\end{remark}

\section{Open problems}\label{sec:open-problems}

The full sufficiency (Conjecture~\ref{conj:suff}) remains open, and so does the Single-Pass Conjecture that implies it (Theorem~\ref{thm:spc-reduction}). The unconditional gap is measured by two problems --- NCT for general $d$ (\S\ref{sec:nct-problem}), which subsumes the split--split fragment of SPC and implies the mixed fragment MPC (Proposition~\ref{prop:nct-implies-mpc}), and PPPC (\S\ref{sec:pppc-problem}), the above-cutoff inert-pair fragment and the principal residual. The working chain of Theorem~\ref{thm:sharpened-chain} proves Conjecture~\ref{conj:suff} from PPPC together with the formally weaker MPC$^{>}$ (\S\ref{sec:mpc}) and the single-exponent statement X1 (Conjecture~\ref{conj:x1}); but \emph{proving} MPC meets the same structural obstruction as NCT --- Observation~\ref{obs:additive-multiplicative} applies unchanged to its split datum --- so the discussion below is phrased for the two principal residuals. Multiplicity is no longer a hypothesis anywhere in the chain: the two-distinct-factor input is unconditional (Lemma~\ref{lem:w1-empty}), and a repeated prime factor of a passing $W_p$ would in any case be a base-$2$ Wieferich prime above $2^{64}$ (\S\ref{sec:wieferich}, \cite{ref:OEISWieferich}) --- a constraint of independent interest, not discussed separately below.

The two principal problems are \emph{technically} independent --- they live on different branches ($r \equiv 1$ versus $r \equiv 3 \pmod 8$), in different algebras ($\F_r$ versus $\F_{r^2}$), and admit different attacks. We record at the outset that they nonetheless meet a single \emph{structural} obstruction, which names precisely what new arithmetic input a closure of either would require. This is a unification of the obstruction, not a reduction of one problem to the other: the split side is delimited by the $\mu_p$-tautology of \S\ref{sec:nct-problem} and the finiteness band of Proposition~\ref{prop:nct-band}; the inert side is delimited by Proposition~\ref{prop:frobenius-blind} and the pair-count partition of Proposition~\ref{prop:refined-second-moment}. Neither side implies the other.

\begin{observation}[The additive-coupling / multiplicative-projection obstruction]\label{obs:additive-multiplicative}
Both problems demand forcing a coincidence between two multiplicative orders attached to a single prime factor $r$ of $W_p$: the order of the \emph{non-unit} $2 = (\sqrt 2)^2$ and the order of a \emph{unit} $u$ of $\Z[\sqrt 2]^\times$ --- $u = \omega_3 = \alpha^2$ (norm $+1$) on the inert branch, $u = \alpha = 1 + \sqrt 2$ (norm $-1$) on the split branch. The obstruction has two levels, not to be conflated.

\emph{(1) Global level.} In $\Z[\sqrt 2]^\times$ the cyclic subgroups $\langle 2 \rangle$ and $\langle u \rangle$ meet trivially, $\langle 2 \rangle \cap \langle u \rangle = \{1\}$, by the valuation $v_{(\sqrt 2)}$ at the ramified prime: $v_{(\sqrt 2)}(2) = 2$ while every unit has $v_{(\sqrt 2)} = 0$, so $2^m = u^n$ forces $m = 0$ and then $n = 0$. Hence there is \emph{no} multiplicative relation between $2$ and $u$; the only relation is \emph{additive} --- $\sqrt 2 = \alpha - 1$, whence $2 = (\alpha - 1)^2$ (inert) and $2 = \alpha - \alpha^{-1}$ (split).

\emph{(2) Residue level.} The datum deciding membership lives in a multiplicative power-residue projection of $\F_r^\times$ or $\F_{r^2}^\times$, and the one available relation --- the additive $\sqrt 2 = \alpha - 1$ --- has \emph{no image} there: $x \mapsto x^{(r-1)/p}$ does not factor through addition, since $(x + y)^{(r-1)/p}$ is not a function of $x^{(r-1)/p}$ and $y^{(r-1)/p}$ (e.g.\ in $\F_{73}$ with $p = 3$, the elements $1$ and $3$ share the projection value $1^{24} = 3^{24} = 1$, yet $(1+1)^{24} = 64 \neq 8 = (3+1)^{24}$). At this level ``multiplicative independence'' is \emph{false} and must not be invoked: in the cyclic group $\F_r^\times$ the subgroups $\langle \alpha \rangle$ and $\langle 2 \rangle$ meet in a subgroup of order $\gcd(\ord_r \alpha, \ord_r 2)$, generically large, and equal to $\{\pm 1\}$ --- never $\{1\}$ --- on the split branch where $\gcd(8 d, 2 p) = 2$. The two branches realize the obstruction differently:
\begin{itemize}
\item \emph{Split, $r \equiv 1 \pmod 8$.} Both $2$ and $\alpha$ lie in $\F_r$; the decisive datum is $j_\alpha = \log_2(\alpha) \bmod p$, read in the $\mu_p$-coordinate $x \mapsto x^{(r-1)/p}$. The three exact $\F_r$-internal identities (i)--(iii) of \S\ref{sec:nct-problem} all descend from $\sqrt 2 = \alpha - 1$ and collapse under $x^{(r-1)/p}$ to $0 \equiv 0 \pmod p$ (the $\mu_p$-tautology), constraining $j_\alpha$ not at all.
\item \emph{Inert, $r \equiv 3 \pmod 8$.} Here $2 \in \F_r$ but $\omega_3 \in \F_{r^2}$ on the norm-$1$ line $\mu_{r+1}$; $r$ must simultaneously divide $2^p + 1$ with $\ord_r(2) = 2 p$ (an $r-1$ datum on the non-unit $2$) and be a primitive divisor of the fixed $V_{2 d_r}$ with $\ord_r(\omega_3) = 4 d_r$, $d_r \mid W_{p-2}$ (an $r+1$ datum on the unit). These sit on opposite sides of $\gcd(r-1, r+1) = 2$; there is no single field element whose $\mu_p$-coordinate decides membership, and no analogue of the tautological trio --- the realization here is the integer order-coincidence itself.
\end{itemize}
Consequently every coupled method surveyed in \S\S\ref{sec:nct-problem}--\ref{sec:pppc-problem} meets the same wall, in three faces, none a claim of mathematical impossibility, only of the reach of presently-known methods:
\begin{itemize}
\item[(R)] reciprocity / Stickelberger / class field theory \emph{relocate} the power-residue symbol over the Galois orbit of $r$ (symmetrising it into a norm) but cannot \emph{create} its value;
\item[(C)] fixed-conductor Galois / Iwasawa / Chebotarev / $\mathrm{abc}$-radical invariants are \emph{blind} to the order datum, because $\ord_r(2)$ is integer data --- ``a difference of infinitely many splitting conditions, not the splitting condition of any one finite extension'' (Proposition~\ref{prop:frobenius-blind}, unconditional);
\item[(A)] additive, resultant and elimination identities --- and the Bugeaud--Corvaja--Zannier gcd bounds for the two recurrences $2^p + 1$ and $V_{2 d}$ --- \emph{die}, because the projection does not factor through addition (and, for the gcd bounds, because they constrain \emph{size} in families-in-$n$, never the \emph{count} at a fixed value).
\end{itemize}
On the inert (PPPC) side a fourth wall, analytic rather than algebraic, compounds these three. The residual there is a \emph{count} --- the pair-sum $\Pi$ --- and the condition carrying its finiteness, ``$r$ is a primitive divisor of the \emph{fixed integer} $V_{2 d}$'', is not a prime-in-arithmetic-progression condition; sieve and effective-Chebotarev bounds count only the latter, so they do not by themselves bound $\Pi$. This is the counting analogue of face~(C), and like (C) it is unconditional as a limitation of fixed-Frobenius data; see \S\ref{sec:weaker-blocked}.

A method closing \emph{either} problem must supply genuinely new arithmetic content living in the multiplicative power-residue projection itself --- an object of odd order $p$ (split) or a determinacy law for the multiset of orders of the divisors of one fixed integer (inert) --- coupling $2$ and the unit \emph{across} the additive bridge $\sqrt 2 = \alpha - 1$, rather than scoring $r$ against either structure alone. No such input is presently known. The asymmetry is real and is the only difference: PPPC is inert ($\F_{r^2}$, two prime factors per composite $W_p$ available, residual a pair-count $\Pi$); NCT is split ($\F_r$, one factor per $W_p$, so two-prime contradictions are structurally unavailable, and every class-field-theoretic angle yields an equivalence, never a forcing, endpoint ``$a \equiv 1 \pmod 8$''). The decoupling integers $\gcd(r-1, r+1) = 2$ (inert) and $\gcd(8 d, 2 p) = 2$ (split) are the same ``$=2$'' obstruction in two guises.
\end{observation}


\medskip\noindent The computational data and certificates underlying the finite ranges quoted above and in Table~\ref{tab:status} are collected in Appendix~\ref{sec:computational}; the residual-obstruction facts needed to delimit the two open problems are collected in Appendix~\ref{sec:residual-appendix}. In brief:

\textbf{Problem 1 (NCT, general $d$):} prove Conjecture~\ref{conj:nct-general}, no compatible triple for any odd $d > 1$ (proved unconditionally for every odd $d \leq 200$ --- Theorem~\ref{thm:nct99} for $d \leq 99$ and complete fixed-$d$ closures at all nine parity-unblocked values $107 \leq d \leq 179$ --- and at all nineteen parity-unblocked values of $(200,400]$, \S\ref{sec:nct-verify}). The split realization of Observation~\ref{obs:additive-multiplicative} confines each $d$ to an effective finite search (Proposition~\ref{prop:nct-band}) but leaves a $\mu_p$-coordinate that the known $\F_r$-internal identities do not constrain.

\textbf{Problem 2 (PPPC):} prove $|T_p^{>}| \leq 1$ for every prime $p$, i.e.\ $\Pi = 0$ (Conjecture~\ref{conj:pppc}), which in Theorem~\ref{thm:sharpened-chain} closes the all-inert branch of Conjecture~\ref{conj:suff} --- the sub-cutoff leg is certified by the Platinum enumeration, the discharges, and X1. The inert realization of Observation~\ref{obs:additive-multiplicative}, the unconditional Proposition~\ref{prop:frobenius-blind}, and the reason prime-in-AP estimates alone do not control $\Pi$ are in \S\ref{sec:pppc-problem}. The unconditional perimeter of the problem is tightened in this version: two members of $T_p^{>}$ are separated at scale $8 p \gcd(d_{r_1}, d_{r_2})$ (Theorem~\ref{thm:pair-separation}), and the diagonal part of $\Pi$ vanishes at every admissible $d \leq 400$ (Propositions~\ref{prop:diagonal-closures} and~\ref{prop:danger-census-400}), so a violation of PPPC is either off-diagonal --- congruence-consistent but blocked at the realizability level, Remark~\ref{rem:no-congruence-kill} --- or diagonal at some admissible $d > 400$, where the residue theory is complete and the remaining content is the anti-aligned corner question of \S\ref{sec:pppc-directions}.

\subsection{Generalization}\label{sec:generalization}

\begin{remarknonum}
The structural ingredients used in Sections~\ref{sec:split}--\ref{sec:inert} --- the $v_2$ classification of $\ord_r(\omega_3)$ (Lemma~\ref{lem:v2-class}), the Pell exclusion numbers $N_d = V_d^2 + 2$ (Proposition~\ref{prop:pell-closed-form}), and the compatible-triple formulation of Conjecture R3 (Proposition~\ref{prop:rank-equiv}) --- are properties of $\omega_3 \in \Z[\sqrt 2]$ and the Pell sequences. Wagstaff-specific content enters through Lemma~\ref{lem:ord2} ($\ord_r(2) = 2 p$) and the identity $W_p + 1 = 4 W_{p-2}$; both have analogues in the generalised Chua family $N_b = (b^p + 1)/(b + 1)$. For each base $b$, the Chebyshev element $\omega_b = b + \sqrt{b^2 - 1}$ plays the role of $\omega_3$, and $N_b + 1$ factors similarly. The extension requires $D = b^2 - 1$ to be non-square, so Pell sequences are replaced by the Lucas sequences of $\Z[\sqrt D]$. SPC as stated transfers verbatim to the generalised bases: the set $P_p$ is defined by the same two conditions.
\end{remarknonum}


\appendix

\section{Residual problems and obstruction summary}\label{sec:residual-appendix}

This appendix records only the residual facts needed to make the conditional reduction precise. The longer diagnostic log of closed approaches, together with independent split-local character refinements, is better treated as supplementary material: it is useful for future work but not load-bearing for Theorem~\ref{thm:sharpened-chain}.

Every ingredient of the reduction that carries no conjectural input can be
assembled into a single statement: the complete list of unconditional
constraints that a counterexample to Conjecture~\ref{conj:suff} must
satisfy. The corollary below is proved unconditionally \emph{modulo the
named computational inputs} of Remark~\ref{rem:shape-inputs} --- the same
certified finite computations already stated as inputs elsewhere in the
paper --- and assumes none of SPC, PPPC, MPC, NCT-general, or X1.

\begin{corollary}[Shape of a counterexample]\label{cor:counterexample-shape}
Let $p \geq 5$ be prime, let $W_p$ be composite, and suppose $W_p$
satisfies Condition~(II). Then:

\begin{enumerate}
\item[(0)] \emph{(Common structure.)} Condition~(II) holds at every prime
factor of $W_p$; every prime factor is $\equiv 1$ or $3 \pmod 8$, and the
total multiplicity of factors $\equiv 3 \pmod 8$ is odd; every prime
factor $r$ satisfies $\ord_r(2) = 2p$ exactly, so $r \equiv 1 \pmod{2p}$
and $r = 2pk + 1$ for an integer $k \geq 1$, and no prime divides two
composite Wagstaff numbers; $W_p$ is not a perfect power, so $W_p$ has at
least two distinct prime factors; and any \emph{repeated} prime factor of
$W_p$ is a base-$2$ Wieferich prime $\equiv 1$ or $3 \pmod 8$ exceeding
$2^{64}$ (base-$2$ \emph{super}-Wieferich if its multiplicity is
$\geq 3$), which, if inert, additionally satisfies the
Pell--Wall--Sun--Sun-type condition
$\ord_{r^2}(\omega_3) = \ord_r(\omega_3)$.
Moreover exactly one of the configurations (a), (b) occurs.

\item[(a)] \emph{(Split branch: some prime factor
$r_s \equiv 1 \pmod 8$.)} Then:
  \begin{enumerate}
  \item[(a1)] a compatible triple $(d, r_s, p)$ exists: $d$ is odd,
  $d \mid W_{p-2}$, and $r_s$ is a primitive divisor of $U_{4d}$ (Pell
  rank $\rho(r_s) = 4d$), hence a divisor of the primitive part
  $\Psi_{4d}$;
  \item[(a2)] every prime factor of $d$ is Class III, and $d > 400$;
  \item[(a3)] $8 p d + 1 \leq r_s \leq (2\sqrt 2)^{2\varphi(d)}$; in
  particular $8 p d \mid r_s - 1$, so $r_s > 3200\,p$, and
  $p < (2\sqrt 2)^{2\varphi(d)}/(8 d)$;
  \item[(a4)] $2$ is a fourth-power residue modulo $r_s$, and a $q$-th
  power residue modulo $r_s$ for every prime $q \mid d$;
  \item[(a5)] $W_p$ also has at least one inert prime factor
  $r_{\mathrm{in}} \equiv 3 \pmod 8$, and every inert prime factor
  satisfies the per-factor constraints (b1)--(b3) below; furthermore,
  either every inert prime factor exceeds $10^{12}$ --- so that $p$
  violates MPC$^{>}$ --- or $p = 10{,}916{,}765{,}939$ and the
  conclusions of (b4) hold verbatim.
  \end{enumerate}

\item[(b)] \emph{(Inert branch: every prime factor
$\equiv 3 \pmod 8$.)} Then $W_p$ has $t \geq 3$ prime factors counted
with multiplicity, of which at least two are distinct, and:
  \begin{enumerate}
  \item[(b1)] every prime factor $r$ has odd order parameter
  $d_r = \ord_r(\omega_3)/4$ with
  $d_r \mid \gcd\!\big((r+1)/4,\, W_{p-2}\big)$, $d_r > 400$
  (Corollary~\ref{cor:inert-floor-400}), and
  $d_r > \log(2p)/\log(3 + 2\sqrt 2)$;
  \item[(b2)] every prime factor of every $d_r$ is Class III;
  \item[(b3)] every prime factor $r$ is a primitive divisor of the
  Pell--Lucas number $V_{2 d_r} = V_{d_r}^2 + 2$ in the $d$-indexed
  family, and its complementary cofactor satisfies
  $W_p / r \equiv 1 \pmod{4 d_r}$;
  \item[(b4)] at most one prime factor is $\leq 10^{12}$; if one exists,
  then $p = 10{,}916{,}765{,}939$, that factor is the Platinum witness
  $r = 152{,}834{,}723{,}147 = 14 p + 1$ (with
  $d_r = 38{,}208{,}680{,}787$), and every other prime factor of $W_p$
  exceeds $2 p \cdot 10^{12} > 2.1 \times 10^{22}$;
  \item[(b5)] if $p \neq 10{,}916{,}765{,}939$, then every prime factor
  exceeds $10^{12}$, and $W_p$ has at least two distinct prime factors
  --- at least three if $W_p$ is squarefree --- all lying in $T_p^{>}$;
  in particular $|T_p^{>}| \geq 2$, a violation of PPPC. Any two of
  them, $r_1 = 2 p k_1 + 1 < r_2 = 2 p k_2 + 1$ with
  $d_i = d_{r_i}$, satisfy: $k_1 \equiv k_2 \equiv p \pmod 4$;
  $\gcd(p k_1 + 1,\, p k_2 + 1) \mid k_2 - k_1$;
  $8 p \gcd(d_1, d_2) \mid r_2 - r_1$; the per-candidate and
  pair-coupled valuation congruences of
  Propositions~\ref{prop:small-k-valuation}
  and~\ref{prop:pair-couplings}; and, if $d_1 = d_2 = d$, then
  $d > 400$, $4 d \mid k_2 - k_1$, hence $k_2 - k_1 \geq 1604$ and
  $r_2 - r_1 \geq 3208\, p$;
  \item[(b6)] if $W_p$ is not squarefree, its repeated factor is a
  base-$2$ Wieferich prime $\equiv 3 \pmod 8$ exceeding $2^{64}$
  satisfying the Pell--Wall--Sun--Sun-type condition of item (0), and
  --- if additionally $p \neq 10{,}916{,}765{,}939$ ---
  $W_p \geq (2^{64})^2 \cdot 10^{12}$, so $p \geq 173$.
  \end{enumerate}
\end{enumerate}
\end{corollary}

\begin{proof}
Pure bookkeeping; every item is an existing result applied to the
hypothesis.

\emph{(0).} Condition~(II) restricts along the ring surjection
$\Z[\sqrt 2]/(W_p) \to \Z[\sqrt 2]/(r)$ for every prime $r \mid W_p$
(Lemma~\ref{lem:restriction}); the residue classes and the
parity of the inert multiplicity are Lemma~\ref{lem:mod8}; the pinned
order is Lemma~\ref{lem:ord2} (equivalently
Theorem~\ref{thm:order-pinning}), and the disjointness across exponents
is Theorem~\ref{thm:multi-pinning}. $W_p$ is not a perfect power by
Lemma~\ref{lem:w1-empty}; with compositeness this gives two distinct
factors. A repeated factor $r$ has
$v_r(2^{\,r-1}-1) = v_r(W_p) \geq 2$ (Lemma~\ref{lem:w1-valuation}), so
$r$ is base-$2$ Wieferich, super-Wieferich when $v_r(W_p) \geq 3$
(Corollary~\ref{cor:multiplicity-wieferich}); its residue class is $1$
or $3 \pmod 8$ (Lemma~\ref{lem:mod8}), and the exhaustive search below
$2^{64}$ finds only $1093 \equiv 5$ and $3511 \equiv 7 \pmod 8$
\cite{ref:OEISWieferich}, so $r > 2^{64}$. For an inert repeated factor,
the global pass gives Condition~(II) modulo $r^2$ and the
Pell--Wall--Sun--Sun condition follows as in Remark~\ref{rem:wieferich}
(cf.\ Remark~\ref{rem:pell-wss-global}).

\emph{(a1).} Condition~(II) at $r_s$ in a composite $W_p$ yields
$\ord(\alpha_0) = 4d$ with $d$ odd, $d \mid W_{p-2}$ (descent of
\S\ref{sec:split}), and Proposition~\ref{prop:rank-equiv}(a) gives
$\rho(r_s) = 4d$; this is a compatible triple
(Definition~\ref{def:compat-triple}), and a primitive divisor of
$U_{4d}$ divides $\Psi_{4d}$ (as in
Corollary~\ref{cor:nct-fixed-d-certificate}).

\emph{(a2).} Each prime $q \mid d$ divides $W_{p-2}$; for $q > 3$,
Theorem~\ref{thm:exact-ap} forces $q \in A$, and $3 \in A$. NCT holds
unconditionally for every odd $d \leq 400$ (Theorem~\ref{thm:nct99} and
the twenty-eight complete fixed-$d$ certificates of
\S\ref{sec:nct-verify}), so $d > 400$.

\emph{(a3).} Proposition~\ref{prop:nct-band}; $r_s > 3200\,p$ follows
from $8 p d \mid r_s - 1$ with $d \geq 401$.

\emph{(a4).} Theorem~\ref{thm:quartic} (its hypotheses --- $r_s \equiv 1
\pmod 8$ a primitive divisor of $U_{4d}$, $d$ odd --- hold by (a1)) and
Proposition~\ref{prop:qth-power}.

\emph{(a5).} By Lemma~\ref{lem:mod8} the inert multiplicity is odd,
hence $\geq 1$. Every inert factor passes Condition~(II) locally (item
(0)), so the descent of \S\ref{sec:gap-structure} gives (b1)'s
divisibilities; $d_r > 400$ by Corollary~\ref{cor:inert-floor-400}
(applicable since $W_p$ is composite and passes at every factor), the Class
III property by Theorem~\ref{thm:parity-obstruction}, the lower bound by
Lemma~\ref{lem:lower-bound}, and (b3) by
Lemma~\ref{lem:inert-primitivity},
Proposition~\ref{prop:pell-closed-form} (with
Theorem~\ref{thm:pell-divisibility}) and
Remark~\ref{rem:cofactor-congruence} --- exactly as in
Remark~\ref{rem:mpc-data}. If every inert factor exceeds $10^{12}$, the
pair of data (i), (ii) with $r_{\mathrm{in}} > 10^{12}$ violates
MPC$^{>}$ (\S\ref{sec:mpc}), as in Case A of
Theorem~\ref{thm:sharpened-chain}. Otherwise some inert factor is
$\leq 10^{12}$: then $p \leq (r-1)/2 < 5 \times 10^{11}$
(Lemma~\ref{lem:ord2}), so $(p, r)$ is a locally passing row of the
Platinum enumeration (which ranges over inert factors of
\emph{composite} $W_p$), $p$ is one of the three witness exponents
(Theorem~\ref{thm:platinum}, \S\ref{sec:secondary}), and
Proposition~\ref{prop:witness-discharges} exhibits failing factors at
the other two --- contradicting item (0) --- so
$p = 10{,}916{,}765{,}939$, and (b4)'s conclusions follow as below.

\emph{(b).} $t \geq 3$ with every $d_i > 43$ and Class III is
Theorem~\ref{thm:multi-factor}; two distinct factors from item (0).
Items (b1)--(b3) are per-factor and were justified under (a5). (b4): at
most one factor $\leq 10^{12}$ is Theorem~\ref{thm:platinum}; if one
exists, the witness-exponent identification and the elimination of the
two discharged exponents are as in (a5), and \S\ref{sec:secondary}
records the witness factor and its $d_r$ at
$p = 10{,}916{,}765{,}939$; the directed scan at this exponent, complete
for $k \leq 10^{12}$, finds no second prime factor in the progression
$r = 2 p k + 1$ (\S\ref{sec:x1}, \S\ref{sec:inert-survey}), and every
prime factor lies in that progression (item (0)), so every other factor
has $k > 10^{12}$, i.e.\ exceeds $2 p \cdot 10^{12} = 2.18 \times
10^{22}$.

\emph{(b5).} If $p \neq 10{,}916{,}765{,}939$, no factor is
$\leq 10^{12}$ (else (b4) would force $p = 10{,}916{,}765{,}939$). Every
distinct prime factor then satisfies all five defining conditions of
$T_p^{>}$ (Conjecture~\ref{conj:pppc}) by (b1) and item (0); there are at
least two distinct factors, and at least three when $W_p$ is squarefree
(Theorem~\ref{thm:multi-factor}: $t \geq 3$ with multiplicity $=$
distinct count), exactly as in Case B of
Theorem~\ref{thm:sharpened-chain}. The pair constraints are
Theorem~\ref{thm:pair-separation}(1)--(3) (its hypotheses hold with
$d_i = d_{r_i} \mid (r_i + 1)/4$ odd),
Propositions~\ref{prop:small-k-valuation}
and~\ref{prop:pair-couplings} (condition (v) holds at both candidates).
On the diagonal $d_1 = d_2 = d$: the two factors form two same-$d$
danger triples (Definition~\ref{def:danger-triple}), so
Theorem~\ref{thm:pair-separation}(4) gives $4 d \mid k_2 - k_1$, and
Propositions~\ref{prop:diagonal-closures} and~\ref{prop:danger-census-400}
exclude every admissible
$d \leq 400$; hence $d \geq 401$ ($d$ odd), $k_2 - k_1 \geq 4 d \geq
1604$ and $r_2 - r_1 = 2p(k_2 - k_1) \geq 3208\,p$ (cf.\
Corollary~\ref{cor:diagonal-k500}).

\emph{(b6).} Item (0) restricted to the all-inert case (every factor
$\equiv 3 \pmod 8$); under (b5) the repeated factor exceeds $2^{64}$ and
a second distinct factor exceeds $10^{12}$, and the repeated factor has
multiplicity $\geq 2$, so $W_p \geq (2^{64})^2 \cdot 10^{12}$, giving
$2^p \geq 3 W_p - 1 > 2^{167}$ and $p \geq 168$, hence $p \geq 173$ ($p$
prime).
\end{proof}

\begin{remark}[Computational inputs of the corollary]\label{rem:shape-inputs}
Corollary~\ref{cor:counterexample-shape} is unconditional modulo exactly
the following certified finite computations, each already a stated input
elsewhere in the paper; no conjecture enters.
\begin{itemize}
\item \emph{Platinum enumeration} (Theorem~\ref{thm:platinum}): the
$684{,}965{,}381$-row enumeration of all inert prime factors
$r \leq 10^{12}$ ($p < 5 \times 10^{11}$), used in (a5), (b4), (b5).
\item \emph{Witness discharges}
(Proposition~\ref{prop:witness-discharges}): the certified failing
factors at $p = 85{,}684{,}865{,}933$ and $p = 1{,}251{,}488{,}009$,
used in (a5), (b4), (b5).
\item \emph{Directed scan at $p = 10{,}916{,}765{,}939$}
(\S\ref{sec:x1}, \S\ref{sec:inert-survey}): completeness of the
second-factor search for $k \leq 10^{12}$, used in (b4).
\item \emph{Base-$2$ Wieferich search below $2^{64}$}
\cite{ref:OEISWieferich}: used in item (0) and (b6).
\item \emph{NCT certificates for $d \leq 400$}
(Theorem~\ref{thm:nct99} and \S\ref{sec:nct-verify}): complete
primitive-part factorizations with every prime factor APR-CL-certified,
used in (a2).
\item \emph{Danger census through $d = 400$}
(Proposition~\ref{prop:danger-census-400} and
Corollary~\ref{cor:inert-floor-400}): the complete danger-triple census
over the $34$ admissible rungs $43 < d \leq 400$ --- complete $V_{2d}$
factorizations, every listed prime APR-CL-certified, including the
secondary-factor verification at the $d = 67$ triple --- used, through
the floor $d_r > 400$, in (a5) and (b1), and in the diagonal exclusion
of (b5).
\item \emph{Diagonal closures} (Proposition~\ref{prop:diagonal-closures}):
the complete APR-CL-certified factorizations of $V_{114}$, $V_{118}$,
$V_{134}$, $V_{162}$, $V_{166}$, $V_{198}$, $V_{214}$, $V_{242}$, used
in (b5).
\end{itemize}
Items (a1), (a4), (b2), (b3) and the pair-separation congruences of
(b5) carry \emph{no} computational input at all: they are fully
algebraic, as is the perfect-power exclusion of item (0), which rests on
the published Nagell--Ljunggren results \cite{ref:BennettLevin,
ref:BugeaudMignotte2002, ref:BugeaudCaoMignotte2001}. The remaining
items are algebraic apart from one clause each: item (0) apart from the
Wieferich bound $r > 2^{64}$, (b1) apart from the census floor
$d_r > 400$, and (a3) apart from the numerical instantiation
$r_s > 3200\,p$, which inherits the NCT floor $d > 400$ of (a2).
\end{remark}

Two readings of Corollary~\ref{cor:counterexample-shape} deserve
emphasis. First, the \emph{quotable contrapositive}: any composite
Condition-(II) passer either contains a split factor --- and then a
compatible triple at some admissible $d > 400$, where each further value
of $d$ is a finite certificate whose cost is a complete factorization of
$\Psi_{4d} \approx (2\sqrt 2)^{2\varphi(d)}$ (the certified range ends,
at present, with the $290$-digit $\Psi_{1516}$ at $d = 379$,
\S\ref{sec:nct-verify}) --- or it is all-inert, and then either
$p = 10{,}916{,}765{,}939$ (the single exponent that X1 addresses) or
$W_p$ realizes a violation of PPPC outright. Second, the \emph{scale}:
the size constraints alone force $W_p \geq 10^{36}$, i.e.\ $p \geq 127$,
in the all-inert branch ($p \geq 173$ in its non-squarefree corner), but
size is not the content --- the content is a coincidence of a type never
yet observed. The $684{,}965{,}381$-row Platinum enumeration contains
exactly three locally passing rows, at three pairwise distinct
exponents; the above-cutoff scans of \S\ref{sec:empirical} tested
$376{,}435$ candidates satisfying conditions (i)--(iv) and found
\emph{zero} passing condition (v); and the $9{,}506{,}750$-rectangle
corner census of Remark~\ref{rem:corner-sweep} found no danger event at
any composite $W_p$. A counterexample in the main branch (b5) requires
two condition-(v) passes above the cutoff \emph{at the same exponent} ---
two events of a kind not observed even once, individually, anywhere in
the computed record --- subject, in addition, to the separation
congruences of (b5). The corollary thus quantifies, without any
conjectural input, how far outside every computed window a
counterexample must live.

\subsection{Problem 1: NCT for general \texorpdfstring{$d$}{d}}\label{sec:nct-problem}

Conjecture~\ref{conj:nct-general} asks for no compatible triple $(d,r,p)$ at any odd $d>1$. Theorem~\ref{thm:nct99} proves the range $d\leq 99$ unconditionally; Appendix~\ref{sec:computational} gives complete fixed-$d$ closures for all nine parity-unblocked values of $100 \leq d \leq 200$ and for every parity-unblocked value of $(200, 400]$, so NCT holds unconditionally for every odd $d \leq 400$. What remains for $d > 400$ is not an unbounded search at a fixed $d$: each $d$ has an explicit finite band.

\begin{proposition}[NCT structural finiteness band]\label{prop:nct-band}
Let $(d, r, p)$ be a compatible triple (Definition~\ref{def:compat-triple}). Then
\[
  8 p d + 1 \leq r \leq (2\sqrt 2)^{2\varphi(d)}.
\]
Consequently $p < (2\sqrt 2)^{2\varphi(d)}/(8d)$, and for each fixed odd $d>1$ the possible pairs $(r,p)$ are finite and explicitly bounded.
\end{proposition}

\begin{proof}
By Proposition~\ref{prop:rank-equiv}, $\ord_r(\alpha)=8d$ for $\alpha=1+\sqrt2$, while compatibility gives $\ord_r(2)=2p$. Since $p\nmid d$ (Lemma~\ref{lem:pndivW}), the two cyclic subgroups of $\F_r^\times$ generated by $\alpha$ and $2$ intersect exactly in $\{\pm1\}$. Their product has order $(8d)(2p)/2=8pd$, hence $8pd\mid r-1$ and $r\geq 8pd+1$.

For the upper bound, $r$ is a primitive divisor of $U_{4d}$ and hence divides the primitive part
\[
  \Psi_{4d}=\prod_{\zeta}(\alpha-\zeta\bar\alpha),
\]
where $\zeta$ runs over primitive $4d$-th roots. Each factor has absolute value at most $\alpha+|\bar\alpha|=2\sqrt2$, so
\[
  r\leq |\Psi_{4d}|\leq (2\sqrt2)^{\varphi(4d)}=(2\sqrt2)^{2\varphi(d)}.
\]
\end{proof}

\begin{corollary}[Fixed-$d$ NCT certificate]\label{cor:nct-fixed-d-certificate}
Fix an odd $d>1$. If a complete factorization of $\Psi_{4d}$ is known, then NCT at this $d$ is checked by the following finite certificate:
\begin{enumerate}
\item retain only prime factors $r\mid\Psi_{4d}$ with $r\equiv1\pmod8$;
\item compute $p_r=\ord_r(2)/2$ and discard every $r$ for which $p_r$ is not prime;
\item for each remaining $r$, test the exact condition $d\mid W_{p_r-2}$, equivalently the valuation-adjusted congruence system modulo $L_d$ when $d$ is admissible.
\end{enumerate}
If no factor survives, no compatible triple exists at this $d$.
\end{corollary}

\begin{proof}
In a compatible triple at $d$, Proposition~\ref{prop:rank-equiv} gives Pell rank $\rho(r)=4d$, so $r$ is a primitive divisor of $U_{4d}$ and hence divides the primitive part $\Psi_{4d}$. The split residue condition is $r\equiv1\pmod8$. Order-Pinning supplies the only possible Wagstaff exponent $p_r=\ord_r(2)/2$, which must be prime, and compatibility additionally requires $d\mid W_{p_r-2}$. Remark~\ref{rem:valuation-adjusted-Ld} gives the equivalent congruence formulation for the last condition when $d$ is admissible; if $d$ is not admissible, Theorem~\ref{thm:exact-ap} already rules out the relevant divisibility at a non-Class-III prime factor.
\end{proof}

The band is effective but not uniform in $d$; the upper bound grows exponentially in $\varphi(d)$. Thus NCT-general is reduced to infinitely many finite searches, not to one finite computation. The algebraic reason the known identities stop here is the split realization of Observation~\ref{obs:additive-multiplicative}: after projecting to the $\mu_p$-coordinate of $\F_r^\times$, the identities coming from $\sqrt2=\alpha-1$ collapse to tautologies and do not determine the odd-order component that Conjecture R3 needs.

\subsection{Problem 2: PPPC}\label{sec:pppc-problem}

PPPC (Conjecture~\ref{conj:pppc}) asserts $|T_p^{>}|\leq 1$ for every prime $p$, equivalently $\Pi=0$ (\S\ref{sec:equiv-forms}). In the working chain (Theorem~\ref{thm:sharpened-chain}) PPPC closes the all-inert branch, the sub-cutoff leg being certified by the Platinum enumeration, the discharges, and X1. The useful way to inspect the residual is to partition by the order parameter.

\begin{proposition}[Refined second-moment identity]\label{prop:refined-second-moment}
With $T_p^{>}$ and $\Pi$ as in \S\ref{sec:pppc} and \S\ref{sec:equiv-forms}(a), for odd admissible $d>43$ let
\[
  P_d(p)=\#\{\,r\in T_p^{>}: d_r=d\,\}.
\]
Then $|T_p^{>}|=\sum_d P_d(p)$ and
\[
  \binom{|T_p^{>}|}{2}
  =\sum_{d_1<d_2}P_{d_1}(p)P_{d_2}(p)
   +\sum_d \binom{P_d(p)}{2}.
\]
After summing over $p$,
\[
  \Pi=\sum_{d_1<d_2}S(d_1,d_2)+\sum_d S_{\mathrm{diag}}(d),
\]
where $S(d_1,d_2)=\sum_pP_{d_1}(p)P_{d_2}(p)$ and $S_{\mathrm{diag}}(d)=\sum_p\binom{P_d(p)}{2}$.
\end{proposition}

\begin{proof}
This is the multinomial identity for $\binom{\sum_d n_d}{2}$ with $n_d=P_d(p)$. All terms are non-negative, so Tonelli's theorem justifies the rearrangement after summing over $p$.
\end{proof}

The diagonal term matters: one Pell--Lucas number $V_{2d}$ can have two distinct primitive divisors, so bounding only the off-diagonal terms does not bound $\Pi$. The diagonal is, however, exactly the part of $\Pi$ where unconditional progress is possible, because both members of a diagonal pair divide the one fixed integer $V_{2d}$. The following package develops this.

\begin{lemma}[Half-pinning]\label{lem:half-pinning}
Let $d$ be odd and let $r > 3$ be a prime with $r \mid V_{2d}$ and $\ord_r(2) = 2m$, $m$ odd (for $r \equiv 3 \pmod 8$ the parity of $m$ is automatic). Then there is a unique sign $\varepsilon = \varepsilon_r \in \{\pm 1\}$ with
\[
  V_d \;\equiv\; \varepsilon \, 2^{(m+1)/2} \pmod r;
\]
equivalently, $r$ divides exactly one of the two integers $H_{\varepsilon} = V_d - \varepsilon\, 2^{(m+1)/2}$.
\end{lemma}

\begin{proof}
$r \mid V_{2 d} = V_d^2 + 2$ (Proposition~\ref{prop:pell-closed-form}) gives $V_d^2 \equiv -2 \pmod r$. From $\ord_r(2) = 2 m$, the element $2^m$ has order $2$ in $\F_r^\times$, so $2^m \equiv -1 \pmod r$. Hence $V_d^2 \equiv 2^m \cdot 2 = 2^{m+1} \pmod r$ with $m + 1$ even, so $\bigl(V_d \, 2^{-(m+1)/2}\bigr)^2 \equiv 1$ and $V_d \equiv \pm\, 2^{(m+1)/2} \pmod r$. The sign is unique: if both signs held, then $r \mid 2^{(m+1)/2 + 1}$, impossible for odd $r$. For $r \equiv 3 \pmod 8$, $v_2(\ord_r(2)) = 1$ as in the proof of Lemma~\ref{lem:mod8}, so $m$ is odd automatically.
\end{proof}

For a fixed pair $(d, p)$ with $d$ odd and $p \geq 5$ prime, write
\[
  \mathcal{L}(d, p) \;=\; \{\, r > 3 \text{ prime} \;:\; r \mid V_{2 d},\;\; 2^p \equiv -1 \pmod r \,\}
\]
for the \emph{order-divisibility locus} at $(d, p)$. Every danger triple $(d, r, p)$ (Definition~\ref{def:danger-triple}) has $r \in \mathcal{L}(d, p)$: primitivity gives $r \mid V_{2 d}$ and the pinned order gives $2^p \equiv -1$. The locus is in general strictly larger than the set of danger factors --- it also captures non-primitive divisors ($\ord_r(\omega_3) = 4 e$ for a proper divisor $e \mid d$) and split divisors $r \equiv 1 \pmod 8$ of $V_{2 d}$. This over-capture is deliberate and harmless: the certificate below is used for \emph{exclusion}, which is a superset test, and any positives are post-filtered by the mod-$8$, primitivity and size requirements.

\begin{proposition}[Corner localization]\label{prop:corner-localization}
Let $d$ be odd, $p \geq 5$ prime, and set
\[
  G(\varepsilon, \varepsilon') \;=\; \gcd\!\left(V_d - \varepsilon\, 2^{(p+1)/2},\;\; 2 U_d - \varepsilon'\right), \qquad \varepsilon, \varepsilon' \in \{\pm 1\}.
\]
Then $\mathcal{L}(d, p)$ is exactly the set of prime divisors $> 3$ of the four integers $G(\varepsilon, \varepsilon')$, and each $r \in \mathcal{L}(d, p)$ divides $G(\varepsilon, \varepsilon')$ for exactly one sign pair $(\varepsilon_r, \varepsilon'_r)$. The four integers are computable from the residues $2^{(p+1)/2} \bmod (2 U_d \mp 1)$ --- at most four modular exponentiations and four gcds on integers of $O(d)$ digits --- with no factorization of $V_{2 d}$.
\end{proposition}

\begin{proof}
For odd $d$, the identities $V_{2 d} = V_d^2 + 2$ (Proposition~\ref{prop:pell-closed-form}) and $V_d^2 - 8 U_d^2 = -4$ (\S\ref{sec:notation}) combine to the rational halving
\[
  V_{2 d} \;=\; 8 U_d^2 - 2 \;=\; 2\, (2 U_d - 1)(2 U_d + 1).
\]
Let $r \in \mathcal{L}(d, p)$. From $2^p \equiv -1 \pmod r$: $\ord_r(2) \mid 2 p$ and $\nmid p$; since $p$ is \emph{prime}, $\ord_r(2) \in \{2, 2 p\}$, and $\ord_r(2) = 2$ forces $r \mid 3$, excluded --- this is the one point where the primality of $p$ is used. So $\ord_r(2) = 2 p$ with $p$ odd, and Lemma~\ref{lem:half-pinning} pins $V_d \equiv \varepsilon_r\, 2^{(p+1)/2} \pmod r$ for a unique $\varepsilon_r$. By the halving, the odd prime $r$ divides $(2 U_d - 1)(2 U_d + 1)$ and divides exactly one of the two factors (their difference is $2$): $2 U_d \equiv \varepsilon'_r \pmod r$ for a unique $\varepsilon'_r$. Hence $r \mid G(\varepsilon_r, \varepsilon'_r)$ and $r$ divides no other $G$. Conversely, let $r > 3$ be a prime divisor of some $G(\varepsilon, \varepsilon')$. Then $r \mid 2 U_d - \varepsilon'$, so $r \mid V_{2 d}$ by the halving, whence $V_d^2 \equiv -2 \pmod r$; and $r \mid V_d - \varepsilon\, 2^{(p+1)/2}$ squares to $V_d^2 \equiv 2^{p+1} \pmod r$. Together $2^{p+1} \equiv -2$, i.e.\ $2^p \equiv -1 \pmod r$, so $r \in \mathcal{L}(d, p)$. Finally $\gcd(A, B) = \gcd(A \bmod B, B)$ reduces each $G$ to the two residues $2^{(p+1)/2} \bmod (2 U_d \mp 1)$.
\end{proof}

\begin{remark}[Pair caps: three of the four sign classes are product-capped]\label{rem:pair-caps}
Fix $(d, p)$ as above, $d \geq 3$, and let $r_1 \neq r_2$ be two primes in $\mathcal{L}(d, p)$ with sign data $(\varepsilon_i, \varepsilon'_i)$.
\begin{itemize}
\item If $\varepsilon'_1 = \varepsilon'_2 = \varepsilon'$, then $r_1 r_2 \mid 2 U_d - \varepsilon'$, so $r_1 r_2 \leq 2 U_d + 1 \sim \sqrt{V_{2 d}/2}$ --- a cap valid for \emph{all} $p$.
\item If $\varepsilon_1 = \varepsilon_2 = \varepsilon$, then $r_1 r_2$ divides $H_\varepsilon = V_d - \varepsilon\, 2^{(p+1)/2} \neq 0$, so $r_1 r_2 \leq |H_\varepsilon| \leq V_d + 2^{(p+1)/2}$. ($H_\varepsilon \neq 0$ because $V_d/2$ is odd and $> 1$ for odd $d \geq 3$, by $(V_d/2)^2 = 2 U_d^2 - 1$, so $V_d$ is not a power of $2$.) This cap improves on the trivial bound $r_1 r_2 \leq V_{2 d}$ precisely in the range $p < 4 \log_2(1 + \sqrt 2)\, d + O(1) \approx 5.0862\, d$.
\end{itemize}
Only \emph{anti-aligned} pairs --- $\varepsilon_1 = -\varepsilon_2$ \emph{and} $\varepsilon'_1 = -\varepsilon'_2$ --- escape both product caps; pairs in the three aligned sign classes must fit inside a single explicit integer.
\end{remark}

These statements stratify and cap the diagonal; they do not by themselves empty it. What does empty it, in the certified range, is composing them with complete factorizations.

\begin{proposition}[Diagonal vanishing through $d = 124$]\label{prop:diagonal-closures}
For every admissible odd $d$ with $43 < d \leq 124$ (the complete list: $57$, $59$, $67$, $81$, $83$, $99$, $107$, $121$), no prime $p$ admits two distinct danger triples $(d, r_1, p)$, $(d, r_2, p)$ sharing the parameter $d$. Consequently $S_{\mathrm{diag}}(d) = 0$ for all $p$ at every admissible $d \leq 124$. The statement is $k$-free: it excludes diagonal pairs at these $d$ for danger factors of every size, not only in a finite $k$-window.
\end{proposition}

\begin{proof}[Verification]
The six Pell--Lucas numbers $V_{118}$, $V_{162}$, $V_{166}$, $V_{198}$, $V_{214}$, $V_{242}$ ($46$, $63$, $64$, $76$, $82$, $93$ decimal digits) are completely factored, with $4$, $11$, $5$, $12$, $5$, $11$ distinct prime factors exceeding $3$ respectively (factorizations assembled from the public factor database \cite{ref:FactorDB} and re-verified by exact product identities); the $42$ distinct primes occurring across the six are each certified prime by APR-CL (PARI/GP \texttt{isprime} flag~$2$). For each such prime $r$ the danger prerequisites of Definition~\ref{def:danger-triple} are decided exactly: the residue $r \equiv 3 \pmod 8$; primitivity $\ord_r(\omega_3) = 4 d$, computed by order descent in $\F_{r^2}$; $v_2(\ord_r(2)) = 1$; and primality of the pinned exponent $p_r = \ord_r(2)/2$. At $d \in \{83, 99, 107, 121\}$ no prime factor survives all four tests; at $d = 59$ exactly one survives ($r = 88{,}499$, pinned to $p = 44{,}249$) and at $d = 81$ exactly one ($r = 5507$, pinned to $p = 2753$). For $d = 57$ and $d = 67$ the same certificate over the complete factorizations of $V_{114}$ and $V_{134}$ (Remark~\ref{rem:known-danger}) leaves the surviving divisors $\{683,\ 171{,}369{,}731{,}867\}$ --- pinned to the \emph{distinct} exponents $11$ and $85{,}684{,}865{,}933$ --- and $\{148{,}937{,}135{,}261{,}632{,}265{,}803\}$, pinned to $24{,}822{,}855{,}876{,}938{,}710{,}967$. No two surviving divisors share a pinned exponent at any of the eight $d$ (condition (v) was not even needed), so no diagonal pair exists at any $p$.
{\sloppy
Certificate scripts (\S\ref{sec:toolchain}): \texttt{scripts/}\allowbreak\texttt{cp354\_samed\_}\allowbreak\texttt{diagonal\_}\allowbreak\texttt{certificates.py} and \texttt{scripts/}\allowbreak\texttt{cp354\_factordb\_}\allowbreak\texttt{v2d\_}\allowbreak\texttt{diagonal.py}.\par}
\end{proof}

\begin{corollary}[Diagonal vanishing through $d = 400$; no pair with $k \leq 1604$]\label{cor:diagonal-k500}
No prime $p$ admits two distinct danger triples $(d, r_1, p)$, $(d, r_2, p)$ with the same admissible $d \leq 400$: the diagonal contribution to $\Pi$ vanishes there for danger factors of every size. In particular no diagonal pair exists with both $k_i = (r_i - 1)/2p \leq 1604$, unconditionally.
\end{corollary}

\begin{proof}
The census of Proposition~\ref{prop:danger-census-400} is exhaustive on $43 < d \leq 400$, and no two of its rows share both the parameter $d$ and the pinned exponent (the two rows at $d = 57$ have exponents $11 \neq 85{,}684{,}865{,}933$). For the $k$-form: by Theorem~\ref{thm:pair-separation}(4), such a pair has $4 d \mid k_2 - k_1$ with $1 \leq k_2 - k_1 \leq 1603$, so $d \leq 400$; $d$ is admissible with $d > 43$ (Definition~\ref{def:danger-triple}), and every such $d$ is closed.
\end{proof}

\begin{remark}[Factorization-free corner census]\label{rem:corner-sweep}
Independently of the factorizations above, the corner certificate of Proposition~\ref{prop:corner-localization} was run over every admissible $d$ with $43 < d \leq 124$ and every prime $p \leq 10^6$ satisfying the condition-(v) congruence for that $d$ ($d \mid W_{p-2}$; non-admissible $d$ admit no such $p$ at all, by Theorem~\ref{thm:exact-ap}) --- $29{,}930$ pairs $(d, p)$ in total, at four modular exponentiations each. The locus $\mathcal{L}(d, p)$ is empty at every one of them except $(d, p) = (57, 11)$, where it consists of the single prime $683 = W_{11}$ --- the known \emph{phantom} triple of Remark~\ref{rem:known-danger}, at which $W_p$ is prime. Because the census tests the locus itself, it is $k$-unbounded --- it excludes danger factors of every size at the swept $(d, p)$, which $p$-indexed scans of the progression $r = 2 p k + 1$ cannot do --- and it over-captures by design. This is finite-range empirical verification within the swept rectangle, not a proof for general $p$. Script: \texttt{scripts/cp354\_diag\_corner\_gcd.py}. An extension of the census to every admissible $d \leq 1000$ and every condition-(v) prime $p \leq 10^8$ --- $9{,}506{,}750$ rectangles, about $320$ times the original sweep --- found exactly four corner hits, all of them the known phantom rows at \emph{prime} $W_p$ ($683 = W_{11}$ at $d = 171$; $2731 = W_{13}$ at $d = 683$; $43691 = W_{17}$ at $d = 331$ and $d = 993$): zero danger events at composite $W_p$ and zero diagonal pairs (script \texttt{scripts/}\allowbreak\texttt{cp358\_corner\_sweep\_d1000.py}).
\end{remark}

\begin{remark}[Same-$d$ triples force alignment]\label{rem:same-d-triples}
Three distinct primes in $\mathcal{L}(d, p)$ carry three sign pairs $(\varepsilon_i, \varepsilon'_i) \in \{\pm 1\}^2$; by pigeonhole some two share $\varepsilon$ \emph{and} some two (possibly a different two) share $\varepsilon'$. A same-$d$ \emph{triple} therefore always contains an $\varepsilon'$-aligned pair --- product-capped by $2 U_d + 1$ at every $p$ (Remark~\ref{rem:pair-caps}; the $\varepsilon$-aligned cap also applies, though it is size-binding only for $p \lesssim 5.0862\, d$) --- so the anti-aligned escape of Remark~\ref{rem:pair-caps} is closed \emph{for triples}; the pair-level anti-aligned corner question of \S\ref{sec:pppc-directions} remains open. Moreover a same-$d$ triple of danger-class factors contains a same-$d$ danger pair, so same-$d$ triples are excluded at every admissible $d \leq 400$ (all $k$; Propositions~\ref{prop:diagonal-closures} and~\ref{prop:danger-census-400}), and whenever two members have $k_i \leq 1604$ (Corollary~\ref{cor:diagonal-k500}).
\end{remark}

Off the diagonal, and on the diagonal beyond the certified range, the obstructions below delimit what presently-known methods can reach.

\begin{proposition}[Frobenius cannot see the order of $2$]\label{prop:frobenius-blind}
Let $M/\Q$ be any finite Galois extension. There exist primes $r\neq r'$, unramified in $M$, with the same Frobenius conjugacy class in $\Gal(M/\Q)$ but $\ord_r(2)\neq\ord_{r'}(2)$.
\end{proposition}

\begin{proof}
Fix a conjugacy class $\mathcal C\subset\Gal(M/\Q)$. Chebotarev gives infinitely many unramified primes with Frobenius class $\mathcal C$. For a fixed integer $n$, only finitely many primes satisfy $\ord_r(2)=n$, since each divides $2^n-1$. Hence $\ord_r(2)$ is unbounded on that infinite Chebotarev set and assumes at least two values there.
\end{proof}

Thus no invariant factoring through a fixed finite Galois extension can decide membership in $T_p^{>}$: the pinned exponent $p=\ord_r(2)/2$ is not finite-level Frobenius data. This is the precise face (C) of Observation~\ref{obs:additive-multiplicative}.

\subsection{Why the AP-only analytic fallback does not close the gap}\label{sec:weaker-blocked}

A weaker target would be $\Pi<\infty$, which (with the certified records and Lemma~\ref{lem:w1-empty}) would imply only finitely many all-inert global passes. The standard second-moment route still fails for a structural reason. For a fixed $d$, every relevant $r$ is a divisor of the fixed integer $V_{2d}$, hence
\[
  \sum_p P_d(p)\leq \omega(V_{2d}),\qquad
  S(d_1,d_2)\leq \omega(V_{2d_1})\omega(V_{2d_2}).
\]
A sieve or effective-Chebotarev estimate bounds primes in arithmetic progressions; it does not bound how the prime divisors of the fixed integer $V_{2d}$ distribute among the pinned orders $\ord_r(2)$. Feeding only the AP condition $d\mid W_{p-2}$ gives infinitely many $p$ in one valuation-adjusted congruence class and no finiteness for $\Pi$.

The correct modulus for composite $d$ is the valuation-adjusted $L_d$ of Remark~\ref{rem:valuation-adjusted-Ld}, not the squarefree shorthand $M_d$. Thus an attempted second-moment summation must control expressions of the shape
\[
  \frac{\omega(V_{2d_1})\omega(V_{2d_2})}
       {\varphi(L_{d_1})\varphi(L_{d_2})}
\]
together with the distribution of the actual primitive divisors of $V_{2d_i}$. This is a sharper and more concrete target than a raw AP count, but it is not presently supplied by Brun--Titchmarsh, effective Chebotarev, or GRH alone.

\medskip
\noindent\textbf{Finding.} \emph{The AP-only fallback does not prove the weaker bound $\Pi<\infty$. The missing input must detect that a prime divides the specific integer $V_{2d}$ while also satisfying the pinned Wagstaff order condition; a sharper prime-in-AP estimate, even under GRH, does not supply that information. A valuation-adjusted average over $L_d$ remains a legitimate open direction, not a closed argument.}
\medskip

The function-field calibration of \S\ref{sec:pppc-directions} sharpens the negative direction once more: in the Carlitz analogue ($q = 3, 5$) there exist moduli with three and with four distinct large primitive prime factors, all of multiplicity one, so the squarefree many-large-factor configurations are realized there, and no technique transferable to $\F_q[t]$ can empty the corresponding strata.

\subsection{Directions}\label{sec:pppc-directions}

The remaining plausible routes are those that keep the two divisibilities coupled across $W_p+1=4W_{p-2}$: sharper divisor bounds for $V_{2d}$ beyond Stewart--Tijdeman; an average distribution theorem for prime divisors of the non-standard sequence $\{V_{2d}\}$; a half-dimensional sieve for primes of the form $a^2+2b^2$ combined with information about the fixed integer $V_{2d}$; or a genuinely new congruence law for the linked pair $(\ord_r(2),\ord_r(\omega_3))$. Each would be new arithmetic input rather than a refinement of the existing reduction.

\medskip
\noindent\emph{The diagonal after the half-pinning package.} On the diagonal --- the sub-sum $\sum_d S_{\mathrm{diag}}(d)$ of $\Pi$ --- the situation is sharper, in both directions. On the one hand, residue algebra on the diagonal is \emph{complete}. At a prime $r \in \mathcal{L}(d, p)$ with $\ord_r(2) = 2p$, every natural residue datum the configuration supplies lies in the single cyclic subgroup $\langle 2 \rangle \leq \F_r^\times$ of order $2 p$: $-1 = 2^p$; $V_d \equiv \varepsilon\, 2^{(p+1)/2}$ (Lemma~\ref{lem:half-pinning}); $2 U_d \equiv \varepsilon'$ (Proposition~\ref{prop:corner-localization}); and, at a prime $\pi$ of $\Z[\sqrt{-2}]$ above $r$ chosen with $\pi \mid V_d + \sqrt{-2}$ (Remark~\ref{rem:zsigmondy}), $\sqrt{-2} \equiv -V_d \pmod \pi$, while of the conjugate elements $\beta_{\pm} = 1 \pm (V_d/2)\sqrt{-2}$ of norm $V_{2d}/2$, the one not divisible by $\pi$ reduces to the trace $\beta_+ + \beta_- = 2$ itself. The subgroup generated by all of these is therefore exactly $\langle 2 \rangle$, so every $\mu_p$-character evaluated on them is tautological --- the inert-diagonal incarnation of face (A) of Observation~\ref{obs:additive-multiplicative} --- and the only non-tautological residue data are the two sign bits $(\varepsilon, \varepsilon')$, both of which Lemma~\ref{lem:half-pinning} and Proposition~\ref{prop:corner-localization} already extract. There is no further residue-level information on the diagonal to be had; what remains is size (the caps of Remark~\ref{rem:pair-caps}) and realizability. On the other hand, this localizes the entire open content of the diagonal in one named question.

\medskip
\noindent\textbf{Anti-aligned corner question.} \emph{Can the integers $G(\varepsilon, \varepsilon')$ and $G(-\varepsilon, -\varepsilon')$ of Proposition~\ref{prop:corner-localization} simultaneously contain primitive inert prime divisors exceeding $10^{12}$ at a single prime $p$ satisfying condition (v)?}
\medskip

The two corner integers are coprime --- each is odd, and an odd common prime divisor would divide both $2 V_d$ and $2^{(p+1)/2 + 1}$ --- so the question concerns the joint large-prime content of two coprime integers attached to one $(d, p)$. A diagonal pair at any admissible $d > 124$ (the range left open by Proposition~\ref{prop:diagonal-closures}) either occupies two anti-aligned corners or is confined by Remark~\ref{rem:pair-caps} to the divisors of a single explicit integer of roughly half the size of $V_{2d}$. We record one honest calibration: every statement of the package --- the identity, the corners, the caps --- uses only the quadratic recurrence, the order of $2$, and integer size, and transfers verbatim to the analogous Lucas data over $\F_q[t]$, a setting in which the corresponding pair statement is false (we use this nowhere; it is recorded as a guard against over-reading). The package can therefore localize the diagonal difficulty but cannot alone close it: closing the anti-aligned corner question requires input of the same fixed-divisor type as faces (C)/(A) above --- a determinacy law for which corners of which $(d, p)$ capture large primitive divisors.

\section{Computational verification}\label{sec:computational}

This appendix keeps only the certification facts used in the paper. Full tables, scripts, and longer diagnostic logs are treated as supplementary material.

\subsection{Known Wagstaff values}

The current record and probable-prime status for Wagstaff numbers is maintained by PrimePages \cite{ref:PrimePagesWagstaff}; it changes over time and is not an input to any implication in the reduction chain (the certified ladder of Corollary~\ref{cor:ladder} cites it only for the underlying primality certificates of the $W_{p-2}$). The finite computations here use explicit exponent lists in the supplementary material. Three certified cases $p\in\{2617,10501,12391\}$ also have ECPP-independent certificates in the companion paper \cite{ref:A}.

\subsection{Split-branch verification}\label{sec:R3-verify}

All 185 known prime factors $r\equiv1\pmod8$ of composite $W_p$ satisfy Conjecture R3. For each listed factor, $r$ is independently certified prime and the order condition $p\mid\ord_{\F_r^\times}(1+s)$ (with $s^2=2$) is checked directly; the case labels record which structural shortcut explains the computation:
\[
\begin{array}{lrr}
\hline
\text{Configuration} & \text{Count} & \text{Reference}\\
\hline
\text{Case A }(r\mid U_p) & 1 & \text{Theorem~\ref{thm:caseA}}\\
\text{Case B }(r\mid V_p) & 6 & \text{Theorem~\ref{thm:caseB}}\\
\text{Case C, }4U_p^2\equiv1 & 46 & \text{Theorem~\ref{thm:caseC}}\\
\text{Case C, inadmissible }d & 43 & \text{Proposition~\ref{prop:rank-equiv}}\\
\text{Case C, general }(d>200) & 89 & \text{direct R3 verification}\\
\hline
\end{array}
\]
Together these cover every known split factor.

\subsection{NCT and primitive-divisor scans}\label{sec:primdiv-survey}\label{sec:nct-verify}

For NCT, Theorem~\ref{thm:nct99} gives the unconditional range $d\leq99$. The parity-unblocked values in $100\leq d\leq200$ are
\[
  \{107,121,129,131,139,163,171,177,179\}.
\]
All nine are closed by complete fixed-$d$ certificates in the sense of Corollary~\ref{cor:nct-fixed-d-certificate}, and the method is uniform across this subsection: the primitive part $\Psi_{4d}$ is completely factored; the factorization is re-verified by an exact product identity against a locally recomputed $\Psi_{4d}$, with every asserted prime factor independently certified (APR-CL, \S\ref{sec:toolchain}); and every split factor $r \equiv 1 \pmod 8$ is excluded at step~2 of the certificate --- odd $\ord_r(2)$, $v_2(\ord_r(2)) \geq 2$, or composite half-order $\ord_r(2)/2$ --- except for the single step-3 instance at $d = 131$ noted below. Per-value facts and certificate scripts are recorded here; the factorizations and half-order decompositions omitted below are reproduced by the named scripts. The worked pair $d = 107$, $121$ illustrates the certificate. For $d=107$:
\[
\begin{aligned}
  \Psi_{428}={}&
  161522909729\cdot1262880768427\cdot17449564219651\\
  &{}\cdot2906272939091099\cdot132168496320829356461361353683.
\end{aligned}
\]
The only split factor is $r=161522909729$, and
\[
  \ord_r(2)/2=10095181858
  =2\cdot107\cdot47173747
\]
is composite. For $d=121$:
\[
\begin{aligned}
  \Psi_{484}={}&
  1451\cdot568699\cdot1197595081\cdot1614230507
  \cdot135855589627\\
  &{}\cdot7499151961531\cdot249507118668323
  \cdot4005463110187137017.
\end{aligned}
\]
The split factors are $1197595081$ and $4005463110187137017$. For the first:
\[
  \ord_{1197595081}(2)/2
  =9979959=3\cdot11^2\cdot19\cdot1447,
\]
and for $r=4005463110187137017$:
\[
\begin{aligned}
  \ord_r(2)/2
  &=500682888773392127\\
  &=11^2\cdot23\cdot3739247\cdot48113327.
\end{aligned}
\]
Both half-orders are composite, so no compatible triple exists at $d=107$ or $d=121$. The script \texttt{scripts/cp350\_nct\_107\_121\_verify.py} verifies the displayed primitive-part factorizations, primality of each listed factor by SymPy and PARI/GP \texttt{isprime}, and the split half-order decompositions.

The four further values $d=129,139,171,177$ are closed by the same certificate over the complete factorizations of their primitive parts ($65$, $106$, $83$ and $89$ decimal digits respectively):
\[
\begin{aligned}
  \Psi_{516}={}&
  75390435818587404514891047039001\\
  &{}\cdot276431598001487177038457370080041,\\[2pt]
  \Psi_{556}={}&
  6568027\cdot574593323881\\
  &{}\cdot197671787019026728856728443573405844788899\\
  &{}\cdot5764333243492029733503944458218529327657872457,\\[2pt]
  \Psi_{684}={}&
  8209\cdot10259\cdot1186057\cdot82658931465203\\
  &{}\cdot300531150542299499\\
  &{}\cdot19269878832774783314803453706733003299,\\[2pt]
  \Psi_{708}={}&
  6362089\cdot21015437931105774861319763619250473457\\
  &{}\cdot490240983802792962966192466615950084458511337.
\end{aligned}
\]
{\sloppy
The per-value facts: at $d=129$ both factors are split, the $32$-digit factor of odd $\ord_r(2)$ and the $33$-digit factor of composite half-order; at $d=139$ the two factors $\equiv3\pmod8$ are inert and irrelevant to the split search, and of the two split factors, $r=574593323881$ has odd $\ord_r(2)$ and the $46$-digit factor has composite half-order; at $d=171$ four factors are inert, the split factor $r=8209$ has half-order $1026=2\cdot3^3\cdot19$ (even, hence composite), and the split factor $r=1186057$ has odd $\ord_r(2)$; at $d=177$ all three factors are split, with composite half-orders. No split factor of any of the four primitive parts has prime half-order, so step 2 of the certificate already excludes every candidate (step 3 is never reached), and no compatible triple exists at $d\in\{129,139,171,177\}$. The script \texttt{scripts/}\allowbreak\texttt{cp351\_nct\_certificate\_129\_139\_171\_177.py} verifies the displayed primitive-part factorizations (exact product identities against $\Psi_{4d}$) and the split half-order decompositions; primality of all fifteen factors is certified by APR-CL (PARI/GP \texttt{isprime} with flag~$2$) in the independent end-to-end recomputation \texttt{scripts/}\allowbreak\texttt{cp352\_verify\_nct\_recompute.py}, which also re-derives the primitive parts from scratch, recomputes every rank of apparition and every order with minimality witnesses, and re-tests the exact condition $d\mid W_{p_r-2}$.\par}

For $d=163$ the $125$-digit primitive part $\Psi_{652}$ is completely factored into seven primes, all certified by APR-CL. Of the three split factors, two have composite half-orders and the $23$-digit factor has odd $\ord_r(2)$; so no compatible triple exists at $d=163$. (The $20$-digit split factor is the prime of Remark~\ref{rem:strengthening-false}.) The certificate is \texttt{scripts/}\allowbreak\texttt{cp353\_nct\_certificate\_163.py}.

For $d=131$ the $100$-digit $\Psi_{524}$ is completely factored (GNFS) into nine primes. Of the five split factors, three have odd $\ord_r(2)$ or half-order divisible by $4$, and the $35$-digit factor has composite half-order. The factor $r=1049$ is the only certificate instance in this paper reaching step 3: its half-order is the \emph{prime} $p_r=131=d$, and the exact condition $d \mid W_{p_r-2}$ fails because $p \nmid W_{p-2}$ (Lemma~\ref{lem:pndivW}) --- the $d = p$ diagonal. No compatible triple exists at $d=131$. Certificate: \texttt{scripts/}\allowbreak\texttt{cp353\_nct\_certificate\_131.py}.

For $d=179$ the $137$-digit $\Psi_{716}$ is completely factored (ECM, $B_1 \leq 11{\times}10^6$, final $58$-digit prime cofactor) into seven primes. All five split factors are excluded at step 2: four have odd $\ord_r(2)$, and the $35$-digit factor has composite half-order. No compatible triple exists at $d=179$. Certificate: \texttt{scripts/}\allowbreak\texttt{cp353\_nct\_certificate\_179.py}. All sixteen prime factors of $\Psi_{524}$ and $\Psi_{716}$ are certified by APR-CL (PARI \texttt{isprime} flag~$2$).

A separate finite-window sieve searched all nine parity-unblocked values above: all primes $r=8dk+1$ with $r<10^{11}$ were enumerated, rank $\rho(r)=4d$ was tested by the Pell recurrence, and surviving primitive divisors were checked for a compatible prime $p$ with $\ord_r(2)=2p$ and $d\mid W_{p-2}$. The run tested about $1.07\times10^8$ primes, found 12 primitive Pell divisors, and found no compatible triple in this window. Since the structural bound of Proposition~\ref{prop:nct-band} is far larger than $10^{11}$ for some of these $d$, this is evidence only, not a complete verification of $100\leq d\leq200$. After the nine complete closures at $d=107,121,129,131,139,163,171,177,179$, no parity-unblocked value of $100 \leq d \leq 200$ remains open: NCT holds unconditionally for every odd $d \leq 200$.

The window $200 < d \leq 300$ contains ten parity-unblocked values: $d = 201$, $209$, $211$, $227$, $243$, $249$, $251$, $281$, $283$, $297$. At six of them the primitive part $\Psi_{4d}$ is completely factored in the public factorization tables (sizes $102$--$215$ decimal digits; the products were re-verified exactly against locally recomputed $\Psi_{4d}$, and all $37$ prime factors were re-certified by APR-CL, PARI \texttt{isprime} flag~$2$), and the fixed-$d$ certificate closes each one: at $d = 201$, $209$, $249$, $251$, $281$, $297$ every split factor is excluded at step~2 --- odd $\ord_r(2)$ or composite half-order; none reaches the condition-(v) test. The multiplicative orders were computed by explicit descent from completely factored $r-1$ (four stubborn cofactors of $37$--$70$ digits were factored by ECM for this purpose; every piece was certified prime and the products verified). \emph{No compatible triple exists at $d = 201$, $209$, $249$, $251$, $281$, or $297$.} Certificates: \texttt{scripts/}\allowbreak\texttt{cp356\_nct\_certificates\_200\_300.py}, \texttt{scripts/}\allowbreak\texttt{cp356\_aprcl\_recert.py}. The value $d = 243$ was closed separately: its $125$-digit $\Psi_{972}$ fell entirely to ECM ($B_1 \leq 11{\times}10^6$, about half an hour of a 12-thread workstation), $\Psi_{972} = 3\cdot971\cdot15114601\cdot2068037891\cdot935133871378249\cdot p_{35}\cdot p_{55}$, all seven factors APR-CL-certified; the three split factors all have odd $\ord_r(2)$, so no compatible triple exists at $d = 243$ (certificate \texttt{scripts/}\allowbreak\texttt{cp358\_nct\_certificate\_243.py}). The value $d = 211$ fell the same way: its $161$-digit $\Psi_{844}$ was completely factored by ECM in about twenty minutes ($146857 \cdot 338779913 \cdot 52213127964665809 \cdot 307892103084433531 \cdot p_{33} \cdot p_{81}$, all APR-CL-certified; every split factor excluded at step~2; certificate \texttt{scripts/}\allowbreak\texttt{cp358\_nct\_certificate\_211.py}). The window $(300, 400]$ contains nine parity-unblocked values. Six --- $d = 307$, $321$, $347$, $361$, $387$, $393$, with $\Psi_{4d}$ of $163$ to $290$ digits --- are completely factored in the public tables; the same re-verification (exact products against locally recomputed $\Psi_{4d}$, $46$ factors APR-CL, orders by descent from completely factored $r - 1$) and the same certificate close all six; and $d = 363$'s $169$-digit $\Psi_{1452}$ fell to a single ECM round ($11 \cdot 2352241 \cdot 69796187 \cdot p_{17} \cdot p_{25} \cdot p_{36} \cdot p_{78}$, seven factors APR-CL, every split factor with composite half-order; certificate \texttt{scripts/}\allowbreak\texttt{cp358\_nct\_certificate\_363.py}): \emph{no compatible triple exists at any of these fifteen values} (certificates \texttt{scripts/}\allowbreak\texttt{cp358\_nct\_certificates\_range.py}). Two further values fell to a Pell half-split: for odd $d$ the Pell--Lucas number factors as $V_{2d} = 8 P_d^2 - 2 = 2\,(2P_d - 1)(2P_d + 1)$ with coprime odd halves, so any cofactor of $\Psi_{4d} \mid V_{2d}$ splits by a gcd with the halves into two pieces of roughly half the digit size. At $d = 227$ the ECM-stalled $C143$ cofactor split as $p_{84} \times p_{60}$, completing $\Psi_{908} = 907 \cdot p_{84} \cdot 57203 \cdot 5580569 \cdot 16552486912925363 \cdot p_{60}$ (six factors APR-CL; the two split factors have odd $\ord_r(2)$, excluded at step~2): \emph{no compatible triple exists at $d = 227$} (certificate \texttt{scripts/}\allowbreak\texttt{cp367\_nct\_certificate\_227.py}). At $d = 379$ the halves brought the $C290$ primitive part down to ECM scale, completing $\Psi_{1516} = 4547 \cdot p_{141} \cdot 351030947611 \cdot p_{29} \cdot p_{30} \cdot p_{75}$ (six factors APR-CL, exact product identity; both split factors, $p_{30}$ and $p_{141}$, have odd $\ord_r(2)$): \emph{no compatible triple exists at $d = 379$} (certificate \texttt{scripts/}\allowbreak\texttt{cp372\_nct\_certificate\_379.py}). The last two parity-unblocked values of $(200, 400]$ fell to the same half-split followed by GNFS on the single remaining half-cofactor (\texttt{cado-nfs}, ECM to $t35$ first). At $d = 283$ the $L$-half residual $C104$ factored as $p_{45} \cdot p_{59}$, completing $\Psi_{1132} = 12451 \cdot p_{45} \cdot p_{59} \cdot 6793 \cdot 11489975459 \cdot 21850712277340062427 \cdot p_{75}$ (seven factors APR-CL, exact product identity; the three split factors are excluded --- $6793$ by composite half-order, $p_{45}$ by $v_2(\ord_r(2)) \geq 2$, and $p_{59}$ by odd $\ord_r(2)$): \emph{no compatible triple exists at $d = 283$} (certificate \texttt{scripts/}\allowbreak\texttt{cp375\_nct\_certificate\_283.py}). At $d = 331$ the $M$-half residual $C106$ factored as $p_{36} \cdot p_{71}$, completing $\Psi_{1324} = 6619 \cdot 43691 \cdot 4510867 \cdot 224936385721 \cdot p_{22} \cdot p_{78} \cdot p'_{22} \cdot p_{36} \cdot p_{71}$ (nine factors APR-CL, exact product identity; all five split factors excluded --- odd $\ord_r(2)$, $v_2 \geq 2$, or composite half-order): \emph{no compatible triple exists at $d = 331$} (certificate \texttt{scripts/}\allowbreak\texttt{cp376\_nct\_certificate\_331.py}). With these, every parity-unblocked value of $(200, 400]$ is closed: \emph{NCT holds unconditionally for every odd $d \leq 400$}.

A $d\leq1000$ primitive-divisor survey on the inert side finds five danger triples; three are phantom because the corresponding $W_p$ is prime, and the two real triples are closed at secondary factors as described below. The survey is empirical context only; the reduction (Theorem~\ref{thm:sharpened-chain}) uses the Platinum enumeration rather than this catalog.

\subsection{Inert enumeration and local passes}\label{sec:inert-survey}\label{sec:secondary}

Among 869 known inert factors $r\equiv3\pmod8$ with $r<2\times10^5$, Condition~(II) fails at every factor. In the extended Platinum enumeration, every inert prime factor $r\leq10^{12}$ of every composite $W_p$ with $p<5\times10^{11}$ was enumerated by the candidate form $r=2pk+1$. The exponent-indexed scan over $k$ reaches $r=10^{12}$ at each $p$; the small-exponent band $p<5\times10^5$, where this requires $k$ as large as ${\sim}\,10^{12}/2p$, was completed by a dedicated sub-scan ($41{,}498$ primes, no additional local pass), so the coverage $r\leq10^{12}$ is exhaustive across the full exponent range. The run recorded $684{,}965{,}381$ pairs $(p,r)$:
\[
\begin{array}{lrr}
\hline
\text{Outcome} & \text{Count} & \text{Fraction}\\
\hline
G_r\leq43 & 669{,}378{,}360 & 97.72\%\\
\text{order-excess exclusion} & 15{,}587{,}018 & 2.28\%\\
\text{local-pass witnesses} & 3 & <10^{-6}\%\\
\hline
\end{array}
\]
The three local-pass witnesses occur at distinct exponents:
\[
\begin{array}{rll}
1{,}251{,}488{,}009 & r=2p+1 & d=41{,}716{,}267,\\
10{,}916{,}765{,}939 & r=152{,}834{,}723{,}147 & d=38{,}208{,}680{,}787,\\
85{,}684{,}865{,}933 & r=2p+1 & d=57.
\end{array}
\]
Of the three Platinum local-pass witnesses, two are now discharged at secondary factors. The $d=57$ witness, at $p=85{,}684{,}865{,}933$, overlaps the $d\leq1000$ primitive-divisor catalog and is closed by the secondary factor below. The $d=41{,}716{,}267$ witness, at $p=1{,}251{,}488{,}009$, is closed by a secondary factor found by a directed scan of the progression $r=2pk+1$ with $k\leq10^{11}$. The third witness exponent $p=10{,}916{,}765{,}939$ remains undischarged --- the analogous directed scan at this exponent, complete for $k\leq10^{12}$, finds no second factor (the witness factor itself, at $k=7$, is the only hit in the progression $r=2pk+1$), its full $W_p$ cofactor is not resolved here, and the expected failure is recorded as the at-a-factor statement X1 (Conjecture~\ref{conj:x1}). Separately, the $d\leq1000$ primitive-divisor catalog contains a second real triple at $d=67$, also closed at a secondary factor:
\[
  \begin{aligned}
  d=57:&\quad r_2=7{,}675{,}646{,}348{,}774{,}307{,}083,\\
       &\quad \text{inert; parity obstruction}\\
       &\quad (d_{r_2}=3\cdot487\cdot1{,}313{,}423{,}399{,}858{,}711),\\
  d=67:&\quad r_2=5{,}474{,}333{,}343{,}676{,}555{,}561{,}818{,}313,\\
       &\quad \text{split; R3 verified directly},\\
  d=41{,}716{,}267:&\quad r_2=33{,}829{,}735{,}778{,}964{,}491\ (k=13{,}515{,}805),\\
       &\quad \text{inert; parity obstruction}\\
       &\quad (d_{r_2}=3\cdot41\cdot68{,}759{,}625{,}567{,}001).
  \end{aligned}
\]
These secondary closures are certified computational facts (Proposition~\ref{prop:witness-discharges} records the two witness-exponent discharges and their verification); they are not counted as unconditional entries in Table~\ref{tab:status}. The local-pass witness rows above remain part of the enumeration record; their discharge status is as stated. The census of Proposition~\ref{prop:danger-census-400} is the $d$-indexed complement of this $r$-indexed record: at composite $W_p$, the only locally passing inert factors with $d_r \leq 400$ --- for $r$ of \emph{any} size --- are the $d = 57$ Platinum witness and the $d = 67$ catalog row, both closed above.

\subsection{Software and supplementary material}\label{sec:toolchain}

Integer factorization uses SymPy/gmpy2 routines, Cunningham tables \cite{ref:Cunningham}, and the online factor database \cite{ref:FactorDB} only to locate candidate factors; every asserted prime factor used in finite exhaustion is independently certified by APR-CL in PARI/GP \cite{ref:CohenLenstra}, and every assembled factorization is re-verified by an exact product identity. BPSW and deterministic Miller--Rabin are preliminary screens only. The supplementary deposit records factor tables, Platinum enumeration outputs, direct Condition-(II) checks, and the longer closed-approaches log omitted from this lean draft. The v4.0 supplementary bundle includes the danger-census artifacts (\texttt{cp412\_danger\_census.json}) alongside the NCT certificates.

\emph{Convention (computational inputs).} Unconditional statements proved by finite computation --- the fixed-$d$ NCT closures, the Pell-sequence exclusion, the danger census, the witness discharges --- rest on explicitly listed factorizations and certificates, archived as above and re-verifiable directly from the artifacts; they are used as theorems, as is standard for certificate-backed results. Separately, the reduction theorems \emph{declare} as stated computational inputs every computational ingredient of the conditional chain: the Platinum enumeration (a $684{,}965{,}381$-row sweep, not re-verifiable without substantial compute) and the witness discharges (small, but part of the chain's ledger) --- so that each conditional statement carries its complete list of non-algebraic ingredients in its own wording, and Theorem~\ref{thm:spc-reduction} can be checked to carry none.

\begin{thebibliography}{99}

\bibitem{ref:AtkinMorain} A.~O.~L. Atkin and F.~Morain, \emph{Elliptic curves and primality proving}, Math. Comp. \textbf{61} (1993), 29--68.

\bibitem{ref:BennettLevin} M.~A.~Bennett and A.~Levin, \emph{The Nagell--Ljunggren equation via Runge's method}, Monatsh.\ Math.\ \textbf{177} (2015), 15--31.

\bibitem{ref:BHV} Yu.~Bilu, G.~Hanrot, and P.~M. Voutier, \emph{Existence of primitive divisors of Lucas and Lehmer numbers}, J. Reine Angew. Math. \textbf{539} (2001), 75--122.

\bibitem{ref:BLS} J.~Brillhart, D.~H. Lehmer, and J.~L. Selfridge, \emph{New primality criteria and factorizations of $2^m \pm 1$}, Math. Comp. \textbf{29} (1975), 620--647.

\bibitem{ref:Cunningham} J.~Brillhart, D.~H. Lehmer, J.~L. Selfridge, B.~Tuckerman, and S.~S. Wagstaff, Jr., \emph{Factorizations of $b^n \pm 1$}, 3rd ed., Contemp. Math. \textbf{22}, AMS, 2002; online updates at \url{https://homes.cerias.purdue.edu/~ssw/cun/}.

\bibitem{ref:BugeaudCaoMignotte2001} Y.~Bugeaud, Z.~Cao and M.~Mignotte, \emph{On simple $K_4$-groups}, J.~Algebra \textbf{241} (2001), 658--668.

\bibitem{ref:BugeaudMignotte2002} Y.~Bugeaud and M.~Mignotte, \emph{On the Diophantine equation $\frac{x^n-1}{x-1} = y^q$ with negative $x$}, in: Number Theory for the Millennium I (Urbana, 2000), A~K~Peters, Natick, MA, 2002, pp.~145--151.

\bibitem{ref:Chua} K.~S. Chua, \emph{Chebyshev polynomials and higher order Lucas--Lehmer algorithm}, arXiv:2010.02677, 2020.

\bibitem{ref:CohenLenstra} H.~Cohen and A.~K. Lenstra, \emph{Implementation of a new primality test}, Math. Comp. \textbf{48} (1987), 103--121.

\bibitem{ref:Grantham} J.~Grantham, \emph{Frobenius pseudoprimes}, Math. Comp. \textbf{70} (2001), 873--891.

\bibitem{ref:HeathBrown} D.~R. Heath-Brown, \emph{Artin's conjecture for primitive roots}, Quart. J. Math. \textbf{37} (1986), 27--38.

\bibitem{ref:Hooley} C.~Hooley, \emph{On Artin's conjecture}, J. Reine Angew. Math. \textbf{225} (1967), 209--220.

\bibitem{ref:Iwaniec} H.~Iwaniec, \emph{Primes of the type $\varphi(x,y)+A$}, Acta Arith. \textbf{21} (1972), 203--234.

\bibitem{ref:Mihailescu} P.~Mih{\u a}ilescu, \emph{Primary cyclotomic units and a proof of Catalan's conjecture}, J. Reine Angew. Math. \textbf{572} (2004), 167--195.

\bibitem{ref:Moree} P.~Moree, \emph{Artin's primitive root conjecture --- a survey}, Integers \textbf{12} (2012), 1305--1416.

\bibitem{ref:OEISWieferich} OEIS Foundation Inc., \emph{A001220: Wieferich primes}, \url{https://oeis.org/A001220/internal}, accessed 2026-06-05.

\bibitem{ref:PrimePagesWagstaff} The PrimePages, \emph{Wagstaff primes}, \url{https://t5k.org/top20/page.php?id=67}, accessed 2026-07-16.

\bibitem{ref:Stewart} C.~L. Stewart, \emph{On divisors of Lucas and Lehmer numbers}, Acta Math. \textbf{211} (2013), 291--314.

\bibitem{ref:FactorDB} M.~Tervooren, \emph{Factor Database}, \url{http://factordb.com/}, accessed 2026-06-11.

\bibitem{ref:VrbaReix} A.~Vrba and T.~Reix, \emph{The Wagstaff primality test conjecture} (with a standing 500-euro reward for a published proof), mersenneforum.org, thread~18657, \url{https://mersenneforum.org/showthread.php?t=18657}, accessed 2026-07-15.

\bibitem{ref:Williams} H.~C. Williams, \emph{\'Edouard Lucas and Primality Testing}, CMS Monographs 22, Wiley, 1998.

\bibitem{ref:A} A.~Dolotov, \emph{Three Brillhart--Lehmer--Selfridge primality proofs for Wagstaff numbers}, Zenodo, 2026, doi:10.5281/zenodo.19645478.

\bibitem{ref:VRNote} A.~Dolotov, \emph{On the exact statement of the Vrba--Reix primality test for Wagstaff numbers}, companion note, 2026.

\end{thebibliography}

\end{document}
